% Barycentres — Mathématiques 1re Tech — Cours signé OnBuch+
% Programme officiel MINESEC. Compiler avec : tectonic N09-barycentres.tex
\newcommand{\DOCMATIERE}{Mathématiques}
\newcommand{\DOCNIVEAU}{1re Tech}
\newcommand{\DOCMODULE}{Mathématiques – classe de Première technique}
\newcommand{\DOCLECON}{N09}
\newcommand{\DOCTITRE}{Barycentres}
% OnBuch+ — préambule « cours signés » ENRICHI (compilé avec tectonic / XeLaTeX)
% Système visuel « Pop » d'OnBuch+ : fond crème (jamais blanc), bordures encre
% épaisses, ombres « dures » décalées, titres Archivo Black en capitales.
% Version MATHÉMATIQUES : graphes (pgfplots/tikz), boîtes pédagogiques
% (définition, propriété, méthode, attention, le savais-tu, exercice résolu,
% exercice, TP, bilan), page de garde et signature OnBuch+ ancrée en bas.
%
% Macros attendues dans main.tex AVANT \input{preamble} :
%   \DOCMATIERE  \DOCNIVEAU  \DOCMODULE  \DOCLECON (numéro)  \DOCTITRE
\documentclass[11pt]{article}

\usepackage{fontspec}
\usepackage[french]{babel}
\usepackage[a4paper,margin=2cm,top=3.4cm,bottom=2.8cm,headheight=1.4cm,footskip=1.3cm]{geometry}
\usepackage{amsmath,amssymb,amsfonts}
\usepackage{mathtools} % \xmapsto, \coloneqq... souvent utilisés par les modèles
\usepackage{cancel} % simplifications barrées
% \abs{z} (module, valeur absolue) et \norm{u} (norme), utilisés par les modèles
\providecommand{\abs}{}\let\abs\relax\DeclarePairedDelimiter{\abs}{\lvert}{\rvert}
\providecommand{\norm}{}\let\norm\relax\DeclarePairedDelimiter{\norm}{\lVert}{\rVert}
\usepackage[table]{xcolor} % [table] : fournit aussi \rowcolors (alternance de lignes)
\usepackage{pagecolor}
\usepackage{tikz}
\usetikzlibrary{arrows.meta,positioning,calc,shapes.geometric,shapes.misc,shapes.arrows,decorations.pathmorphing,decorations.pathreplacing,decorations.markings,patterns,shadows,mindmap,trees,fit,backgrounds,matrix,intersections,angles,quotes}
\usepackage{pgfplots}
\usepackage{pgf-pie} % diagrammes circulaires \pie (statistiques)
\pgfplotsset{compat=1.18}
\usepgfplotslibrary{fillbetween,groupplots} % aires entre courbes (calcul intégral)
\usepackage{enumitem}
\usepackage{fancyhdr}
% [most] charge toutes les bibliothèques tcolorbox (listings, minted, raster,
% documentation...) et gonfle inutilement la mémoire TeX sur un document long
% (~300 boîtes) : on ne charge que ce qui est réellement utilisé.
\usepackage[skins,breakable]{tcolorbox}
\usepackage{graphicx}
\usepackage{colortbl}
\usepackage{booktabs}
\usepackage{tabularx}
\usepackage{diagbox} % case d'en-tête diagonale des tableaux à double entrée
% Colonnes extensibles centrée / gauche / droite (C, L, R), souvent utilisées par les modèles
\newcolumntype{C}{>{\centering\arraybackslash}X}
\newcolumntype{L}{>{\raggedright\arraybackslash}X}
\newcolumntype{R}{>{\raggedleft\arraybackslash}X}
\usepackage{array}
\usepackage{multicol}
\usepackage{siunitx}
% parse-numbers=false : accepte aussi \SI{2,5 \times 10^{-3}}{...}
\sisetup{locale=FR,output-decimal-marker={,},group-minimum-digits=4,parse-numbers=false}
\DeclareSIUnit\ohm{\ensuremath{\symup{\Omega}}} % U+2126 absent de la police
\usepackage{qrcode}
% \degree : commande fréquemment utilisée par les modèles pour un angle ou une
% température en dehors de \SI{}{} (non fournie par les paquets chargés).
\providecommand{\degree}{^{\circ}}
% \qed (fin de démonstration), utilisé par les modèles sans amsthm
\providecommand{\qed}{\hfill$\square$}
% \barre : double barre des valeurs interdites dans un tableau de variations (tkz-tab)
\providecommand{\barre}{\ensuremath{\|}}
% environnement proof (amsthm non chargé) utilisé par les modèles
\ifdefined\proof\else\newenvironment{proof}[1][Démonstration]{\par\noindent\textit{#1.}\ }{\qed\par}\fi
\usepackage{lastpage}
\usepackage{hyperref}
\hypersetup{hidelinks}

% ============================================================
% Palette « Pop » OnBuch+ — mêmes hex que la classe P (plus_ui.dart)
% ============================================================
\definecolor{poporange}{HTML}{F59321}
\definecolor{popdark}{HTML}{E07A0C}
\definecolor{popcream}{HTML}{FFF6E8}
\definecolor{popink}{HTML}{1C1714}
\definecolor{poppurple}{HTML}{7A5AE0}
\definecolor{popgreen}{HTML}{2E9E6A}
\definecolor{poppink}{HTML}{E0457B}
\definecolor{popblue}{HTML}{2F7DE1}
\definecolor{popgold}{HTML}{C98A12}
\definecolor{popmuted}{HTML}{8A7F76}
\definecolor{popline}{HTML}{EADFD2}
\definecolor{popgray}{HTML}{8A7F76}
\colorlet{popgrey}{popgray}
% Alias tolérants (le modèle nomme parfois les couleurs différemment)
\colorlet{popwhite}{white}
\colorlet{popblack}{popink}
\colorlet{popred}{poppink}
\colorlet{popyellow}{popgold}
% Teintes claires pour fonds de tableaux / zones
\colorlet{popblueL}{popblue!10!white}
\colorlet{poporangeL}{poporange!14!white}
\colorlet{popgreenL}{popgreen!12!white}
\colorlet{poppurpleL}{poppurple!12!white}
\colorlet{poppinkL}{poppink!12!white}
\colorlet{popgoldL}{popgold!14!white}

\pagecolor{popcream}

% ============================================================
% Typographie
% ============================================================
\defaultfontfeatures{Ligatures=TeX}
\setmainfont{PlusJakartaSans-Medium.ttf}[Path=../../../fonts/,BoldFont=PlusJakartaSans-Bold.ttf]
% Maths en sans-serif (Fira Math) avec chiffres et lettres droites de Plus
% Jakarta, pour que les formules se fondent dans le texte.
\usepackage[math-style=ISO,bold-style=ISO]{unicode-math}
\setmathfont{FiraMath-Regular.otf}[Path=../../../fonts/]
\setmathfont{PlusJakartaSans-Medium.ttf}[Path=../../../fonts/,range={up/num,up/{Latin,latin},\mathup}]
\usepackage{icomma} % virgule décimale sans espace en mode math
% Caractères mathématiques Unicode absents des polices de texte (Plus Jakarta,
% Archivo Black) : rendus par leur équivalent mathématique, en texte comme en
% formule — sinon ils s'affichent en carré vide (ex. « Arithmétique dans ℤ »).
\usepackage{newunicodechar}
\newunicodechar{ℝ}{\ensuremath{\mathbb{R}}}
\newunicodechar{ℤ}{\ensuremath{\mathbb{Z}}}
\newunicodechar{ℂ}{\ensuremath{\mathbb{C}}}
\newunicodechar{ℕ}{\ensuremath{\mathbb{N}}}
\newunicodechar{ℚ}{\ensuremath{\mathbb{Q}}}
\newunicodechar{𝔻}{\ensuremath{\mathbb{D}}}
\newunicodechar{α}{\ensuremath{\alpha}}
\newunicodechar{β}{\ensuremath{\beta}}
\newunicodechar{γ}{\ensuremath{\gamma}}
\newunicodechar{δ}{\ensuremath{\delta}}
\newunicodechar{ε}{\ensuremath{\varepsilon}}
\newunicodechar{θ}{\ensuremath{\theta}}
\newunicodechar{λ}{\ensuremath{\lambda}}
\newunicodechar{μ}{\ensuremath{\mu}}
\newunicodechar{σ}{\ensuremath{\sigma}}
\newunicodechar{τ}{\ensuremath{\tau}}
\newunicodechar{φ}{\ensuremath{\varphi}}
\newunicodechar{ψ}{\ensuremath{\psi}}
\newunicodechar{ω}{\ensuremath{\omega}}
\newunicodechar{Σ}{\ensuremath{\Sigma}}
% majuscules produites par \MakeUppercase dans les titres (α-aminés -> Α)
\newunicodechar{Α}{\ensuremath{\alpha}}
\newunicodechar{Β}{\ensuremath{\beta}}
\newunicodechar{Γ}{\ensuremath{\Gamma}}
\newunicodechar{Δ}{\ensuremath{\Delta}}
\newunicodechar{Ε}{\ensuremath{\varepsilon}}
\newunicodechar{Θ}{\ensuremath{\Theta}}
\newunicodechar{Λ}{\ensuremath{\Lambda}}
\newunicodechar{Μ}{\ensuremath{\mu}}
\newunicodechar{Τ}{\ensuremath{\tau}}
\newunicodechar{Φ}{\ensuremath{\Phi}}
\newunicodechar{Ψ}{\ensuremath{\Psi}}
\newunicodechar{Ω}{\ensuremath{\Omega}}
\newunicodechar{↦}{\ensuremath{\mapsto}}
\newunicodechar{∈}{\ensuremath{\in}}
\newunicodechar{∉}{\ensuremath{\notin}}
\newunicodechar{∧}{\ensuremath{\wedge}}
\newunicodechar{⇔}{\ensuremath{\Leftrightarrow}}
\newunicodechar{⟺}{\ensuremath{\Longleftrightarrow}}
\newunicodechar{⇒}{\ensuremath{\Rightarrow}}
\newunicodechar{⟹}{\ensuremath{\Longrightarrow}}
\newunicodechar{∘}{\ensuremath{\circ}}
\newunicodechar{∪}{\ensuremath{\cup}}
\newunicodechar{∩}{\ensuremath{\cap}}
\newunicodechar{≡}{\ensuremath{\equiv}}
\newunicodechar{⊂}{\ensuremath{\subset}}
\newunicodechar{⊥}{\ensuremath{\perp}}
\newunicodechar{∃}{\ensuremath{\exists}}
\newunicodechar{∀}{\ensuremath{\forall}}
\newunicodechar{∛}{\ensuremath{\sqrt[3]{\,}}}
\newunicodechar{∜}{\ensuremath{\sqrt[4]{\,}}}
\newunicodechar{⃗}{\ensuremath{{}^{\to}}}
\newunicodechar{ⁿ}{\ifmmode^{n}\else\textsuperscript{\ensuremath{n}}\fi}
\newunicodechar{ˣ}{\ifmmode^{x}\else\textsuperscript{\ensuremath{x}}\fi}
\newunicodechar{ᵇ}{\ifmmode^{b}\else\textsuperscript{\ensuremath{b}}\fi}
\newunicodechar{ʳ}{\ifmmode^{r}\else\textsuperscript{\ensuremath{r}}\fi}
\newunicodechar{ᵘ}{\ifmmode^{u}\else\textsuperscript{\ensuremath{u}}\fi}
\newunicodechar{ʸ}{\ifmmode^{y}\else\textsuperscript{\ensuremath{y}}\fi}
\newunicodechar{ᵃ}{\ifmmode^{a}\else\textsuperscript{\ensuremath{a}}\fi}
\newunicodechar{ᵉ}{\ifmmode^{e}\else\textsuperscript{\ensuremath{e}}\fi}
\newunicodechar{ᵏ}{\ifmmode^{k}\else\textsuperscript{\ensuremath{k}}\fi}
\newunicodechar{ᵖ}{\ifmmode^{p}\else\textsuperscript{\ensuremath{p}}\fi}
\newunicodechar{ᴺ}{\ifmmode^{N}\else\textsuperscript{\ensuremath{N}}\fi}
\newunicodechar{ᶜ}{\ifmmode^{c}\else\textsuperscript{\ensuremath{c}}\fi}
\newunicodechar{ʰ}{\ifmmode^{h}\else\textsuperscript{\ensuremath{h}}\fi}
\newunicodechar{ᵠ}{\ifmmode^{q}\else\textsuperscript{\ensuremath{q}}\fi}
\newunicodechar{ₙ}{\ifmmode_{n}\else\textsubscript{\ensuremath{n}}\fi}
\newunicodechar{ₐ}{\ifmmode_{a}\else\textsubscript{\ensuremath{a}}\fi}
\newunicodechar{ᵢ}{\ifmmode_{i}\else\textsubscript{\ensuremath{i}}\fi}
\newunicodechar{ᵣ}{\ifmmode_{r}\else\textsubscript{\ensuremath{r}}\fi}
\newunicodechar{ₖ}{\ifmmode_{k}\else\textsubscript{\ensuremath{k}}\fi}
\newunicodechar{ᵦ}{\ifmmode_{\beta}\else\textsubscript{\ensuremath{\beta}}\fi}
\newunicodechar{ₓ}{\ifmmode_{x}\else\textsubscript{\ensuremath{x}}\fi}
\newfontfamily\popdisplay{ArchivoBlack-Regular.ttf}[Path=../../../fonts/]
\newfontfamily\poplogo{SpaceGrotesk-Bold.ttf}[Path=../../../fonts/]
\newfontfamily\popsemi{PlusJakartaSans-SemiBold.ttf}[Path=../../../fonts/]
\newfontfamily\popxbold{PlusJakartaSans-ExtraBold.ttf}[Path=../../../fonts/]
\newfontfamily\popxboldls{PlusJakartaSans-ExtraBold.ttf}[Path=../../../fonts/,LetterSpace=3.0]

\setlist[enumerate]{leftmargin=1.7em,itemsep=5pt,topsep=4pt}
\setlist[itemize]{leftmargin=1.7em,itemsep=4pt,topsep=4pt,label={\color{poporange}\textbullet}}
\setlength{\parindent}{0pt}
\setlength{\parskip}{5pt}
\renewcommand{\arraystretch}{1.25}

% pgfplots : style Pop par défaut
\pgfplotsset{
  every axis/.append style={axis line style={popink,line width=1pt},tick style={popink},
    grid style={popline,line width=0.5pt},label style={font=\small},tick label style={font=\footnotesize}},
  popcurve/.style={poporange,line width=1.8pt,no markers,smooth},
}
% Même style disponible dans un \draw TikZ simple (hors pgfplots)
\tikzset{popcurve/.style={poporange,line width=1.8pt}}

% ============================================================
% Pill / squiggle
% ============================================================
\newcommand{\poppill}[3]{%
  \tikz[baseline=-0.55ex]\node[draw=popink,line width=1.2pt,rounded corners=2.6pt,%
    fill=#2,inner xsep=6.5pt,inner ysep=3pt]{%
    \popxboldls\fontsize{7}{7}\selectfont\color{#3}\MakeUppercase{#1}};%
}
\newcommand{\popsquiggle}[1]{%
  \tikz[baseline=-0.4ex]{
    \draw[#1,line width=2.6pt,line cap=round]
      plot [smooth] coordinates {(0,0.09) (0.34,0.01) (0.68,0.21) (1.02,0.01) (1.36,0.10)};
  }%
}
% Mot-clé surligné (marqueur orange sous le texte)
\newcommand{\cle}[1]{{\popxbold\color{popdark}#1}}

% ============================================================
% En-tête / pied
% ============================================================
\pagestyle{fancy}
\fancyhf{}
\renewcommand{\headrulewidth}{0pt}
\renewcommand{\footrulewidth}{0pt}
\lhead{\raisebox{-0.15ex}{\poplogo\fontsize{13}{13}\selectfont\color{popink}OnBuch\color{poporange}+}}
\rhead{\poppill{\DOCMATIERE\ \textbullet\ \DOCNIVEAU\ \textbullet\ Leçon \DOCLECON}{white}{popink}}
\lfoot{\popxbold\fontsize{7.6}{7.6}\selectfont\color{popmuted}\MakeUppercase{Cours signé OnBuch+}\ \textbullet\ {\popsemi\DOCTITRE}}
\rfoot{\tikz[baseline=-0.6ex]\node[rounded corners=8.5pt,fill=popink,minimum height=17pt,minimum width=17pt,inner xsep=5pt,inner sep=0]{\popxbold\fontsize{8}{8}\selectfont\color{popcream}\thepage\,/\,\pageref*{LastPage}};}

% ============================================================
% Titres
% ============================================================
\newcommand{\chaptitre}[2]{%
  \begin{tcolorbox}[enhanced,colback=poporange,colframe=popink,boxrule=2.6pt,arc=11pt,
      drop shadow=popink,left=15pt,right=15pt,top=12pt,bottom=11pt,before skip=3pt,after skip=18pt]
    \poppill{#2}{popink}{popcream}\par\vspace{9pt}
    {\raggedright\hyphenpenalty=10000\exhyphenpenalty=10000\popdisplay\fontsize{22}{24}\selectfont\color{popcream}\MakeUppercase{#1}\par}
  \end{tcolorbox}
}
% Section numérotée : pastille orange avec le numéro + titre Archivo
\newcounter{popsec}
\newcommand{\coursec}[1]{%
  \stepcounter{popsec}\setcounter{popsub}{0}%
  \par\vspace{16pt}\noindent
  \begin{minipage}{\linewidth}
  \tikz[baseline=-0.7ex]\node[draw=popink,line width=1.6pt,fill=poporange,rounded corners=4pt,minimum size=22pt,inner sep=0]{\popdisplay\fontsize{12}{12}\selectfont\color{popcream}\thepopsec};%
  \hspace{8pt}{\popdisplay\fontsize{15}{17}\selectfont\color{popink}\MakeUppercase{#1}}\par
  \vspace{4pt}\noindent{\color{poporange}\rule{36pt}{3.6pt}}
  \end{minipage}\par\nopagebreak\vspace{8pt}
}
\newcounter{popsub}
\newcommand{\courssub}[1]{%
  \stepcounter{popsub}%
  \par\vspace{9pt}\noindent{\popxbold\fontsize{11.5}{13}\selectfont\color{popink}{\color{poporange}\thepopsec.\thepopsub}\hspace{6pt}#1}\par\nopagebreak\vspace{4pt}
}

% ============================================================
% Boîtes pédagogiques — même recette : fond blanc, bordure épaisse colorée,
% ombre dure encre, badge plein. Titre optionnel [#1].
% ============================================================
\tcbset{popbase/.style={breakable,enhanced,colback=white,boxrule=2.3pt,arc=9pt,
  shadow={3pt}{-3pt}{0pt}{popink},left=14pt,right=14pt,top=11pt,bottom=11pt,before skip=11pt,after skip=15pt,
  fontupper=\popsemi\fontsize{10.6}{14.5}\selectfont\color{popink},
  fontlower=\popsemi\fontsize{10.6}{14.5}\selectfont\color{popink}}}
% Coche dessinée (le glyphe ✓ n'existe pas dans les polices du document)
\newcommand{\popcheck}[1]{\tikz[baseline=-0.5ex]\draw[#1,line width=1.6pt,line cap=round,line join=round](-0.13,0.0)--(-0.04,-0.09)--(0.13,0.1);}
\newcommand{\popboxhead}[3]{\poppill{#1}{#2}{white}\ifx\relax#3\relax\else\hspace{6pt}{\popxbold\fontsize{10.6}{12}\selectfont\color{popink}#3}\fi\par\vspace{7pt}}

\newtcolorbox{aretenir}[1][]{popbase,colframe=popgreen,before upper={\popboxhead{À retenir}{popgreen}{#1}}}
\newtcolorbox{exemplebox}[1][]{popbase,colframe=poporange,before upper={\popboxhead{Exemple}{poporange}{#1}}}
\newtcolorbox{definition}[1][]{popbase,colframe=popblue,before upper={\popboxhead{Définition}{popblue}{#1}}}
\newtcolorbox{propriete}[1][]{popbase,colframe=poppurple,before upper={\popboxhead{Propriété}{poppurple}{#1}}}
\newtcolorbox{methode}[1][]{popbase,colframe=popgold,before upper={\popboxhead{Méthode}{popgold}{#1}}}
\newtcolorbox{attention}[1][]{popbase,colframe=poppink,before upper={\popboxhead{Attention}{poppink}{#1}}}
\newtcolorbox{experience}[1][]{popbase,colframe=popblue,colback=popblueL,before upper={\popboxhead{Expérience}{popblue}{#1}}}
% « Le savais-tu ? » : carte violette pleine (contexte, culture, Cameroun)
\newtcolorbox{savaistu}[1][]{popbase,colback=poppurple,colframe=popink,
  fontupper=\popsemi\fontsize{10.4}{14.2}\selectfont\color{white},
  before upper={\poppill{Le savais-tu\,?}{white}{poppurple}\ifx\relax#1\relax\else\hspace{6pt}{\popxbold\color{white}#1}\fi\par\vspace{7pt}}}
% Exercice résolu : énoncé en haut, \tcblower puis la solution sur fond crème
\newcounter{popexo}\newcounter{popexr}
\newtcolorbox{exoresolu}[1][]{popbase,colframe=poporange,colbacklower=popcream,
  segmentation style={popink,line width=1.2pt,dashed},
  before upper={\stepcounter{popexr}\popboxhead{Exercice résolu \thepopexr}{poporange}{#1}},
  before lower={\poppill{Solution}{popink}{popcream}\par\vspace{6pt}}}
% Exercice (non résolu en place) — la correction est dans la partie Corrigés
\newtcolorbox{exercice}[1][]{popbase,colframe=popink,
  before upper={\stepcounter{popexo}\popboxhead{Exercice \thepopexo}{popink}{#1}}}
% Corrigé
\newtcolorbox{corrige}[1][]{popbase,colframe=popgreen,colback=popgreenL,
  before upper={\popboxhead{Corrigé}{popgreen}{#1}}}
% Objectifs / prérequis / bilan
\newtcolorbox{objectifs}{popbase,colframe=popink,colback=poporangeL,
  before upper={\popboxhead{Objectifs de la leçon}{popink}{}}}
\newtcolorbox{prerequis}{popbase,colframe=popink,colback=popblueL,
  before upper={\popboxhead{Prérequis}{popblue}{}}}
\newtcolorbox{bilan}{popbase,colframe=popink,colback=white,boxrule=2.8pt,
  before upper={\popboxhead{Fiche bilan}{poporange}{L'essentiel à connaître pour l'examen}}}
% Figure encadrée avec légende
\newtcolorbox{popfigure}[1][]{enhanced,breakable=false,colback=white,colframe=popline,boxrule=1.4pt,arc=8pt,
  left=10pt,right=10pt,top=10pt,bottom=8pt,before skip=10pt,after skip=12pt,halign=center,
  lower separated=false,halign lower=center,
  fontlower=\popsemi\fontsize{9}{11.5}\selectfont\color{popmuted},
  #1}
\newcommand{\legende}[1]{\par\vspace{6pt}{\popsemi\fontsize{9}{11.5}\selectfont\color{popmuted}#1}}

% ============================================================
% Signature OnBuch+
% ============================================================
\newcommand{\poplogoinline}[1]{{\poplogo\fontsize{#1}{#1}\selectfont\color{popink}OnBuch\color{poporange}+}}
\newcommand{\signatureonbuch}{%
  \begin{tcolorbox}[enhanced,colback=popink,colframe=popink,boxrule=2.2pt,arc=10pt,
      drop shadow=poporange,left=14pt,right=14pt,top=10pt,bottom=10pt,before skip=0pt,after skip=0pt]
    \tikz[baseline=-0.7ex]\node[circle,fill=popgreen,draw=popcream,line width=1.3pt,minimum size=22pt,inner sep=0]{\popcheck{white}};%
    \hspace{9pt}\begin{minipage}[c]{0.62\linewidth}
      {\popxbold\fontsize{12}{13}\selectfont\color{popcream}Signé\ \textbullet\ {\poplogo OnBuch\color{poporange}+}}\par\vspace{2pt}
      {\raggedright\popsemi\fontsize{8}{10.5}\selectfont\color{popline}\MakeUppercase{Équipe pédagogique OnBuch+}\\\MakeUppercase{Conforme au programme officiel MINESEC}\par}
    \end{minipage}\hfill
    {\poplogo\fontsize{22}{22}\selectfont\color{popcream}OnBuch\color{poporange}+}
  \end{tcolorbox}%
}

% ============================================================
% Page de garde
% #1 titre · #2 matière · #3 niveau · #4 module · #5 n° leçon · #6 durée ·
% #7 description / objectifs (texte ou itemize)
% ============================================================
\newcommand{\pagegarde}[7]{%
  \thispagestyle{empty}
  \newgeometry{a4paper,margin=1.7cm,top=1.6cm,bottom=1.4cm}%
  \begin{tikzpicture}[remember picture,overlay]
    \fill[poporange!18] ([xshift=-3.2cm,yshift=-3.6cm]current page.north east) circle (4.6cm);
    \fill[poppurple!14] ([xshift=2.4cm,yshift=7.6cm]current page.south west) circle (3.8cm);
  \end{tikzpicture}%
  {\poplogo\fontsize{30}{30}\selectfont\color{popink}OnBuch\color{poporange}+}\hfill
  \raisebox{0.6ex}{\poppill{Cours signé \textbullet\ Premium}{popink}{popcream}}\par
  \vspace{1pt}\popsquiggle{poppurple}\par
  \vspace{8pt}
  \begin{tcolorbox}[enhanced,colback=poporange,colframe=popink,boxrule=2.8pt,arc=13pt,
      drop shadow=popink,left=18pt,right=18pt,top=12pt,bottom=13pt,after skip=10pt]
    \poppill{#2 \textbullet\ #3}{popink}{popcream}\hspace{5pt}\poppill{#4}{white}{popink}\par\vspace{8pt}
    {\popdisplay\fontsize{14}{14}\selectfont\color{popink}LEÇON #5}\par\vspace{4pt}
    {\raggedright\hyphenpenalty=10000\exhyphenpenalty=10000\popdisplay\fontsize{25}{27}\selectfont\color{popcream}\MakeUppercase{#1}\par}
    \vspace{8pt}
    {\popxbold\fontsize{9.5}{9.5}\selectfont\color{popink}\MakeUppercase{Durée conseillée : #6}}
  \end{tcolorbox}
  \begin{tcolorbox}[enhanced,colback=white,colframe=popink,boxrule=2.2pt,arc=10pt,
      drop shadow=popink,left=16pt,right=16pt,top=10pt,bottom=10pt,after skip=8pt]
    \poppill{Dans cette leçon}{poppurple}{white}\par\vspace{5pt}
    \popsemi\fontsize{9.3}{12.6}\selectfont\color{popink}#7
  \end{tcolorbox}
  \noindent\begin{minipage}[t]{0.487\linewidth}
    \begin{tcolorbox}[enhanced,colback=popgreen,colframe=popink,boxrule=2.2pt,arc=10pt,
        drop shadow=popink,left=11pt,right=11pt,top=10pt,bottom=10pt,height=3.1cm,valign=center]
      \tcbox[colback=white,colframe=popink,boxrule=1.3pt,arc=5pt,boxsep=0pt,nobeforeafter,
        left=3pt,right=3pt,top=3pt,bottom=3pt,valign=center]{\qrcode[height=1.9cm]{https://play.google.com/store/apps/details?id=com.luvvix.onbuch&hl=fr}}%
      \hspace{8pt}\begin{minipage}[c]{0.5\linewidth}
        {\popdisplay\fontsize{11}{12.5}\selectfont\color{white}\MakeUppercase{Télécharge OnBuch}}\par\vspace{4pt}
        {\raggedright\popsemi\fontsize{8.4}{11}\selectfont\color{white}Gratuit \textbullet\ Google Play\\Tuteur IA, annales, concours, résultats\par}
      \end{minipage}
    \end{tcolorbox}
  \end{minipage}\hfill
  \begin{minipage}[t]{0.487\linewidth}
    \begin{tcolorbox}[enhanced,colback=poppurple,colframe=popink,boxrule=2.2pt,arc=10pt,
        drop shadow=popink,left=11pt,right=11pt,top=10pt,bottom=10pt,height=3.1cm,valign=center]
      \tikz[baseline=-0.7ex]\node[circle,fill=white,draw=popink,line width=1.3pt,minimum size=1.6cm,inner sep=0]{\popdisplay\fontsize{24}{24}\selectfont\color{poppurple}+};%
      \hspace{8pt}\begin{minipage}[c]{0.6\linewidth}
        {\popdisplay\fontsize{11}{12.5}\selectfont\color{white}\MakeUppercase{Passe à OnBuch+}}\par\vspace{4pt}
        {\raggedright\popsemi\fontsize{8.4}{11}\selectfont\color{white}Tous les cours signés, Tuteur IA illimité, fiches PDF — onglet OnBuch+ de l'app\par}
      \end{minipage}
    \end{tcolorbox}
  \end{minipage}
  \vspace{8pt}
  \popcontenu
  \vfill
  \signatureonbuch
  \restoregeometry
  \newpage
}

% Rangée « Ce cours contient » (page de garde)
\newcommand{\poptile}[3]{% #1 couleur #2 gros texte #3 légende
  \begin{tcolorbox}[enhanced,colback=#1,colframe=popink,boxrule=2pt,arc=9pt,shadow={2.5pt}{-2.5pt}{0pt}{popink},
      left=6pt,right=6pt,top=9pt,bottom=9pt,halign=center,valign=center,height=1.75cm,width=\linewidth,nobeforeafter]
    {\popdisplay\fontsize{12.5}{13}\selectfont\color{white}\MakeUppercase{#2}}\par\vspace{3pt}
    {\centering\popsemi\fontsize{7.8}{9.5}\selectfont\color{white}#3\par}
  \end{tcolorbox}}
\newcommand{\popcontenu}{%
  \par\noindent{\popxboldls\fontsize{7.5}{7.5}\selectfont\color{popmuted}CE COURS SIGNÉ CONTIENT}\par\vspace{7pt}
  \noindent\begin{minipage}[t]{0.235\linewidth}\poptile{popblue}{Cours}{complet, illustré, pas à pas}\end{minipage}\hfill
  \begin{minipage}[t]{0.235\linewidth}\poptile{poporange}{Méthodes}{et exercices résolus}\end{minipage}\hfill
  \begin{minipage}[t]{0.235\linewidth}\poptile{poppink}{Exercices}{type examen + corrigés}\end{minipage}\hfill
  \begin{minipage}[t]{0.235\linewidth}\poptile{popgreen}{Fiche bilan}{l'essentiel à retenir}\end{minipage}\par}

% Sommaire
\newtcolorbox{sommaire}{enhanced,colback=white,colframe=popink,boxrule=2.2pt,arc=10pt,
  drop shadow=popink,left=18pt,right=18pt,top=13pt,bottom=13pt,before skip=6pt,after skip=20pt,
  before upper={\poppill{Sommaire}{poppurple}{white}\par\vspace{9pt}\popsemi\fontsize{10.6}{14.5}\selectfont\color{popink}}}

% Signature de fin de cours — poussée tout en bas de la dernière page
\newcommand{\findecours}{%
  \par\vfill\nopagebreak
  \begin{center}\popsquiggle{poporange}\end{center}\vspace{4pt}
  \signatureonbuch
}
\AtBeginDocument{\def\llbracket{\mathopen{[\mkern-3.5mu[}}\def\rrbracket{\mathclose{]\mkern-3.5mu]}}} % intervalles d'entiers (glyphe ⟦ absent de la police)
% Glyphes absents de Fira Math : redéfinis à partir de glyphes disponibles
\AtBeginDocument{%
  \def\setminus{\mathbin{\backslash}}\let\smallsetminus\setminus
  \def\vdots{\mathord{\vbox{\baselineskip3.2pt\lineskiplimit0pt\kern4pt\hbox{.}\hbox{.}\hbox{.}}}}%
  \def\ddots{\mathinner{\mkern1mu\raise6.5pt\hbox{.}\mkern2mu\raise3.6pt\hbox{.}\mkern2mu\raise0.7pt\hbox{.}\mkern1mu}}%
  \def\triangle{\mathord{\scalebox{0.9}{$\symup{\Delta}$}}}%
  \def\nabla{\mathord{\raisebox{\depth}{\scalebox{1}[-1]{$\symup{\Delta}$}}}}%
  \def\checkmark{\popcheck{popdark}}%
  \def\star{\mathbin{\ast}}\def\bigstar{\mathord{\ast}}%
  \def\diamond{\mathbin{\scalebox{0.6}{$\Diamond$}}}%
  \def\bigcup{\mathop{\mathchoice{\raisebox{-0.3em}{\scalebox{1.7}{$\cup$}}}{\scalebox{1.25}{$\cup$}}{\cup}{\cup}}\displaylimits}%
  \def\bigcap{\mathop{\mathchoice{\raisebox{-0.3em}{\scalebox{1.7}{$\cap$}}}{\scalebox{1.25}{$\cap$}}{\cap}{\cap}}\displaylimits}%
  \def\lesseqqgtr{\mathrel{\lessgtr}}\def\bigoplus{\mathop{\oplus}\displaylimits}%
}
\providecommand{\ssi}{si et seulement si} % abréviation fréquente
\AtBeginDocument{\providecommand{\wideparen}{\overparen}} % arc de cercle

%% FIGFIX-BEGIN v5 — correctifs globaux des figures (docs/onbuch-generation/tools/figfix)
\usepackage{adjustbox}\usepackage{etoolbox}\tcbuselibrary{raster}
% toute figure TikZ/pgfplots plus large que la ligne est réduite à la largeur disponible
\BeforeBeginEnvironment{tikzpicture}{\begin{adjustbox}{max width=\linewidth}}
\AfterEndEnvironment{tikzpicture}{\end{adjustbox}}
% légendes pgfplots lisibles par-dessus les courbes
\pgfplotsset{every axis/.append style={legend style={fill=white,fill opacity=0.92,text opacity=1,draw=black!15,inner sep=3pt}}}
% étiquettes d'arêtes (syntaxe edge["..."]) avec fond blanc
\tikzset{every edge quotes/.append style={fill=white,inner sep=1.2pt}}
%% FIGFIX-END
\begin{document}

\pagegarde{Barycentres}{Mathématiques}{1re Tech}{Mathématiques – classe de Première technique}{N09}{2 séances (fiche de progression harmonisée)}{Cette leçon introduit le barycentre de deux, trois et quatre points pondérés, ses propriétés fondamentales (associativité, homogénéité, réduction vectorielle) et son application à la détermination de lignes de niveau. Les élèves apprendront à construire, localiser et exploiter le barycentre dans des situations géométriques et des problèmes concrets de la vie courante.
\vspace{4pt}
\begin{multicols}{2}\small
\begin{itemize}
\item À la fin de cette leçon, je sais définir le barycentre de deux points pondérés et le caractériser vectoriellement
\item À la fin de cette leçon, je sais construire et placer le barycentre de deux points pondérés sur une figure
\item À la fin de cette leçon, je sais déterminer l'ensemble des points M tels que αMA→ + βMB→ soit colinéaire à un vecteur donné ou nul
\item À la fin de cette leçon, je sais définir et construire le barycentre de trois ou quatre points pondérés
\item À la fin de cette leçon, je sais utiliser l'associativité du barycentre (barycentres partiels) pour simplifier une construction ou une démonstration
\item À la fin de cette leçon, je sais démontrer qu'un point est barycentre de points donnés et justifier des alignements ou des concours de droites
\end{itemize}\end{multicols}}

\begin{sommaire}
\begin{multicols}{2}
\begin{enumerate}
\item Barycentre de deux points pondérés
\item Propriétés fondamentales du barycentre
\item Barycentre de trois points pondérés
\item Barycentre de quatre points et applications géométriques
\item Lignes de niveau
\item Applications à la vie courante et synthèse
\item Activité d'intégration
\item Méthodes et exercices résolus
\item Exercices
\item Corrigés des exercices
\item Fiche bilan
\end{enumerate}
\end{multicols}
\end{sommaire}

\begin{objectifs}
\begin{itemize}
\item À la fin de cette leçon, je sais définir le barycentre de deux points pondérés et le caractériser vectoriellement
\item À la fin de cette leçon, je sais construire et placer le barycentre de deux points pondérés sur une figure
\item À la fin de cette leçon, je sais déterminer l'ensemble des points M tels que αMA→ + βMB→ soit colinéaire à un vecteur donné ou nul
\item À la fin de cette leçon, je sais définir et construire le barycentre de trois ou quatre points pondérés
\item À la fin de cette leçon, je sais utiliser l'associativité du barycentre (barycentres partiels) pour simplifier une construction ou une démonstration
\item À la fin de cette leçon, je sais démontrer qu'un point est barycentre de points donnés et justifier des alignements ou des concours de droites
\item À la fin de cette leçon, je sais déterminer des lignes de niveau du type M ↦ αMA² + βMB² = k ou M ↦ MA/MB = k
\item À la fin de cette leçon, je sais modéliser une situation concrète d'équilibre (levier, centre de gravité) par un barycentre et rédiger une démarche justifiée
\end{itemize}
\end{objectifs}

\begin{prerequis}
\begin{itemize}
\item Vecteurs du plan : somme, produit par un réel, colinéarité (Seconde)
\item Relation de Chasles et réduction de sommes vectorielles
\item Coordonnées d'un point et d'un vecteur dans un repère
\item Milieux, médianes et centre de gravité d'un triangle (géométrie du triangle)
\item Équations de droites et de cercles dans un repère orthonormé
\end{itemize}
\end{prerequis}

\newpage

\coursec{Barycentre de deux points pondérés}

\courssub{Situation d'accroche et problématique}

Mama Ngo Bell, soudeuse à Douala, vient de recevoir une barre métallique de \SI{2}{m} qu'elle doit partager entre ses deux ateliers. Mais voilà : à une extrémité de la barre est soudé un moteur, deux fois plus lourd que l'outil fixé à l'autre extrémité. Pour transporter la barre sans qu'elle ne bascule, elle cherche le \cle{point d'équilibre} où poser le support. Intuitivement, elle sent bien que ce point n'est pas le milieu de la barre : il doit être plus proche du moteur, du côté le plus lourd. Mais où exactement ?

À Yaoundé, un menuisier se pose une question voisine : où se trouve le centre de gravité d'une table dont les pieds supportent des charges différentes ?

Dans les deux cas, il s'agit de trouver un point qui « résume » la position de plusieurs points auxquels on a attaché des \emph{poids} différents. Ce point porte un nom : le \cle{barycentre}. C'est l'outil mathématique du point d'équilibre, et nous allons voir qu'il permet aussi de résoudre élégamment de nombreux problèmes de géométrie.

\textbf{Problématique :} comment définir, localiser et calculer précisément le point d'équilibre de deux points affectés de masses ?

\courssub{Notion de point pondéré}

\begin{definition}[Point pondéré]
Un \cle{point pondéré} est un couple $(A, \alpha)$ où $A$ est un point du plan (ou de l'espace) et $\alpha$ un nombre réel appelé \cle{coefficient} ou \cle{masse} du point $A$.
\end{definition}

L'intuition est celle de la physique : on imagine qu'au point $A$ on a accroché une masse de $\alpha$ kilogrammes. Le coefficient peut être :
\begin{itemize}
\item \textbf{positif} : le point « attire » le barycentre vers lui ;
\item \textbf{négatif} : le point « repousse » le barycentre (comme un ballon gonflé à l'hélium qui soulèverait la tige) ;
\item \textbf{nul} : le point n'a aucune influence.
\end{itemize}

\begin{attention}[Un coefficient nul ne supprime pas le point]
Si $\alpha = 0$, le point $(A, 0)$ n'a simplement aucun poids dans le calcul : le barycentre ne dépend plus de la position de $A$. Attention à ne pas confondre « coefficient nul » et « point inexistant ».
\end{attention}

\courssub{Définition du barycentre de deux points pondérés}

Reprenons la barre de Mama Ngo Bell. Posons $A$ l'extrémité avec le moteur (masse $\alpha$) et $B$ l'autre extrémité (masse $\beta$). Le point d'équilibre $G$ est le point où les « efforts » des deux masses se compensent exactement. Traduite en langage vectoriel, cette compensation s'écrit très simplement.

\begin{definition}[Barycentre de deux points pondérés]
Soient $(A, \alpha)$ et $(B, \beta)$ deux points pondérés tels que $\alpha + \beta \neq 0$.
On appelle \cle{barycentre} de $(A, \alpha)$ et $(B, \beta)$ l'unique point $G$ vérifiant :
\[ \alpha\,\overrightarrow{GA} + \beta\,\overrightarrow{GB} = \vec{0}. \]
On note $G = \mathrm{bar}\{(A, \alpha)\,;\,(B, \beta)\}$.
\end{definition}

\begin{attention}[La condition $\alpha + \beta \neq 0$ est indispensable]
Si $\alpha + \beta = 0$ (par exemple $\alpha = 1$ et $\beta = -1$), la relation devient $\overrightarrow{GA} - \overrightarrow{GB} = \vec{0}$, c'est-à-dire $\overrightarrow{BA} = \vec{0}$, ce qui est impossible dès que $A \neq B$. Il n'existe alors \textbf{aucun} barycentre. Avant tout calcul, vérifie toujours que la somme des coefficients est non nulle.
\end{attention}

\courssub{Existence, unicité et position du barycentre}

\begin{propriete}[Existence et unicité du barycentre — démonstration exigible]
Soient $(A, \alpha)$ et $(B, \beta)$ deux points pondérés avec $\alpha + \beta \neq 0$. Alors il existe un \textbf{unique} point $G$ tel que $\alpha\,\overrightarrow{GA} + \beta\,\overrightarrow{GB} = \vec{0}$, et ce point vérifie :
\[ \overrightarrow{AG} = \frac{\beta}{\alpha + \beta}\,\overrightarrow{AB}. \]
En particulier, $G$ appartient toujours à la droite $(AB)$.
\end{propriete}

\begin{experience}[Démonstration guidée : réduction vectorielle]
Partons de la relation de définition et transformons-la pas à pas, en faisant apparaître le point $A$ grâce à la relation de Chasles.
\begin{enumerate}
\item La relation de Chasles donne $\overrightarrow{GB} = \overrightarrow{GA} + \overrightarrow{AB}$.
\item Remplaçons dans la définition :
\[ \alpha\,\overrightarrow{GA} + \beta\left(\overrightarrow{GA} + \overrightarrow{AB}\right) = \vec{0}. \]
\item Regroupons les termes en $\overrightarrow{GA}$ :
\[ (\alpha + \beta)\,\overrightarrow{GA} + \beta\,\overrightarrow{AB} = \vec{0}. \]
\item Comme $\alpha + \beta \neq 0$, on peut diviser :
\[ \overrightarrow{GA} = -\frac{\beta}{\alpha + \beta}\,\overrightarrow{AB}. \]
\item En multipliant par $-1$ (ce qui change le sens du vecteur) :
\[ \overrightarrow{AG} = \frac{\beta}{\alpha + \beta}\,\overrightarrow{AB}. \]
\end{enumerate}
Cette dernière égalité définit le point $G$ de manière \textbf{unique} : à partir de $A$, on parcourt la fraction $\dfrac{\beta}{\alpha+\beta}$ du vecteur $\overrightarrow{AB}$. Le barycentre existe donc, il est unique, et comme $\overrightarrow{AG}$ est colinéaire à $\overrightarrow{AB}$, le point $G$ est bien sur la droite $(AB)$. \hfill $\blacksquare$
\end{experience}

\begin{propriete}[Position de $G$ selon les signes des coefficients]
Soit $G = \mathrm{bar}\{(A, \alpha)\,;\,(B, \beta)\}$ avec $\alpha + \beta \neq 0$.
\begin{itemize}
\item $G$ appartient \textbf{toujours} à la droite $(AB)$ ;
\item si $\alpha$ et $\beta$ sont de \textbf{même signe}, alors $G$ appartient au \textbf{segment} $[AB]$ ;
\item si $\alpha$ et $\beta$ sont de \textbf{signes contraires}, alors $G$ est \textbf{à l'extérieur} du segment $[AB]$, du côté du point dont la valeur absolue du coefficient est la plus grande ;
\item plus la valeur absolue du coefficient d'un point est grande, plus $G$ est \textbf{proche} de ce point.
\end{itemize}
\end{propriete}

Cette propriété confirme l'intuition de Mama Ngo Bell : le point d'équilibre se rapproche de l'extrémité la plus lourde.

\begin{popfigure}
\begin{center}
\begin{tikzpicture}[scale=1.0]
% Cas 1 : coefficients positifs
\draw[popline, thick] (0,0) -- (5,0);
\fill[popblue] (0,0) circle (2.2pt) node[below=2pt] {$A$};
\fill[popblue] (5,0) circle (2.2pt) node[below=2pt] {$B$};
\fill[poporange] (3,0) circle (2.6pt) node[below=2pt] {$G_1$};
\node[above=3pt, popdark] at (2.5,0) {\small $\alpha=2,\ \beta=3$ : $G_1\in[AB]$};
% Cas 2 : isobarycentre
\draw[popline, thick] (0,-1.6) -- (5,-1.6);
\fill[popblue] (0,-1.6) circle (2.2pt) node[below=2pt] {$A$};
\fill[popblue] (5,-1.6) circle (2.2pt) node[below=2pt] {$B$};
\fill[popgreen] (2.5,-1.6) circle (2.6pt) node[below=2pt] {$G_2$};
\node[above=3pt, popdark] at (2.5,-1.6) {\small $\alpha=\beta=1$ : $G_2$ milieu de $[AB]$};
% Cas 3 : coefficient négatif
\draw[popline, thick] (0,-3.2) -- (5,-3.2);
\draw[popline, dashed, thick] (5,-3.2) -- (7.5,-3.2);
\fill[popblue] (0,-3.2) circle (2.2pt) node[below=2pt] {$A$};
\fill[popblue] (5,-3.2) circle (2.2pt) node[below=2pt] {$B$};
\fill[poppurple] (7.5,-3.2) circle (2.6pt) node[below=2pt] {$G_3$};
\node[above=3pt, popdark] at (3.75,-3.2) {\small $\alpha=-1,\ \beta=2$ : $G_3$ hors de $[AB]$};
\end{tikzpicture}
\end{center}
\legende{Trois positions du barycentre selon les signes des coefficients : à l'intérieur du segment (masses de même signe), au milieu (isobarycentre), à l'extérieur (masses de signes contraires).}
\end{popfigure}

\courssub{Cas particulier : l'isobarycentre}

\begin{definition}[Isobarycentre]
Lorsque les deux coefficients sont \textbf{égaux} ($\alpha = \beta \neq 0$), le barycentre est appelé \cle{isobarycentre} de $A$ et $B$. La relation devient :
\[ \overrightarrow{AG} = \frac{\alpha}{\alpha + \alpha}\,\overrightarrow{AB} = \frac{1}{2}\,\overrightarrow{AB}. \]
L'isobarycentre de $A$ et $B$ est donc le \textbf{milieu du segment} $[AB]$.
\end{definition}

Le milieu d'un segment, que tu connais depuis le collège, est donc un barycentre particulier : celui où les deux points « pèsent » le même poids. Le barycentre généralise le milieu.

\courssub{Méthode de construction du barycentre}

\begin{methode}[Placer le barycentre de $(A, \alpha)$ et $(B, \beta)$]
\begin{enumerate}
\item Vérifier que $\alpha + \beta \neq 0$.
\item Calculer le rapport $k = \dfrac{\beta}{\alpha + \beta}$.
\item Écrire $\overrightarrow{AG} = k\,\overrightarrow{AB}$ : on part de $A$ et on parcourt la fraction $k$ du segment $[AB]$ (en sens inverse si $k < 0$).
\item Placer $G$ sur la droite $(AB)$ et vérifier la cohérence : $G$ doit être plus proche du point au plus gros coefficient (en valeur absolue).
\end{enumerate}
\end{methode}

\begin{attention}[Erreur fréquente : inverser les coefficients]
Dans la formule $\overrightarrow{AG} = \dfrac{\beta}{\alpha + \beta}\,\overrightarrow{AB}$, c'est le coefficient de \textbf{l'autre point} ($\beta$, celui de $B$) qui apparaît au numérateur quand on part de $A$. C'est logique : plus $B$ est lourd, plus $G$ est attiré vers $B$, donc plus le rapport est grand. Une bonne façon de vérifier : si $\beta = 0$, alors $\overrightarrow{AG} = \vec{0}$ et $G = A$ : sans masse en $B$, tout le poids est en $A$.
\end{attention}

\begin{exoresolu}[Construction et coordonnées de $G = \mathrm{bar}\{(A, 2)\,;\,(B, 3)\}$]
\textbf{Énoncé.} Soient $A$ et $B$ deux points tels que $AB = \SI{5}{cm}$.
\begin{enumerate}
\item Construire le point $G = \mathrm{bar}\{(A, 2)\,;\,(B, 3)\}$ et calculer $AG$.
\item Dans un repère, on donne $A(1\,;\,2)$ et $B(6\,;\,7)$. Calculer les coordonnées de $G$.
\end{enumerate}
\tcblower
\textbf{Solution détaillée.}
\begin{enumerate}
\item La somme des coefficients est $2 + 3 = 5 \neq 0$ : le barycentre existe. On applique la formule :
\[ \overrightarrow{AG} = \frac{\beta}{\alpha + \beta}\,\overrightarrow{AB} = \frac{3}{5}\,\overrightarrow{AB}. \]
On place donc $G$ sur $[AB]$, aux $\dfrac{3}{5}$ du segment en partant de $A$ :
\[ AG = \frac{3}{5} \times AB = \frac{3}{5} \times 5 = \SI{3}{cm}. \]
Vérification : le coefficient de $B$ (3) est plus grand que celui de $A$ (2), donc $G$ est plus proche de $B$ : en effet $GB = 5 - 3 = \SI{2}{cm} < AG$. Cohérent.
\item On utilise les formules des coordonnées du barycentre :
\begin{align*}
x_G &= \frac{\alpha\, x_A + \beta\, x_B}{\alpha + \beta} = \frac{2 \times 1 + 3 \times 6}{5} = \frac{20}{5} = 4, \\
y_G &= \frac{\alpha\, y_A + \beta\, y_B}{\alpha + \beta} = \frac{2 \times 2 + 3 \times 7}{5} = \frac{25}{5} = 5.
\end{align*}
Donc $G(4\,;\,5)$. On remarque que $\overrightarrow{AG} = (3\,;\,3) = \dfrac{3}{5}\overrightarrow{AB}$ puisque $\overrightarrow{AB} = (5\,;\,5)$ : les deux méthodes concordent.
\end{enumerate}
\end{exoresolu}

\courssub{Coordonnées du barycentre dans un repère}

\begin{propriete}[Coordonnées du barycentre]
Dans un repère $(O\,;\,\vec{\imath}, \vec{\jmath})$, soient $A(x_A\,;\,y_A)$ et $B(x_B\,;\,y_B)$. Le barycentre $G$ de $(A, \alpha)$ et $(B, \beta)$, avec $\alpha + \beta \neq 0$, a pour coordonnées :
\[ x_G = \frac{\alpha\, x_A + \beta\, x_B}{\alpha + \beta} \qquad \text{et} \qquad y_G = \frac{\alpha\, y_A + \beta\, y_B}{\alpha + \beta}. \]
\end{propriete}

Ces formules se démontrent en projetant la relation vectorielle $\alpha\,\overrightarrow{GA} + \beta\,\overrightarrow{GB} = \vec{0}$ sur chaque axe. On retient qu'il s'agit d'une \emph{moyenne pondérée} des coordonnées : chaque point « vote » pour sa position avec un poids égal à son coefficient. C'est exactement le même principe que le calcul de ta moyenne trimestrielle avec des coefficients de matières !

\courssub{Interprétation physique : le point d'équilibre}

\begin{aretenir}[Barycentre et centre de gravité]
Physiquement, le barycentre de $(A, \alpha)$ et $(B, \beta)$ (avec $\alpha, \beta > 0$) est le \cle{point d'équilibre} d'une tige rigide sans masse portant une masse $\alpha$ en $A$ et une masse $\beta$ en $B$. Posée sur un appui placé en $G$, la tige reste horizontale. C'est la loi du levier : $\alpha \times AG = \beta \times BG$.
\end{aretenir}

\begin{popfigure}
\begin{center}
\begin{tikzpicture}[scale=1.0]
% Fléau
\draw[popdark, very thick] (0,0) -- (6,0);
% Masses
\fill[popblue!70] (0,-0.9) rectangle (1.2,0);
\node[white] at (0.6,-0.45) {\small $m_A$};
\fill[poporange!80] (5.2,-0.6) rectangle (6,0);
\node[white] at (5.6,-0.3) {\small $m_B$};
% Appui (triangle) en G
\fill[popgreen!70] (3.6,-1.1) -- (4.4,-1.1) -- (4,0) -- cycle;
\fill[poppurple] (4,0) circle (2.4pt) node[above=3pt] {$G$};
% Points A et B
\fill[popblue] (0.6,0) circle (2pt) node[above=3pt] {$A$};
\fill[popblue] (5.6,0) circle (2pt) node[above=3pt] {$B$};
% Cotes
\draw[<->, popmuted] (0.6,-1.5) -- (4,-1.5) node[midway, below] {\small $AG$};
\draw[<->, popmuted] (4,-1.5) -- (5.6,-1.5) node[midway, below] {\small $GB$};
\end{tikzpicture}
\end{center}
\legende{Fléau en équilibre : l'appui $G$ est le barycentre des points $A$ et $B$ pondérés par les masses $m_A$ et $m_B$. La masse la plus lourde est la plus proche de l'appui.}
\end{popfigure}

Revenons à Mama Ngo Bell : le moteur en $A$ est deux fois plus lourd que l'outil en $B$, donc $\alpha = 2\beta$. Le point d'équilibre vérifie :
\[ \overrightarrow{AG} = \frac{\beta}{2\beta + \beta}\,\overrightarrow{AB} = \frac{1}{3}\,\overrightarrow{AB}. \]
Pour une barre de \SI{2}{m}, le support doit être posé à $\dfrac{1}{3} \times 2 \approx \SI{0,67}{m}$ de l'extrémité $A$ (celle du moteur). Le problème concret est résolu par un simple calcul de barycentre.

\begin{savaistu}[D'où vient le mot « barycentre » ?]
Le mot \emph{barycentre} vient du grec \emph{barus} (« lourd ») et \emph{kentron} (« centre ») : littéralement, le « centre des pesanteurs ». Le grand savant grec \textbf{Archimède} (III\textsuperscript{e} siècle av. J.-C.) utilisait déjà cette idée pour énoncer ses célèbres lois du levier : « Donnez-moi un point d'appui, et je soulèverai le monde ! » Aujourd'hui, le barycentre est partout : les ingénieurs l'utilisent pour équilibrer les avions et les grues de chantier, les astronomes pour décrire comment la Terre et la Lune tournent autour de leur barycentre commun, et même les informaticiens pour lisser les images numériques.
\end{savaistu}

\begin{exercice}[À toi de jouer]
Sur une droite graduée, on place $A$ d'abscisse $-2$ et $B$ d'abscisse $4$.
\begin{enumerate}
\item Placer $G = \mathrm{bar}\{(A, 3)\,;\,(B, 1)\}$ et calculer son abscisse.
\item Placer $H = \mathrm{bar}\{(A, 5)\,;\,(B, -2)\}$ et calculer son abscisse. Que remarque-t-on sur sa position ?
\item Déterminer le réel $\alpha$ tel que $\mathrm{bar}\{(A, \alpha)\,;\,(B, 2)\}$ soit le point d'abscisse $0$.
\end{enumerate}
\end{exercice}

\begin{corrige}[Exercice 1]
\begin{enumerate}
\item $\overrightarrow{AG} = \dfrac{1}{3+1}\overrightarrow{AB} = \dfrac{1}{4}\overrightarrow{AB}$, donc $x_G = -2 + \dfrac{1}{4}\times 6 = -0{,}5$.
\item $\overrightarrow{AH} = \dfrac{-2}{5-2}\overrightarrow{AB} = -\dfrac{2}{3}\overrightarrow{AB}$, donc $x_H = -2 - \dfrac{2}{3}\times 6 = -6$. Les coefficients sont de signes contraires : $H$ est à l'extérieur de $[AB]$, du côté de $A$ (car $|\alpha| > |\beta|$).
\item On veut $x_G = \dfrac{\alpha \times (-2) + 2 \times 4}{\alpha + 2} = 0$, soit $-2\alpha + 8 = 0$, d'où $\alpha = 4$. Vérification : $\alpha + \beta = 6 \neq 0$. 
\end{enumerate}
\end{corrige}

\coursec{Propriétés fondamentales du barycentre}

Dans la section précédente, nous avons défini le barycentre de deux points pondérés : $G = \text{bar}\{(A,\alpha),(B,\beta)\}$ est l'unique point tel que $\alpha\overrightarrow{GA}+\beta\overrightarrow{GB}=\vec{0}$, lorsque $\alpha+\beta\neq 0$. Nous avons aussi vu que $\overrightarrow{AG}=\dfrac{\beta}{\alpha+\beta}\overrightarrow{AB}$. Ces définitions sont utiles, mais ce qui fait toute la puissance du barycentre, ce sont ses \cle{propriétés}. C'est l'objet de cette section : nous allons établir quatre propriétés fondamentales (homogénéité, réduction vectorielle, invariance par isométrie, caractérisation) et apprendre à les utiliser pour transformer des problèmes compliqués en problèmes simples.

\courssub{Homogénéité du barycentre}

\textbf{Intuition.} Imagine une balance (un levier d'atelier) avec une masse de \SI{1}{kg} en $A$ et une masse de \SI{2}{kg} en $B$. Le point d'équilibre $G$ est à une position bien précise. Si tu remplaces les masses par \SI{3}{kg} et \SI{6}{kg}, ou par \SI{0,5}{kg} et \SI{1}{kg}, le point d'équilibre \emph{ne bouge pas} : ce qui compte, c'est le \emph{rapport} des masses, pas leur valeur absolue. C'est exactement le contenu de la propriété d'homogénéité.

\begin{propriete}[Homogénéité du barycentre]
Soient $A$ et $B$ deux points, $\alpha$ et $\beta$ deux réels tels que $\alpha+\beta\neq 0$, et $k$ un réel \textbf{non nul}. Alors :
\[
\text{bar}\{(A,\alpha),(B,\beta)\}=\text{bar}\{(A,k\alpha),(B,k\beta)\}.
\]
Autrement dit, on peut multiplier (ou diviser) tous les coefficients par un même nombre non nul sans changer le barycentre.
\end{propriete}

\begin{experience}[Démonstration guidée de l'homogénéité]
Notons $G=\text{bar}\{(A,\alpha),(B,\beta)\}$. Montrons que $G$ est aussi le barycentre de $(A,k\alpha)$ et $(B,k\beta)$.
\begin{enumerate}
\item Par définition de $G$, on a : $\alpha\overrightarrow{GA}+\beta\overrightarrow{GB}=\vec{0}$.
\item Multiplions les deux membres par le réel $k$ : $k\left(\alpha\overrightarrow{GA}+\beta\overrightarrow{GB}\right)=k\,\vec{0}=\vec{0}$.
\item En distribuant : $k\alpha\overrightarrow{GA}+k\beta\overrightarrow{GB}=\vec{0}$.
\item Comme $k\neq 0$ et $\alpha+\beta\neq 0$, on a $k\alpha+k\beta=k(\alpha+\beta)\neq 0$. La relation de l'étape 3 signifie donc exactement que $G=\text{bar}\{(A,k\alpha),(B,k\beta)\}$. \hfill$\square$
\end{enumerate}
\end{experience}

\begin{exemplebox}[Simplification de coefficients]
\begin{itemize}
\item $\text{bar}\{(A,3),(B,6)\}=\text{bar}\{(A,1),(B,2)\}$ (on divise par $3$).
\item $\text{bar}\left\{\left(A,\dfrac{1}{2}\right),\left(B,\dfrac{1}{3}\right)\right\}=\text{bar}\{(A,3),(B,2)\}$ (on multiplie par $6$).
\item $\text{bar}\{(A,-2),(B,4)\}=\text{bar}\{(A,1),(B,-2)\}$ (on divise par $-2$).
\end{itemize}
L'homogénéité permet de se ramener à des coefficients entiers et petits, ce qui simplifie tous les calculs.
\end{exemplebox}

\begin{popfigure}
\begin{center}
\begin{tikzpicture}[scale=1.1]
% Segment AB
\draw[popline, thick] (0,0) -- (6,0);
\fill[popblue] (0,0) circle (2pt) node[below left] {$A$};
\fill[poporange] (6,0) circle (2pt) node[below right] {$B$};
% G = bar{(A,1),(B,2)} : AG = 2/3 AB
\fill[popink] (4,0) circle (2.5pt) node[below] {$G$};
% Coefficients version 1
\node[popblue, above] at (0,0.35) {coefficient $1$};
\node[poporange, above] at (6,0.35) {coefficient $2$};
% Coefficients version 2
\node[popblue!70, below=14pt] at (0,0) {ou $3$};
\node[poporange!80, below=14pt] at (6,0) {ou $6$};
% Braces distances
\draw[<->, popmuted] (0,-0.9) -- (4,-0.9) node[midway, below] {$AG=\dfrac{2}{3}AB$};
\draw[<->, popmuted] (4,-0.9) -- (6,-0.9) node[midway, below] {$GB=\dfrac{1}{3}AB$};
\end{tikzpicture}
\end{center}
\legende{Le barycentre de $(A,1)$ et $(B,2)$ est le même que celui de $(A,3)$ et $(B,6)$ : seul le rapport des coefficients compte. Le point $G$ est deux fois plus proche de $B$ que de $A$ car $B$ porte le coefficient le plus grand.}
\end{popfigure}

\begin{attention}[Multiplier par zéro est interdit]
La condition $k\neq 0$ est essentielle : si l'on multipliait les coefficients par $0$, on obtiendrait $(A,0)$ et $(B,0)$, dont la somme des coefficients est nulle : le barycentre n'existerait plus. De même, on ne peut multiplier \emph{qu'un seul} des deux coefficients : $\text{bar}\{(A,1),(B,2)\}\neq\text{bar}\{(A,2),(B,2)\}$ en général.
\end{attention}

\courssub{La propriété de réduction vectorielle}

\textbf{Intuition.} Voici l'outil le plus puissant de la leçon. Souvent, un problème donne une condition portant sur un point $M$ \emph{quelconque} du plan, par exemple une expression comme $2\overrightarrow{MA}+3\overrightarrow{MB}$. Cette expression semble dépendre de la position de $M$... mais la propriété de réduction affirme qu'elle se \emph{réduit} toujours à un seul vecteur passant par le barycentre $G$. Tout se passe comme si les deux points $A$ et $B$ pouvaient être « remplacés » par le point $G$ affecté de la somme des coefficients.

\begin{propriete}[Théorème fondamental de réduction]
Soient $A$ et $B$ deux points, $\alpha$ et $\beta$ deux réels tels que $\alpha+\beta\neq 0$, et $G=\text{bar}\{(A,\alpha),(B,\beta)\}$. Alors, \textbf{pour tout point $M$ du plan} :
\[
\alpha\overrightarrow{MA}+\beta\overrightarrow{MB}=(\alpha+\beta)\overrightarrow{MG}.
\]
Cette démonstration est exigible au Probatoire technique.
\end{propriete}

\begin{experience}[Démonstration pas à pas de la réduction]
Soit $M$ un point quelconque du plan. L'idée est de faire apparaître le point $G$ grâce à la relation de Chasles.
\begin{enumerate}
\item Pour chaque vecteur, on passe par $G$ : $\overrightarrow{MA}=\overrightarrow{MG}+\overrightarrow{GA}$ et $\overrightarrow{MB}=\overrightarrow{MG}+\overrightarrow{GB}$.
\item On remplace dans la somme :
\begin{align*}
\alpha\overrightarrow{MA}+\beta\overrightarrow{MB}
&=\alpha\left(\overrightarrow{MG}+\overrightarrow{GA}\right)+\beta\left(\overrightarrow{MG}+\overrightarrow{GB}\right)\\
&=\alpha\overrightarrow{MG}+\beta\overrightarrow{MG}+\alpha\overrightarrow{GA}+\beta\overrightarrow{GB}\\
&=(\alpha+\beta)\overrightarrow{MG}+\underbrace{\left(\alpha\overrightarrow{GA}+\beta\overrightarrow{GB}\right)}_{=\,\vec{0}\ \text{car}\ G\ \text{est le barycentre}}\\
&=(\alpha+\beta)\overrightarrow{MG}. \qquad \square
\end{align*}
\item La clé de la démonstration est donc la relation de Chasles, suivie de la définition même du barycentre.
\end{enumerate}
\end{experience}

\begin{popfigure}
\begin{center}
\begin{tikzpicture}[scale=1.15]
\coordinate (M) at (0,0);
\coordinate (A) at (1.2,2.2);
\coordinate (B) at (5,1.6);
% G = bar{(A,2),(B,3)} => AG = 3/5 AB
\coordinate (G) at ($(A)!0.6!(B)$);
\fill[popink] (M) circle (2pt) node[below left] {$M$};
\fill[popblue] (A) circle (2pt) node[above] {$A$};
\fill[poporange] (B) circle (2pt) node[right] {$B$};
\fill[popgreen] (G) circle (2.5pt) node[above right] {$G$};
\draw[->, popblue, very thick] (M) -- (A) node[midway, left] {$\overrightarrow{MA}$};
\draw[->, poporange, very thick] (M) -- (B) node[midway, below right] {$\overrightarrow{MB}$};
\draw[->, popgreen, very thick] (M) -- (G) node[midway, above, sloped] {$\overrightarrow{MG}$};
\draw[dashed, popmuted] (A) -- (B);
\end{tikzpicture}
\end{center}
\legende{Pour tout point $M$, la combinaison $\alpha\overrightarrow{MA}+\beta\overrightarrow{MB}$ se réduit au seul vecteur $(\alpha+\beta)\overrightarrow{MG}$ : les points $A$ et $B$ « s'effacent » au profit du barycentre $G$.}
\end{popfigure}

\begin{exoresolu}[Réduire $2\overrightarrow{MA}+3\overrightarrow{MB}$]
Soient $A$ et $B$ deux points et $G=\text{bar}\{(A,2),(B,3)\}$. Simplifier l'expression $2\overrightarrow{MA}+3\overrightarrow{MB}$ pour tout point $M$.
\tcblower
La somme des coefficients est $2+3=5\neq 0$, donc le barycentre $G$ existe et la réduction s'applique directement :
\[
2\overrightarrow{MA}+3\overrightarrow{MB}=(2+3)\overrightarrow{MG}=5\overrightarrow{MG}.
\]
Ainsi, pour \emph{n'importe quel} point $M$ du plan, la somme $2\overrightarrow{MA}+3\overrightarrow{MB}$ est égale à $5\overrightarrow{MG}$ : on a remplacé deux vecteurs par un seul.
\end{exoresolu}

\begin{attention}[Condition d'application : $\alpha+\beta\neq 0$]
La réduction n'est utilisable que si $\alpha+\beta\neq 0$, car c'est la condition d'existence du barycentre. Avant d'écrire $\alpha\overrightarrow{MA}+\beta\overrightarrow{MB}=(\alpha+\beta)\overrightarrow{MG}$, vérifie toujours que la somme des coefficients n'est pas nulle. Si $\alpha+\beta=0$, la situation est totalement différente : voir le paragraphe suivant.
\end{attention}

\courssub{Cas particulier : somme des coefficients nulle}

Que se passe-t-il si $\alpha+\beta=0$, c'est-à-dire si $\beta=-\alpha$ ? Reprenons le calcul de la démonstration : pour tout point $M$,
\[
\alpha\overrightarrow{MA}-\alpha\overrightarrow{MB}=\alpha\left(\overrightarrow{MA}-\overrightarrow{MB}\right)=\alpha\overrightarrow{BA}.
\]

\begin{propriete}[Somme constante]
Si $\alpha+\beta=0$ (avec $\alpha\neq 0$), alors pour tout point $M$ du plan :
\[
\alpha\overrightarrow{MA}+\beta\overrightarrow{MB}=\alpha\overrightarrow{BA},
\]
qui est un \cle{vecteur constant}, indépendant du point $M$.
\end{propriete}

\begin{exemplebox}[Exemple chiffré]
Pour tout point $M$ : $2\overrightarrow{MA}-2\overrightarrow{MB}=2\overrightarrow{BA}$. Que $M$ soit en $A$, en $B$ ou n'importe où ailleurs, le résultat est toujours le même vecteur $2\overrightarrow{BA}$. C'est d'ailleurs pour cela qu'il n'existe \emph{aucun} point $G$ vérifiant $2\overrightarrow{GA}-2\overrightarrow{GB}=\vec{0}$ : cela imposerait $2\overrightarrow{BA}=\vec{0}$, c'est-à-dire $A=B$.
\end{exemplebox}

\courssub{Caractérisation du barycentre}

La propriété de réduction permet de démontrer immédiatement l'équivalence fondamentale suivante, qui sert de définition pratique du barycentre.

\begin{propriete}[Caractérisation du barycentre]
Soient $\alpha$ et $\beta$ deux réels tels que $\alpha+\beta\neq 0$. Pour tout point $G$ du plan :
\[
G=\text{bar}\{(A,\alpha),(B,\beta)\}\quad\Longleftrightarrow\quad \alpha\overrightarrow{GA}+\beta\overrightarrow{GB}=\vec{0}.
\]
\end{propriete}

En effet, en appliquant la réduction au point $M=G$, on obtient $\alpha\overrightarrow{GA}+\beta\overrightarrow{GB}=(\alpha+\beta)\overrightarrow{GG}=\vec{0}$. Réciproquement, si un point $G$ vérifie $\alpha\overrightarrow{GA}+\beta\overrightarrow{GB}=\vec{0}$, alors pour tout point $M$ :
\[
\alpha\overrightarrow{MA}+\beta\overrightarrow{MB}=\alpha\overrightarrow{MG}+\beta\overrightarrow{MG}+\alpha\overrightarrow{GA}+\beta\overrightarrow{GB}=(\alpha+\beta)\overrightarrow{MG},
\]
ce qui prouve que $G$ joue bien le rôle du barycentre.

\begin{aretenir}[Les trois écritures équivalentes]
Pour $\alpha+\beta\neq 0$, les trois affirmations suivantes sont équivalentes :
\begin{enumerate}
\item $G=\text{bar}\{(A,\alpha),(B,\beta)\}$ ;
\item $\alpha\overrightarrow{GA}+\beta\overrightarrow{GB}=\vec{0}$ (caractérisation vectorielle) ;
\item $\overrightarrow{AG}=\dfrac{\beta}{\alpha+\beta}\overrightarrow{AB}$ (position de $G$ sur la droite $(AB)$).
\end{enumerate}
On choisit celle qui est la plus adaptée au problème : la (2) pour les calculs vectoriels, la (3) pour construire $G$.
\end{aretenir}

\courssub{Invariance par isométrie}

\begin{propriete}[Conservation par les isométries (admise)]
Le barycentre est conservé par les transformations usuelles du plan : \cle{translation}, \cle{rotation}, \cle{symétrie} axiale ou centrale (et plus généralement par toute isométrie). Autrement dit, si $G=\text{bar}\{(A,\alpha),(B,\beta)\}$ et si une isométrie envoie $A$ sur $A'$, $B$ sur $B'$ et $G$ sur $G'$, alors :
\[
G'=\text{bar}\{(A',\alpha),(B',\beta)\}.
\]
Ce résultat est admis ; on l'illustre ci-dessous.
\end{propriete}

\begin{exemplebox}[Illustration avec une symétrie]
Soit $G=\text{bar}\{(A,1),(B,2)\}$. Effectuons la symétrie d'axe $(\Delta)$ : $A$ devient $A'$, $B$ devient $B'$, $G$ devient $G'$. La symétrie conserve les distances, donc $\dfrac{A'G'}{G'B'}=\dfrac{AG}{GB}=\dfrac{2}{1}$ : le point $G'$ est bien le barycentre de $(A',1)$ et $(B',2)$. Concrètement, si une pièce métallique homogène est découpée puis retournée, son centre de gravité « suit » le mouvement.
\end{exemplebox}

\begin{savaistu}[Le barycentre, outil des métiers techniques]
La conservation du barycentre par isométrie est utilisée en mécanique : le \emph{centre de gravité} d'une pièce est un barycentre de points pondérés par des masses. Quand on déplace ou fait pivoter la pièce (translation, rotation), le centre de gravité subit le même mouvement. C'est ce principe qui permet, par exemple, de prévoir l'équilibre d'une grue de chantier ou la stabilité d'un véhicule, très étudiés au lycée technique.
\end{savaistu}

\courssub{Utilisation de la réduction pour déterminer des ensembles de points}

Voici l'application la plus importante de la réduction, qui prépare la section sur les lignes de niveau. Le principe : une condition portant sur un point $M$ \emph{inconnu} et faisant intervenir $A$ et $B$ se transforme, grâce au barycentre, en une condition ne faisant intervenir que $M$ et $G$ — beaucoup plus facile à interpréter géométriquement.

\begin{methode}[Déterminer un ensemble de points grâce à la réduction]
Pour trouver l'ensemble des points $M$ vérifiant une condition sur $\alpha\overrightarrow{MA}+\beta\overrightarrow{MB}$ :
\begin{enumerate}
\item Vérifier que $\alpha+\beta\neq 0$ (sinon, la somme est un vecteur constant et la condition ne dépend plus de $M$).
\item Introduire le barycentre $G=\text{bar}\{(A,\alpha),(B,\beta)\}$ et le placer sur la figure (ou le construire avec $\overrightarrow{AG}=\dfrac{\beta}{\alpha+\beta}\overrightarrow{AB}$).
\item Remplacer : $\alpha\overrightarrow{MA}+\beta\overrightarrow{MB}=(\alpha+\beta)\overrightarrow{MG}$.
\item La condition devient une relation simple sur $\overrightarrow{MG}$ ou sur la distance $MG$ : on reconnaît alors un cercle, une droite, un segment, etc.
\item Conclure en décrivant précisément l'ensemble (centre, rayon, position).
\end{enumerate}
\end{methode}

\begin{exoresolu}[Un premier ensemble de points]
Le plan est muni de deux points $A$ et $B$ tels que $AB=\SI{6}{cm}$. Déterminer l'ensemble $(\mathcal{E})$ des points $M$ tels que :
\[
\left\|\,2\overrightarrow{MA}+3\overrightarrow{MB}\,\right\|=\SI{10}{cm}.
\]
\tcblower
\textbf{Étape 1.} La somme des coefficients est $2+3=5\neq 0$ : on peut introduire le barycentre.

\textbf{Étape 2.} Soit $G=\text{bar}\{(A,2),(B,3)\}$. On a $\overrightarrow{AG}=\dfrac{3}{5}\overrightarrow{AB}$, donc $AG=\dfrac{3}{5}\times \SI{6}{cm}=\SI{3,6}{cm}$ : $G$ est le point du segment $[AB]$ situé à \SI{3,6}{cm} de $A$.

\textbf{Étape 3.} Par la propriété de réduction, pour tout point $M$ :
\[
2\overrightarrow{MA}+3\overrightarrow{MB}=5\overrightarrow{MG}.
\]

\textbf{Étape 4.} La condition devient :
\[
\left\|5\overrightarrow{MG}\right\|=\SI{10}{cm}\quad\Longleftrightarrow\quad 5\,MG=\SI{10}{cm}\quad\Longleftrightarrow\quad MG=\SI{2}{cm}.
\]

\textbf{Étape 5.} L'ensemble $(\mathcal{E})$ est donc le \textbf{cercle de centre $G$ et de rayon \SI{2}{cm}}. Le problème, qui faisait intervenir deux points et deux vecteurs, s'est ramené à une simple distance à un point fixe.
\end{exoresolu}

\begin{attention}[Pièges fréquents dans les ensembles de points]
\begin{itemize}
\item \textbf{Oublier le facteur $(\alpha+\beta)$} : $\left\|2\overrightarrow{MA}+3\overrightarrow{MB}\right\|=\SI{10}{cm}$ donne $5\,MG=\SI{10}{cm}$ et non $MG=\SI{10}{cm}$. Le rayon est \SI{2}{cm}, pas \SI{10}{cm}.
\item \textbf{Se tromper de barycentre} : pour réduire $2\overrightarrow{MA}+3\overrightarrow{MB}$, le barycentre est $\text{bar}\{(A,2),(B,3)\}$ : le coefficient de $\overrightarrow{MA}$ va avec le point $A$.
\item \textbf{Appliquer la réduction quand $\alpha+\beta=0$} : dans ce cas, la somme est constante et l'ensemble cherché est soit le plan tout entier, soit l'ensemble vide.
\end{itemize}
\end{attention}

\begin{exercice}[À faire]
Soient $A$ et $B$ deux points tels que $AB=\SI{4}{cm}$.
\begin{enumerate}
\item Construire le point $G=\text{bar}\{(A,1),(B,3)\}$.
\item Montrer que pour tout point $M$, $\overrightarrow{MA}+3\overrightarrow{MB}=4\overrightarrow{MG}$.
\item En déduire l'ensemble des points $M$ tels que $\left\|\overrightarrow{MA}+3\overrightarrow{MB}\right\|=\SI{8}{cm}$.
\item Que dire de l'ensemble des points $M$ tels que $\overrightarrow{MA}-\overrightarrow{MB}=\overrightarrow{BA}$ ?
\end{enumerate}
\end{exercice}

\begin{corrige}[Exercice]
\begin{enumerate}
\item $\overrightarrow{AG}=\dfrac{3}{4}\overrightarrow{AB}$, donc $G\in[AB]$ avec $AG=\SI{3}{cm}$.
\item La somme des coefficients vaut $1+3=4\neq 0$ : la réduction donne $\overrightarrow{MA}+3\overrightarrow{MB}=4\overrightarrow{MG}$.
\item La condition équivaut à $4\,MG=\SI{8}{cm}$, soit $MG=\SI{2}{cm}$ : c'est le cercle de centre $G$ et de rayon \SI{2}{cm}.
\item La somme des coefficients est $1-1=0$ : pour tout point $M$, $\overrightarrow{MA}-\overrightarrow{MB}=\overrightarrow{BA}$ (vecteur constant). La condition est donc vérifiée par \emph{tous} les points du plan.
\end{enumerate}
\end{corrige}

\begin{aretenir}[Bilan de la section]
\begin{itemize}
\item \textbf{Homogénéité} : $\text{bar}\{(A,\alpha),(B,\beta)\}=\text{bar}\{(A,k\alpha),(B,k\beta)\}$ pour tout $k\neq 0$.
\item \textbf{Réduction} : si $\alpha+\beta\neq 0$ et $G=\text{bar}\{(A,\alpha),(B,\beta)\}$, alors pour tout point $M$ : $\alpha\overrightarrow{MA}+\beta\overrightarrow{MB}=(\alpha+\beta)\overrightarrow{MG}$.
\item \textbf{Somme nulle} : si $\alpha+\beta=0$, alors $\alpha\overrightarrow{MA}+\beta\overrightarrow{MB}=\alpha\overrightarrow{BA}$ est constant.
\item \textbf{Caractérisation} : $G=\text{bar}\{(A,\alpha),(B,\beta)\}\Longleftrightarrow\alpha\overrightarrow{GA}+\beta\overrightarrow{GB}=\vec{0}$.
\item \textbf{Isométries} : translations, rotations et symétries conservent le barycentre.
\end{itemize}
Ces propriétés s'étendront naturellement à trois et quatre points dans les sections suivantes, et la réduction sera l'outil central de l'étude des lignes de niveau.
\end{aretenir}

\coursec{Barycentre de trois points pondérés}

Après avoir maîtrisé le \cle{barycentre} de deux points et ses propriétés vectorielles fondamentales, nous étendons naturellement la définition à trois points pondérés. Cette généralisation n'est pas seulement théorique : elle permet de modéliser le centre de gravité d'un système de trois masses, de localiser le centre de gravité d'un triangle (\cle{isobarycentre}) et de résoudre des problèmes de mécanique ou de géométrie analytique en réduisant le problème à une succession de barycentres à deux points.

\courssub{Définition, existence et coordonnées}

\begin{definition}[Barycentre de trois points pondérés]
Soient $A$, $B$, $C$ trois points du plan et $\alpha$, $\beta$, $\gamma$ trois nombres réels \textbf{non nuls} tels que $\alpha+\beta+\gamma \neq 0$. On appelle \textbf{barycentre} du système pondéré $\{(A,\alpha),\,(B,\beta),\,(C,\gamma)\}$ l'unique point $G$ vérifiant :
\[
\alpha\,\overrightarrow{GA} + \beta\,\overrightarrow{GB} + \gamma\,\overrightarrow{GC} = \vec{0}.
\]
On note $G = \operatorname{bar}\{(A,\alpha),\,(B,\beta),\,(C,\gamma)\}$.
\end{definition}

\begin{propriete}[Existence, unicité et expression vectorielle]
Pour tout triplet de coefficients $\alpha,\beta,\gamma$ non nuls de somme non nulle, le barycentre $G$ existe et est unique. De plus, pour tout point $O$ du plan :
\[
\overrightarrow{OG} = \frac{\alpha\,\overrightarrow{OA} + \beta\,\overrightarrow{OB} + \gamma\,\overrightarrow{OC}}{\alpha+\beta+\gamma}.
\]
En particulier, en choisissant $O=A$, on obtient la relation fondamentale (souvent utilisée pour les constructions) :
\[
\overrightarrow{AG} = \frac{\beta\,\overrightarrow{AB} + \gamma\,\overrightarrow{AC}}{\alpha+\beta+\gamma}.
\]
\end{propriete}

\begin{experience}[Démonstration par réduction à deux points]
On cherche $G$ tel que $\alpha\vec{GA}+\beta\vec{GB}+\gamma\vec{GC}=\vec{0}$.
\begin{enumerate}
\item On isole les termes en $A$ et $B$ : $\alpha\vec{GA}+\beta\vec{GB} = -\,\gamma\vec{GC}$.
\item On introduit le barycentre partiel $G'$ de $\{(A,\alpha),(B,\beta)\}$ (existe car $\alpha+\beta \neq 0$ ; si $\alpha+\beta=0$, voir la propriété d'associativité plus loin).
\item Par définition de $G'$ : $\alpha\vec{G'A}+\beta\vec{G'B}=\vec{0}$.
   En écrivant $\vec{GA} = \vec{GG'} + \vec{G'A}$ et $\vec{GB} = \vec{GG'} + \vec{G'B}$ (relation de Chasles), on obtient :
   $\alpha\vec{GA}+\beta\vec{GB} = (\alpha+\beta)\vec{GG'} + \underbrace{(\alpha\vec{G'A}+\beta\vec{G'B})}_{=\vec{0}} = (\alpha+\beta)\vec{GG'}$.
\item L'équation initiale devient $(\alpha+\beta)\vec{GG'} + \gamma\vec{GC} = \vec{0}$.
\item Or $\vec{GC} = \vec{GG'} + \vec{G'C}$. Donc $(\alpha+\beta)\vec{GG'} + \gamma(\vec{GG'}+\vec{G'C}) = \vec{0}$.
\item Soit $(\alpha+\beta+\gamma)\vec{GG'} = -\gamma\vec{G'C}$, d'où $\vec{GG'} = \frac{-\gamma}{\alpha+\beta+\gamma}\vec{G'C} = \frac{\gamma}{\alpha+\beta+\gamma}\vec{CG'}$.
\item Cela montre que $G$ est le barycentre de $\{(G',\alpha+\beta),\,(C,\gamma)\}$. Il existe, est unique et appartient à la droite $(G'C)$.
\end{enumerate}
\end{experience}

\begin{propriete}[Coordonnées du barycentre dans un repère]
Dans un repère orthonormé $(O;\vec{\imath},\vec{\jmath})$, si $A(x_A,y_A)$, $B(x_B,y_B)$, $C(x_C,y_C)$, alors les coordonnées de $G = \operatorname{bar}\{(A,\alpha),(B,\beta),(C,\gamma)\}$ sont les \textbf{moyennes pondérées} des coordonnées :
\[
x_G = \frac{\alpha x_A + \beta x_B + \gamma x_C}{\alpha+\beta+\gamma}, \qquad
y_G = \frac{\alpha y_A + \beta y_B + \gamma y_C}{\alpha+\beta+\gamma}.
\]
\end{propriete}

\courssub{Position du barycentre et cas particulier : l'isobarycentre}

\begin{propriete}[Position de $G$ dans le plan]
Le barycentre $G$ de trois points $A,B,C$ appartient toujours au plan $(ABC)$.
\begin{itemize}
\item Si les trois coefficients $\alpha,\beta,\gamma$ sont de \textbf{même signe} (tous positifs ou tous négatifs) \textbf{et si $A,B,C$ ne sont pas alignés}, alors $G$ se trouve \textbf{à l'intérieur du triangle $ABC$} (plus précisément dans l'intérieur de l'enveloppe convexe).
\item Si les coefficients sont de signes différents, $G$ se trouve à l'extérieur du triangle (éventuellement à l'infini si $\alpha+\beta+\gamma=0$, cas exclu par la définition).
\end{itemize}
\end{propriete}

\begin{propriete}[Isobarycentre et centre de gravité du triangle]
L'\textbf{isobarycentre} de trois points $A,B,C$ est le barycentre à coefficients égaux : $G = \operatorname{bar}\{(A,1),(B,1),(C,1)\}$.
Il coïncide avec le \textbf{centre de gravité (centroïde)} du triangle $ABC$, c'est-à-dire l'intersection des trois médianes. Ses coordonnées sont les moyennes arithmétiques :
\[
x_G = \frac{x_A+x_B+x_C}{3}, \qquad y_G = \frac{y_A+y_B+y_C}{3}.
\]
\emph{Preuve exigible au Probatoire :} Par associativité (voir plus bas), $G$ est le barycentre de $\{(A,1),(B,1)\}$ (milieu $I$ de $[AB]$) et $(C,1)$. Donc $G$ appartient à la médiane $[CI]$ et $\vec{CG} = \frac{2}{3}\vec{CI}$. Par symétrie, $G$ appartient aux trois médianes.
\end{propriete}

\begin{popfigure}
\begin{tikzpicture}[scale=1.1, >=Stealth]
% Triangle
\coordinate (A) at (0,0);
\coordinate (B) at (5,0);
\coordinate (C) at (1.5,4);
\draw[popblue, thick] (A) -- (B) -- (C) -- cycle;
% Medians
\coordinate (I) at ($(A)!0.5!(B)$); % Midpoint AB
\coordinate (J) at ($(B)!0.5!(C)$); % Midpoint BC
\coordinate (K) at ($(C)!0.5!(A)$); % Midpoint CA
\draw[poporange, dashed] (C) -- (I);
\draw[poporange, dashed] (A) -- (J);
\draw[poporange, dashed] (B) -- (K);
% Centroid G
\coordinate (G) at ({(0+5+1.5)/3}, {(0+0+4)/3});
\fill[poppink] (G) circle (2.5pt);
% Labels
\node[below left] at (A) {$A$};
\node[below right] at (B) {$B$};
\node[above] at (C) {$C$};
\node[below] at (I) {$I$};
\node[right] at (J) {$J$};
\node[left] at (K) {$K$};
\node[above right] at (G) {$G$};
% Vectors 2/3
\draw[->, popgreen, thick] (C) -- ($(C)!0.6666!(I)$) node[midway, left] {$\frac{2}{3}\vec{CI}$};
\draw[->, popgreen, thick] (A) -- ($(A)!0.6666!(J)$) node[midway, above] {$\frac{2}{3}\vec{AJ}$};
\end{tikzpicture}
\legende{Isobarycentre $G$ (centre de gravité) : intersection des médianes $[CI], [AJ], [BK]$. $G$ est aux deux tiers de chaque médiane en partant du sommet.}
\end{popfigure}

\courssub{Associativité et construction par barycentres partiels}

C'est la propriété la plus puissante pour le calcul et la construction géométrique : on peut « regrouper » les points deux à deux.

\begin{propriete}[Théorème : Associativité du barycentre — Règle des barycentres partiels]
Soient $A,B,C$ trois points et $\alpha,\beta,\gamma$ trois réels non nuls tels que $\alpha+\beta+\gamma \neq 0$.
Si $\alpha+\beta \neq 0$ et si l'on pose $G' = \operatorname{bar}\{(A,\alpha),(B,\beta)\}$, alors :
\[
G = \operatorname{bar}\{(A,\alpha),(B,\beta),(C,\gamma)\} = \operatorname{bar}\{(G',\alpha+\beta),\,(C,\gamma)\}.
\]
De même, si $\beta+\gamma \neq 0$ ou $\gamma+\alpha \neq 0$, on peut regrouper les autres paires.
\end{propriete}

\begin{attention}[Condition cruciale sur les sommes partielles]
L'associativité ne s'applique \textbf{que si la somme des coefficients du groupe regroupé est non nulle}.
\begin{itemize}
\item Si $\alpha+\beta = 0$, le point $G' = \operatorname{bar}\{(A,\alpha),(B,\beta)\}$ \textbf{n'existe pas} (barycentre à somme nulle = point à l'infini / direction). On ne peut pas faire cette réduction-là.
\item Il faut alors regrouper $(B,C)$ si $\beta+\gamma\neq 0$, ou $(C,A)$ si $\gamma+\alpha\neq 0$. Au moins une de ces trois sommes est non nulle car $(\alpha+\beta)+(\beta+\gamma)+(\gamma+\alpha)=2(\alpha+\beta+\gamma)\neq 0$.
\end{itemize}
\textbf{Piège fréquent :} Appliquer la formule $\vec{AG} = \frac{\beta\vec{AB}+\gamma\vec{AC}}{\alpha+\beta+\gamma}$ sans vérifier que le dénominateur n'est pas nul (déjà exclu par définition) \textbf{et} sans vérifier que le regroupement choisi pour la construction géométrique est valide ($\alpha+\beta\neq 0$).
\end{attention}

\begin{methode}[Construction géométrique de $G = \operatorname{bar}\{(A,\alpha),(B,\beta),(C,\gamma)\}$]
\emph{Supposons $\alpha+\beta \neq 0$ (sinon permuter les points).}
\begin{enumerate}
\item Construire le barycentre partiel $G' = \operatorname{bar}\{(A,\alpha),(B,\beta)\}$.
      \begin{itemize}
      \item Si $\alpha$ et $\beta$ même signe : $G' \in [AB]$ avec $\frac{G'A}{G'B} = \frac{\beta}{\alpha}$ (rapport signé).
      \item Si $\alpha$ et $\beta$ signes contraires : $G'$ est à l'extérieur de $[AB]$ (division extérieure).
      \end{itemize}
\item Construire $G = \operatorname{bar}\{(G',\alpha+\beta),\,(C,\gamma)\}$.
      \item $G$ appartient au segment $[G'C]$ si $(\alpha+\beta)$ et $\gamma$ même signe ; sinon $G$ est sur la droite $(G'C)$ à l'extérieur du segment.
      \item Le rapport est $\frac{GG'}{GC} = \frac{\gamma}{\alpha+\beta}$ (signes respectés).
\end{enumerate}
\end{methode}

\begin{popfigure}
\begin{tikzpicture}[scale=1.2, >=Stealth]
% Points
\coordinate (A) at (0,0);
\coordinate (B) at (5,0);
\coordinate (C) at (1.5,4);
% Coefficients example: alpha=1, beta=2, gamma=3 -> alpha+beta=3
% G' = bar{(A,1),(B,2)} -> G' divides [AB] in ratio 2:1 (AG'/G'B = 2/1) -> G' at 2/3 from A.
\coordinate (Gp) at ($(A)!0.66666!(B)$); 
% G = bar{(G', 3), (C, 3)} -> Isobarycentre of G' and C -> Midpoint of [G'C]
\coordinate (G) at ($(Gp)!0.5!(C)$);

% Triangle
\draw[popblue, thick] (A) -- (B) -- (C) -- cycle;
% Construction lines
\draw[poporange, dashed] (C) -- (Gp);
\draw[popgreen, thick] (A) -- (Gp);
\draw[popgreen, thick] (Gp) -- (B);
% Points
\fill[popblue] (A) circle (2pt) node[below left] {$A(\alpha)$};
\fill[popblue] (B) circle (2pt) node[below right] {$B(\beta)$};
\fill[popblue] (C) circle (2pt) node[above] {$C(\gamma)$};
\fill[poporange] (Gp) circle (2.5pt) node[below] {$G'(\alpha+\beta)$};
\fill[poppink] (G) circle (2.5pt) node[above right] {$G$};
% Ratio labels
\node[popgreen] at ($(A)!0.3!(Gp)$) {$\beta$};
\node[popgreen] at ($(Gp)!0.3!(B)$) {$\alpha$};
\node[poporange] at ($(Gp)!0.5!(C)$) [right] {$\gamma$};
\node[poporange] at ($(G)!0.5!(Gp)$) [left] {$\alpha+\beta$};
\end{tikzpicture}
\legende{Construction par barycentres partiels : $G'$ barycentre de $(A,\alpha)$ et $(B,\beta)$, puis $G$ barycentre de $(G',\alpha+\beta)$ et $(C,\gamma)$. Ici $\alpha,\beta,\gamma>0$, donc $G'$ dans $[AB]$ et $G$ dans $[G'C]$.}
\end{popfigure}

\courssub{Exemple chiffré complet}

\begin{exoresolu}[Construction et coordonnées de $G = \operatorname{bar}\{(A,1),(B,2),(C,3)\}$]
Soient $A(0;0)$, $B(6;0)$, $C(0;4)$ dans un repère orthonormé. Construire $G$ et déterminer ses coordonnées.

\textbf{1. Vérification des conditions}
$\alpha=1$, $\beta=2$, $\gamma=3$. Somme $S = 1+2+3 = 6 \neq 0$. $\alpha+\beta = 3 \neq 0$. On peut regrouper $A$ et $B$ en premier.

\textbf{2. Calcul des coordonnées (méthode analytique)}
\begin{align*}
x_G &= \frac{1\times 0 + 2\times 6 + 3\times 0}{6} = \frac{12}{6} = 2, \\
y_G &= \frac{1\times 0 + 2\times 0 + 3\times 4}{6} = \frac{12}{6} = 2.
\end{align*}
Donc $G(2;2)$.

\textbf{3. Construction géométrique par barycentres partiels (méthode synthétique)}
\begin{itemize}
\item \textbf{Étape 1 :} $G' = \operatorname{bar}\{(A,1),(B,2)\}$.
      Coefficients de même signe ($>0$) $\Rightarrow G' \in [AB]$.
      Rapport : $\frac{G'A}{G'B} = \frac{\beta}{\alpha} = \frac{2}{1}$. $G'$ est aux $\frac{2}{3}$ de $A$ vers $B$ (plus près de $B$ qui a le poids 2).
      Coordonnées : $G'\left(\frac{1\cdot0+2\cdot6}{3}; \frac{1\cdot0+2\cdot0}{3}\right) = G'(4;0)$.
\item \textbf{Étape 2 :} $G = \operatorname{bar}\{(G',3),\,(C,3)\}$.
      Coefficients égaux ($3$ et $3$) $\Rightarrow G$ est le \textbf{milieu} de $[G'C]$.
      $G'\left(4;0\right)$, $C\left(0;4\right) \Rightarrow G\left(\frac{4+0}{2}; \frac{0+4}{2}\right) = G(2;2)$.
\end{itemize}
Les deux méthodes donnent le même résultat : \textbf{$G(2;2)$}.

\textbf{4. Position relative}
Tous les coefficients sont positifs et $A,B,C$ non alignés $\Rightarrow G$ est à l'intérieur du triangle $ABC$. Vérification : $A(0,0), B(6,0), C(0,4)$. Le point $(2,2)$ satisfait $x>0, y>0, \frac{x}{6}+\frac{y}{4}=\frac{1}{3}+\frac{1}{2}=\frac{5}{6}<1$, il est bien dans le triangle rectangle.
\end{exoresolu}

\begin{savaistu}[Le centre de gravité en technique]
En mécanique et en construction, le centre de gravité d'une plaque homogène triangulaire (ou d'une poutre en treillis triangulaire) est exactement l'isobarycentre de ses trois sommets. C'est le point où l'on peut poser un support unique pour équilibrer la structure. Les grues de chantier utilisent ce principe : le crochet doit être aligné verticalement avec le centre de gravité de la charge pour éviter le basculement. Plus généralement, le barycentre à trois coefficients permet de déterminer le point d'application de la résultante de trois forces parallèles (ex: réactions d'appuis sur une poutre portée en trois points).
\end{savaistu}

\coursec{Barycentre de quatre points et applications géométriques}

Dans la section précédente, nous avons appris à construire et à utiliser le barycentre de trois points pondérés. Nous avons vu en particulier que le centre de gravité d'un triangle est l'isobarycentre de ses trois sommets. Mais dans les ateliers et sur les chantiers, les objets ont rarement trois points remarquables : une plaque rectangulaire a quatre coins, une poutre est soutenue en quatre points, un châssis de machine repose sur quatre pieds. Il est donc naturel d'étendre la notion de barycentre à \textbf{quatre points pondérés}. Bonne nouvelle : toutes les propriétés déjà rencontrées (existence, unicité, homogénéité, associativité) se généralisent sans difficulté. C'est l'associativité qui fera tout le travail et qui nous permettra de démontrer élégamment des théorèmes classiques de géométrie.

\courssub{Définition du barycentre de quatre points pondérés}

L'intuition est la même qu'avec deux ou trois points : on imagine quatre masses posées en quatre points d'une plaque, et l'on cherche le point d'équilibre, le point où la plaque « tient » sur le doigt.

\begin{definition}[Barycentre de quatre points pondérés]
Soient $A$, $B$, $C$, $D$ quatre points du plan et $\alpha$, $\beta$, $\gamma$, $\delta$ quatre réels tels que
\[
\alpha+\beta+\gamma+\delta \neq 0.
\]
Il existe un unique point $G$ tel que
\[
\alpha\overrightarrow{GA}+\beta\overrightarrow{GB}+\gamma\overrightarrow{GC}+\delta\overrightarrow{GD}=\vec{0}.
\]
Ce point $G$ est appelé \cle{barycentre} du système de points pondérés $(A,\alpha)$, $(B,\beta)$, $(C,\gamma)$, $(D,\delta)$. On note
\[
G=\mathrm{bar}\{(A,\alpha),(B,\beta),(C,\gamma),(D,\delta)\}.
\]
Si $\alpha=\beta=\gamma=\delta$, on dit que $G$ est l'\cle{isobarycentre} des quatre points.
\end{definition}

Comme pour trois points, on dispose d'une formule de \emph{réduction} très pratique pour les calculs : pour tout point $M$ du plan,
\[
\alpha\overrightarrow{MA}+\beta\overrightarrow{MB}+\gamma\overrightarrow{MC}+\delta\overrightarrow{MD}=(\alpha+\beta+\gamma+\delta)\overrightarrow{MG}.
\]
En choisissant $M=A$ (ou l'origine d'un repère), on obtient une expression de $\overrightarrow{AG}$ en fonction de $\overrightarrow{AB}$, $\overrightarrow{AC}$ et $\overrightarrow{AD}$, ce qui permet de placer $G$.

\begin{propriete}[Propriétés généralisées (admises)]
Les propriétés vues pour deux et trois points restent vraies pour quatre points, avec des démonstrations identiques en principe :
\begin{itemize}
\item \textbf{Existence et unicité} : si $\alpha+\beta+\gamma+\delta\neq 0$, le barycentre existe et est unique.
\item \textbf{Homogénéité} : on ne change pas le barycentre en multipliant tous les coefficients par un même réel non nul $k$ :
\[
\mathrm{bar}\{(A,\alpha),(B,\beta),(C,\gamma),(D,\delta)\}=\mathrm{bar}\{(A,k\alpha),(B,k\beta),(C,k\gamma),(D,k\delta)\}.
\]
\item \textbf{Associativité} : on peut regrouper les points en sous-systèmes. Par exemple, si $G_1=\mathrm{bar}\{(A,\alpha),(B,\beta)\}$ avec $\alpha+\beta\neq 0$ et $G_2=\mathrm{bar}\{(C,\gamma),(D,\delta)\}$ avec $\gamma+\delta\neq 0$, alors
\[
G=\mathrm{bar}\{(G_1,\alpha+\beta),(G_2,\gamma+\delta)\}.
\]
On peut aussi remplacer trois des points par leur barycentre partiel affecté de la somme de leurs coefficients.
\end{itemize}
\end{propriete}

\begin{attention}[La condition sur la somme des coefficients]
La condition $\alpha+\beta+\gamma+\delta\neq 0$ est \textbf{indispensable}. Si la somme est nulle, le vecteur $\alpha\overrightarrow{MA}+\beta\overrightarrow{MB}+\gamma\overrightarrow{MC}+\delta\overrightarrow{MD}$ est un vecteur constant non nul : il n'existe aucun point $G$ vérifiant la définition. Avant tout calcul de barycentre, \textbf{commence toujours par vérifier que la somme des coefficients est non nulle}, et signale-le dans ta rédaction : c'est un point de justification attendu au Probatoire.
\end{attention}

\courssub{Cas particulier fondamental : l'isobarycentre des sommets d'un parallélogramme}

Voici le résultat géométrique le plus important de cette section. Il justifie à lui seul l'intérêt des barycentres.

\begin{propriete}[Isobarycentre des sommets d'un parallélogramme]
Soit $ABCD$ un parallélogramme. L'isobarycentre des quatre sommets $(A,1)$, $(B,1)$, $(C,1)$, $(D,1)$ est le \cle{point d'intersection des diagonales} $[AC]$ et $[BD]$, c'est-à-dire le centre du parallélogramme.
\end{propriete}

\textbf{Démonstration.} Notons $O$ le milieu de $[AC]$. Comme $ABCD$ est un parallélogramme, ses diagonales se coupent en leur milieu, donc $O$ est aussi le milieu de $[BD]$. Utilisons l'associativité :
\begin{itemize}
\item $O=\mathrm{bar}\{(A,1),(C,1)\}$ car $O$ est le milieu de $[AC]$ ;
\item $O=\mathrm{bar}\{(B,1),(D,1)\}$ car $O$ est le milieu de $[BD]$.
\end{itemize}
Par associativité, l'isobarycentre $G$ de $(A,1),(B,1),(C,1),(D,1)$ vérifie
\[
G=\mathrm{bar}\{(O,2),(O,2)\}=O.
\]
Donc $G=O$ : l'isobarycentre des quatre sommets est bien l'intersection des diagonales. $\blacksquare$

\begin{exemplebox}[Démonstration : les diagonales d'un parallélogramme se coupent en leur milieu]
Les barycentres permettent de \emph{démontrer} cette propriété classique sans la supposer. Soit $ABCD$ un parallélogramme, c'est-à-dire tel que $\overrightarrow{AB}=\overrightarrow{DC}$. Posons $G=\mathrm{bar}\{(A,1),(B,1),(C,1),(D,1)\}$ (la somme des coefficients vaut $4\neq 0$).
\begin{itemize}
\item Par associativité, en regroupant $(A,1)$ avec $(C,1)$ et $(B,1)$ avec $(D,1)$, on introduit $I$ milieu de $[AC]$ et $J$ milieu de $[BD]$. On a $G=\mathrm{bar}\{(I,2),(J,2)\}$, donc $G$ est le milieu de $[IJ]$.
\item Montrons que $I=J$. On a $\overrightarrow{IJ}=\overrightarrow{IA}+\overrightarrow{AJ}$. Or $\overrightarrow{IA}=\frac{1}{2}\overrightarrow{CA}$ et $\overrightarrow{AJ}=\frac{1}{2}\overrightarrow{AD}$ (car $J$ milieu de $[BD]$ et $ABCD$ parallélogramme $\Rightarrow$ $J$ milieu de $[BD]$ mais il faut relier $A$ et $J$...). 
Mieux : $\overrightarrow{IJ} = \frac{1}{2}(\overrightarrow{AB}+\overrightarrow{CD})$ (relation de Chasles et milieux). Comme $\overrightarrow{CD}=-\overrightarrow{DC}=-\overrightarrow{AB}$, on obtient $\overrightarrow{IJ}=\frac{1}{2}(\overrightarrow{AB}-\overrightarrow{AB})=\vec{0}$.
\end{itemize}
Donc $I=J$. Les diagonales ont le même milieu : elles se coupent en leur milieu commun, qui est l'isobarycentre $G$ des quatre sommets.
\end{exemplebox}

\begin{popfigure}
\begin{tikzpicture}[scale=1.5]
\coordinate (A) at (0,0);
\coordinate (B) at (3.2,0);
\coordinate (D) at (1,1.8);
\coordinate (C) at ($(B)+(D)$);
\draw[thick,popblue] (A)--(B)--(C)--(D)--cycle;
\draw[poporange,thick] (A)--(C);
\draw[poppurple,thick] (B)--(D);
\coordinate (G) at ($(A)!0.5!(C)$);
\fill[popink] (G) circle (1.6pt) node[below right=1pt] {$G$};
\foreach \p/\pos in {A/below left,B/below right,C/above right,D/above left}{
  \fill[popblue] (\p) circle (1.4pt) node[\pos] {$\p$};
}
\node[poporange] at ($(A)!0.5!(G)+(0,-0.28)$) {\footnotesize diagonale $[AC]$};
\node[poppurple] at ($(B)!0.5!(G)+(0.15,0.25)$) {\footnotesize $[BD]$};
\end{tikzpicture}
\legende{Parallélogramme $ABCD$ : l'isobarycentre $G$ des quatre sommets est l'intersection des diagonales, milieu commun de $[AC]$ et $[BD]$.}
\end{popfigure}

\courssub{Application 1 : démontrer que trois points sont alignés}

L'idée est simple mais puissante : si un point $P$ est le barycentre de deux points $Q$ et $R$ (avec des coefficients de somme non nulle), alors $P$ appartient à la droite $(QR)$. En effet, la relation $\beta\overrightarrow{PQ}+\gamma\overrightarrow{PR}=\vec{0}$ montre que $\overrightarrow{PQ}$ et $\overrightarrow{PR}$ sont colinéaires.

\begin{methode}[Prouver un alignement par les barycentres]
Pour montrer que trois points $P$, $Q$, $R$ sont alignés :
\begin{enumerate}
\item Exprimer l'un des points, par exemple $P$, comme barycentre des deux autres : $P=\mathrm{bar}\{(Q,\beta),(R,\gamma)\}$ avec $\beta+\gamma\neq 0$.
\item Pour cela, partir de la définition d'un barycentre global (souvent un barycentre de trois ou quatre points) et utiliser l'\textbf{associativité} pour faire apparaître $P$ comme barycentre partiel.
\item Conclure : $P\in(QR)$, donc $P$, $Q$, $R$ sont alignés.
\end{enumerate}
\end{methode}

\begin{exoresolu}[Alignement de trois points]
\textbf{Énoncé.} Soit $ABC$ un triangle. On définit $G=\mathrm{bar}\{(A,2),(B,1),(C,1)\}$ et $H=\mathrm{bar}\{(A,4),(B,-1),(C,-1)\}$. Montrer que les points $A$, $G$, $H$ sont alignés.
\tcblower
\textbf{Solution.} Vérifions d'abord les sommes : pour $G$, $2+1+1=4\neq 0$ ; pour $H$, $4-1-1=2\neq 0$. Les deux barycentres existent.

Notons $A'$ le milieu de $[BC]$. On a $A'=\mathrm{bar}\{(B,1),(C,1)\}$.

\textbf{Pour $G$ :} par associativité, en regroupant $(B,1)$ et $(C,1)$ :
\[
G=\mathrm{bar}\{(A,2),(A',2)\}.
\]
$G$ est donc le barycentre de $A$ et $A'$ (coefficients $2$ et $2$), c'est-à-dire le \textbf{milieu de $[AA']$}. En particulier, $G\in(AA')$.

\textbf{Pour $H$ :} par associativité, en regroupant $(B,-1)$ et $(C,-1)$ :
\[
\mathrm{bar}\{(B,-1),(C,-1)\}=A' \quad\text{(coefficient $-2$)}.
\]
Donc $H=\mathrm{bar}\{(A,4),(A',-2)\}$.
$H$ est le barycentre de $A$ (coefficient $4$) et $A'$ (coefficient $-2$). La somme des coefficients est $2\neq 0$, donc $H$ existe et appartient à la droite $(AA')$.

\textbf{Conclusion :} $G$ et $H$ appartiennent tous deux à la droite $(AA')$ qui passe par $A$. Les points $A$, $G$, $H$ sont alignés.
\end{exoresolu}

\begin{attention}[Bien choisir le regroupement]
Dans un problème d'alignement, le regroupement par associativité doit faire apparaître \textbf{le même point pivot} (souvent un sommet commun ou un milieu) dans les deux expressions. Si tes coefficients ne permettent pas un tel regroupement, reviens au calcul vectoriel avec la formule de réduction : $\overrightarrow{AG}=\dfrac{1}{\alpha+\beta+\gamma+\delta}\left(\beta\overrightarrow{AB}+\gamma\overrightarrow{AC}+\delta\overrightarrow{AD}\right)$, puis cherche un coefficient de colinéarité.
\end{attention}

\courssub{Application 2 : démontrer que des droites sont concourantes}

C'est l'application la plus spectaculaire de l'associativité. Pour montrer que trois droites sont concourantes, on montre qu'un \textbf{même point} (un barycentre global) appartient à chacune des trois droites, en l'exprimant de trois façons différentes.

\begin{methode}[Prouver un concours de droites]
Pour montrer que trois droites $(d_1)$, $(d_2)$, $(d_3)$ sont concourantes :
\begin{enumerate}
\item Introduire un barycentre global $G$ de trois ou quatre points pondérés bien choisis.
\item Par associativité, écrire $G$ comme barycentre de deux points dont l'un est sur $(d_1)$ : alors $G\in(d_1)$.
\item Recommencer avec un autre regroupement pour obtenir $G\in(d_2)$, puis $G\in(d_3)$.
\item Conclure : $G$ est commun aux trois droites, qui sont donc concourantes en $G$.
\end{enumerate}
\end{methode}

\courssub{Application 3 : le théorème des médianes (démonstration guidée)}

Nous allons démontrer rigoureusement le théorème annoncé dans la section précédente. Cette démonstration est \textbf{formatrice et exigible} : elle illustre parfaitement la méthode du concours par associativité.

\begin{propriete}[Théorème des médianes — démonstration guidée en classe]
Dans un triangle $ABC$, les trois médianes sont concourantes en un point $G$ appelé \cle{centre de gravité} du triangle. De plus, si $A'$ est le milieu de $[BC]$, alors
\[
\overrightarrow{AG}=\frac{2}{3}\overrightarrow{AA'},
\]
c'est-à-dire que $G$ est situé aux \textbf{deux tiers} de chaque médiane en partant du sommet (partage $2/3$ -- $1/3$).
\end{propriete}

\begin{experience}[Démonstration guidée du théorème des médianes]
Soit $ABC$ un triangle. Notons $A'$, $B'$, $C'$ les milieux respectifs de $[BC]$, $[CA]$, $[AB]$. Les médianes sont les droites $(AA')$, $(BB')$, $(CC')$.

\textbf{Étape 1 : introduire le bon barycentre.} Considérons l'isobarycentre $G=\mathrm{bar}\{(A,1),(B,1),(C,1)\}$. La somme des coefficients vaut $3\neq 0$, donc $G$ existe et est unique.

\textbf{Étape 2 : premier regroupement.} Comme $A'$ est le milieu de $[BC]$, on a $A'=\mathrm{bar}\{(B,1),(C,1)\}$. Par associativité :
\[
G=\mathrm{bar}\{(A,1),(A',2)\}.
\]
La relation de définition donne $\overrightarrow{GA}+2\overrightarrow{GA'}=\vec{0}$, soit $\overrightarrow{AG}=2\overrightarrow{GA'}$. En écrivant $\overrightarrow{AA'}=\overrightarrow{AG}+\overrightarrow{GA'}=\overrightarrow{AG}+\tfrac{1}{2}\overrightarrow{AG}=\tfrac{3}{2}\overrightarrow{AG}$, on obtient
\[
\overrightarrow{AG}=\frac{2}{3}\overrightarrow{AA'}.
\]
En particulier $G\in(AA')$ : le point $G$ est sur la première médiane, aux deux tiers depuis $A$.

\textbf{Étape 3 : deuxième regroupement.} De même, $B'=\mathrm{bar}\{(C,1),(A,1)\}$, donc par associativité $G=\mathrm{bar}\{(B,1),(B',2)\}$, d'où $G\in(BB')$ et $\overrightarrow{BG}=\tfrac{2}{3}\overrightarrow{BB'}$.

\textbf{Étape 4 : troisième regroupement.} Enfin, $C'=\mathrm{bar}\{(A,1),(B,1)\}$ donne $G=\mathrm{bar}\{(C,1),(C',2)\}$, d'où $G\in(CC')$ et $\overrightarrow{CG}=\tfrac{2}{3}\overrightarrow{CC'}$.

\textbf{Conclusion.} Le point $G$ appartient aux trois médianes : elles sont concourantes en $G$, et $G$ partage chaque médiane dans le rapport $2/3$ -- $1/3$ à partir du sommet. $\blacksquare$
\end{experience}

\begin{popfigure}
\begin{tikzpicture}[scale=1.6]
\coordinate (A) at (0,2.6);
\coordinate (B) at (-1.6,0);
\coordinate (C) at (2,0);
\coordinate (A') at ($(B)!0.5!(C)$);
\coordinate (B') at ($(A)!0.5!(C)$);
\coordinate (C') at ($(A)!0.5!(B)$);
\draw[thick,popblue] (A)--(B)--(C)--cycle;
\draw[poporange,thick] (A)--(A');
\draw[poppurple,thick] (B)--(B');
\draw[popgreen,thick] (C)--(C');
\coordinate (G) at ($(A)!0.6667!(A')$);
\fill[popink] (G) circle (1.8pt) node[right=2pt] {$G$};
\fill[popblue] (A) circle (1.4pt) node[above] {$A$};
\fill[popblue] (B) circle (1.4pt) node[below left] {$B$};
\fill[popblue] (C) circle (1.4pt) node[below right] {$C$};
\fill[poporange] (A') circle (1.3pt) node[below] {$A'$};
\fill[poppurple] (B') circle (1.3pt) node[above right] {$B'$};
\fill[popgreen] (C') circle (1.3pt) node[above left] {$C'$};
\draw[<->,popink] ($(A)+(0.12,0)$)--($(G)+(0.12,0)$) node[midway,right] {\footnotesize $\frac{2}{3}$};
\draw[<->,popink] ($(G)+(0.12,0)$)--($(A')+(0.12,0)$) node[midway,right] {\footnotesize $\frac{1}{3}$};
\end{tikzpicture}
\legende{Les trois médianes du triangle $ABC$ sont concourantes en $G$, avec $\overrightarrow{AG}=\dfrac{2}{3}\overrightarrow{AA'}$.}
\end{popfigure}

\begin{aretenir}[L'essentiel sur le centre de gravité]
\begin{itemize}
\item Le centre de gravité $G$ du triangle $ABC$ est l'\textbf{isobarycentre} des trois sommets : $\overrightarrow{GA}+\overrightarrow{GB}+\overrightarrow{GC}=\vec{0}$.
\item $G$ est situé aux $\dfrac{2}{3}$ de chaque médiane à partir du sommet : $\overrightarrow{AG}=\dfrac{2}{3}\overrightarrow{AA'}$.
\item Pour quatre points, l'isobarycentre des sommets d'un parallélogramme est l'intersection des diagonales.
\end{itemize}
\end{aretenir}

\courssub{Lien avec la physique : centre de gravité d'une plaque et d'un système de masses}

En physique, le \cle{centre de gravité} d'un système matériel est exactement le barycentre des points où sont concentrées les masses, affectées de coefficients proportionnels à ces masses. C'est le point d'application de la force de pesanteur globale.

\begin{exemplebox}[Plaque rectangulaire homogène]
Une plaque métallique rectangulaire homogène $ABCD$ (épaisseur et densité constantes) peut être modélisée par quatre masses égales placées aux quatre sommets. Son centre de gravité est donc l'isobarycentre des quatre sommets : d'après la propriété démontrée plus haut, c'est \textbf{l'intersection des diagonales}. C'est exactement là que l'ouvrier doit placer le crochet du palan pour que la plaque reste horizontale pendant le levage. Si la plaque est percée ou renforcée en un point, on ajoute un point pondéré (masse négative pour un trou, masse supplémentaire pour un renfort) et le centre de gravité se déplace : on le recalcule par associativité.
\end{exemplebox}

\begin{exoresolu}[Centre de gravité d'un système de masses]
\textbf{Énoncé.} Sur une barre rigide, on place quatre masses : $2$ kg en $A$, $1$ kg en $B$, $1$ kg en $C$ et $4$ kg en $D$. Déterminer la position du centre de gravité $G$ du système sachant que $ABCD$ est un parallélogramme.
\tcblower
\textbf{Solution.} La somme des coefficients est $2+1+1+4=8\neq 0$ : le barycentre $G=\mathrm{bar}\{(A,2),(B,1),(C,1),(D,4)\}$ existe.

Regroupons $(B,1)$ et $(C,1)$ : leur barycentre est $J$, milieu de $[BC]$, affecté du coefficient $2$.
Alors $G_1=\mathrm{bar}\{(A,2),(J,2)\}$ est le \textbf{milieu de $[AJ]$}, affecté du coefficient $4$.
Finalement, $G=\mathrm{bar}\{(G_1,4),(D,4)\}$ : $G$ est le \textbf{milieu du segment} $[G_1D]$.

Concrètement : on place $J$ milieu de $[BC]$, puis $G_1$ milieu de $[AJ]$, puis $G$ milieu de $[G_1D]$. La masse de $4$ kg en $D$ « attire » fortement le centre de gravité vers $D$, ce qui est conforme à l'intuition physique : c'est vers $D$ qu'il faut déplacer le point de levage pour équilibrer la barre.
\end{exoresolu}

\begin{attention}[Coefficients négatifs : attention à l'interprétation]
Un coefficient négatif n'a pas de sens physique direct (pas de « masse négative »), mais il est très utile en géométrie et pour modéliser un \emph{enlèvement de matière} (un trou dans une plaque). Dans ce cas, le barycentre peut se trouver \textbf{à l'extérieur} du système de points : ce n'est pas une erreur de calcul. En revanche, si tous les coefficients sont positifs, le barycentre est nécessairement « entouré » par les points (à l'intérieur du polygone qu'ils forment).
\end{attention}

\begin{savaistu}[Le centre de gravité au chantier naval de Limbé]
Les ingénieurs des chantiers navals, comme ceux de Limbé au pied du mont Cameroun, utilisent le centre de gravité pour vérifier la \textbf{stabilité des embarcations}. Une pirogue ou un bateau est stable si son centre de gravité $G$ (barycentre de toutes les masses : coque, moteur, cargaison, passagers) est suffisamment \emph{bas} et bien \emph{centré}. Si la cargaison est mal répartie, $G$ se déplace latéralement et l'embarcation risque le chavirement. Avant la mise à l'eau, on calcule $G$ exactement comme un barycentre de points pondérés : chaque élément du bateau est une « masse » placée en un point, et l'associativité permet de regrouper les calculs par zones (proue, poupe, pont, cale). Les mathématiques de cette leçon sont donc un véritable outil de sécurité en mer !
\end{savaistu}

\begin{exercice}[Pour s'entraîner]
Soit $ABCD$ un parallélogramme de centre $O$.
\begin{enumerate}
\item Justifier que $O=\mathrm{bar}\{(A,1),(B,1),(C,1),(D,1)\}$.
\item On considère $G=\mathrm{bar}\{(A,1),(B,2),(C,1),(D,2)\}$. Montrer, par associativité, que $G$ est confondu avec $O$.
\item Soit $M$ un point quelconque. Montrer que $\overrightarrow{MA}+\overrightarrow{MB}+\overrightarrow{MC}+\overrightarrow{MD}=4\overrightarrow{MO}$.
\end{enumerate}
\end{exercice}

\begin{corrige}[Exercice]
\begin{enumerate}
\item $O$ est le milieu de $[AC]$ et de $[BD]$ (propriété du parallélogramme). Par associativité, $\mathrm{bar}\{(A,1),(B,1),(C,1),(D,1)\}=\mathrm{bar}\{(O,2),(O,2)\}=O$.
\item La somme des coefficients est $1+2+1+2=6\neq 0$. Regroupons : $\mathrm{bar}\{(A,1),(C,1)\}=O$ (coefficient $2$), et $\mathrm{bar}\{(B,2),(D,2)\}=O$ (coefficient $4$). Donc $G=\mathrm{bar}\{(O,2),(O,4)\}=O$. Le point $G$ est confondu avec le centre $O$ du parallélogramme.
\item Par la formule de réduction appliquée à l'isobarycentre $O$ des quatre sommets : pour tout point $M$,
\[
\overrightarrow{MA}+\overrightarrow{MB}+\overrightarrow{MC}+\overrightarrow{MD}=(1+1+1+1)\overrightarrow{MO}=4\overrightarrow{MO}.
\]
\end{enumerate}
\end{corrige}

En résumé, le passage à quatre points pondérés ne demande aucune notion nouvelle : c'est l'\textbf{associativité} qui, par des regroupements astucieux, transforme des démonstrations géométriques parfois délicates (alignements, concours de droites, théorème des médianes) en simples calculs de barycentres. Dans la prochaine section, nous utiliserons ces outils pour étudier les \textbf{lignes de niveau}, c'est-à-dire les ensembles de points vérifiant une relation portant sur les distances aux sommets d'un système pondéré.

\coursec{Lignes de niveau}

Dans la section précédente, nous avons vu que le barycentre $G$ de points pondérés permet de \emph{réduire} une somme de vecteurs : $\alpha\overrightarrow{MA}+\beta\overrightarrow{MB}=(\alpha+\beta)\overrightarrow{MG}$. Cette identité est un véritable outil de simplification. Nous allons maintenant l'exploiter pour répondre à une question typique du Probatoire : \emph{quel est l'ensemble des points $M$ du plan vérifiant une condition donnée ?} Ces ensembles s'appellent des \cle{lignes de niveau}.

\courssub{Notion de ligne de niveau}

\textbf{Intuition.} Sur une carte topographique, une courbe de niveau relie tous les points situés à la même altitude. En météorologie, les isobares relient les points de même pression. En mathématiques, l'idée est identique : on regroupe les points $M$ du plan pour lesquels une certaine quantité $f(M)$ (une distance, une somme de carrés de distances, un rapport de distances...) prend une valeur fixée $k$.

\begin{definition}[Ligne de niveau]
Soit $f$ une fonction qui à tout point $M$ du plan associe un nombre réel $f(M)$, et soit $k$ un réel fixé. La \cle{ligne de niveau} $k$ de $f$ est l'ensemble des points $M$ du plan tels que
\[ f(M)=k. \]
On la note souvent $\mathcal{L}_k$ ou $\Gamma_k$.
\end{definition}

\begin{exemplebox}[Premiers exemples]
\begin{itemize}
\item Si $f(M)=OM$ (distance à un point fixe $O$), la ligne de niveau $k=3$ est le cercle de centre $O$ et de rayon $3$.
\item Si $f(M)=MA-MB$, la ligne de niveau $k=0$ est l'ensemble des points équidistants de $A$ et $B$ : nous verrons que c'est la médiatrice de $[AB]$.
\end{itemize}
\end{exemplebox}

\begin{methode}[Déterminer une ligne de niveau en trois étapes]
\begin{enumerate}
\item \textbf{Réduire} l'expression à l'aide du barycentre $G$ des points pondérés concernés (relation de Chasles en introduisant $G$).
\item \textbf{Isoler} la quantité $MG$ ou $\overrightarrow{MG}$ dans l'égalité obtenue.
\item \textbf{Conclure} sur la nature géométrique de l'ensemble : droite (vecteur colinéaire à un vecteur fixe, ou $MG$ perpendiculaire à une direction), cercle ($MG=r$), point ($MG=0$) ou ensemble vide.
\end{enumerate}
\end{methode}

\courssub{Lignes de niveau vectorielles}

\begin{propriete}[Ligne de niveau de $M\mapsto\alpha\overrightarrow{MA}+\beta\overrightarrow{MB}$]
Soient $A$ et $B$ deux points distincts, $\alpha$ et $\beta$ deux réels avec $\alpha+\beta\neq 0$, et $G$ le barycentre de $(A,\alpha)$, $(B,\beta)$. L'ensemble des points $M$ tels que $\alpha\overrightarrow{MA}+\beta\overrightarrow{MB}$ soit colinéaire à un vecteur fixe non nul $\vec{u}$ est la \textbf{droite passant par $G$} et de direction $\vec{u}$.
\end{propriete}

\textbf{Justification.} Par réduction, $\alpha\overrightarrow{MA}+\beta\overrightarrow{MB}=(\alpha+\beta)\overrightarrow{MG}$. Dire que ce vecteur est colinéaire à $\vec{u}$ revient à dire que $\overrightarrow{MG}$ est colinéaire à $\vec{u}$, c'est-à-dire que $M$ appartient à la droite passant par $G$ dirigée par $\vec{u}$.

\begin{propriete}[Ligne de niveau de $M\mapsto\|\alpha\overrightarrow{MA}+\beta\overrightarrow{MB}\|$]
Avec les mêmes hypothèses ($\alpha+\beta\neq 0$), l'ensemble des points $M$ tels que $\|\alpha\overrightarrow{MA}+\beta\overrightarrow{MB}\|=r$ (avec $r>0$) est le \textbf{cercle de centre $G$} et de rayon $\dfrac{r}{|\alpha+\beta|}$.
\end{propriete}

\textbf{Justification.} $\|\alpha\overrightarrow{MA}+\beta\overrightarrow{MB}\|=|\alpha+\beta|\times MG$, donc l'équation équivaut à $MG=\dfrac{r}{|\alpha+\beta|}$ : c'est un cercle de centre $G$.

\courssub{Ligne de niveau de $M\mapsto\alpha MA^2+\beta MB^2$}

\textbf{Intuition.} La quantité $MA^2+MB^2$ mesure en quelque sorte « l'éloignement total » de $M$ par rapport à $A$ et $B$. Les points où cet éloignement est constant devraient former une courbe fermée autour du « centre » de $A$ et $B$ : nous allons démontrer qu'il s'agit d'un cercle centré au barycentre.

\begin{propriete}[Théorème — démonstration exigible]
Soient $A$ et $B$ deux points, $\alpha$ et $\beta$ deux réels tels que $\alpha+\beta\neq 0$, et $G$ le barycentre de $(A,\alpha)$, $(B,\beta)$. Pour tout point $M$ :
\[ \alpha MA^2+\beta MB^2=(\alpha+\beta)MG^2+\alpha GA^2+\beta GB^2. \]
L'ensemble des points $M$ tels que $\alpha MA^2+\beta MB^2=k$ est donc :
\begin{itemize}
\item le \textbf{cercle de centre $G$} et de rayon $\sqrt{\dfrac{k-(\alpha GA^2+\beta GB^2)}{\alpha+\beta}}$ si cette quantité sous la racine est strictement positive ;
\item le \textbf{point $G$ seul} si elle est nulle ;
\item l'\textbf{ensemble vide} si elle est négative.
\end{itemize}
\end{propriete}

\begin{experience}[Démonstration guidée du théorème]
On part de $MA^2=\overrightarrow{MA}^2$ et on introduit $G$ par la relation de Chasles.
\begin{enumerate}
\item Écris $\overrightarrow{MA}=\overrightarrow{MG}+\overrightarrow{GA}$, puis développe :
\[ MA^2=\overrightarrow{MA}^2=MG^2+2\,\overrightarrow{MG}\cdot\overrightarrow{GA}+GA^2. \]
\item Fais de même pour $MB^2=MG^2+2\,\overrightarrow{MG}\cdot\overrightarrow{GB}+GB^2$.
\item Multiplie la première égalité par $\alpha$, la seconde par $\beta$, et additionne :
\[ \alpha MA^2+\beta MB^2=(\alpha+\beta)MG^2+2\,\overrightarrow{MG}\cdot(\alpha\overrightarrow{GA}+\beta\overrightarrow{GB})+\alpha GA^2+\beta GB^2. \]
\item Comme $G$ est le barycentre de $(A,\alpha)$, $(B,\beta)$, on a $\alpha\overrightarrow{GA}+\beta\overrightarrow{GB}=\vec{0}$. Le terme du milieu disparaît :
\[ \alpha MA^2+\beta MB^2=(\alpha+\beta)MG^2+\alpha GA^2+\beta GB^2. \]
\item L'équation $\alpha MA^2+\beta MB^2=k$ devient $MG^2=\dfrac{k-(\alpha GA^2+\beta GB^2)}{\alpha+\beta}$. Selon le signe du second membre, on obtient un cercle, un point ou l'ensemble vide.
\end{enumerate}
\end{experience}

\begin{exoresolu}[Ligne de niveau de $M\mapsto MA^2+MB^2$]
\textbf{Énoncé.} Soit $[AB]$ un segment de longueur $AB=6$. Déterminer l'ensemble $\Gamma$ des points $M$ du plan tels que $MA^2+MB^2=50$.
\tcblower
\textbf{Solution détaillée.} On applique la méthode en trois étapes.

\textbf{Étape 1 : réduire avec $G$.} Ici $\alpha=\beta=1$, donc $\alpha+\beta=2\neq 0$ et $G$ est le milieu de $[AB]$. On a $GA=GB=3$, donc $GA^2=GB^2=9$.

\textbf{Étape 2 : isoler $MG$.} Par le théorème :
\[ MA^2+MB^2=2\,MG^2+GA^2+GB^2=2\,MG^2+18. \]
L'équation $MA^2+MB^2=50$ équivaut à $2\,MG^2=32$, soit $MG^2=16$, c'est-à-dire $MG=4$.

\textbf{Étape 3 : conclure.} $\Gamma$ est le cercle de centre $G$ (milieu de $[AB]$) et de rayon $4$.

\textbf{Vérification.} Pour $M=A$ : $MA^2+MB^2=0+36=36<50$ ; le point $A$ est donc à l'intérieur du cercle, ce qui est cohérent car $GA=3<4$.
\end{exoresolu}

\begin{popfigure}
\begin{tikzpicture}[scale=0.62]
\draw[popline, dashed] (-4.2,0) -- (4.2,0);
\coordinate (A) at (-3,0);
\coordinate (B) at (3,0);
\coordinate (G) at (0,0);
\draw[popblue, thick] (G) circle (1.2);
\draw[poporange, thick] (G) circle (2.4);
\draw[poppurple, thick] (G) circle (3.6);
\fill[popink] (A) circle (2.2pt) node[below left] {$A$};
\fill[popink] (B) circle (2.2pt) node[below right] {$B$};
\fill[popdark] (G) circle (2.2pt) node[below] {$G$};
\node[popblue] at (0,1.55) {\small $k=20{,}88$};
\node[poporange] at (0,2.75) {\small $k=29{,}52$};
\node[poppurple] at (0,3.95) {\small $k=43{,}92$};
\end{tikzpicture}
\legende{Lignes de niveau de $M\mapsto MA^2+MB^2=k$ pour $AB=6$ : ce sont des cercles concentriques centrés en $G$, milieu de $[AB]$. Plus $k$ est grand, plus le rayon est grand.}
\end{popfigure}

\begin{attention}[Discuter l'existence selon le signe du second membre]
L'erreur classique consiste à écrire directement « c'est un cercle de rayon $\sqrt{\cdots}$ » sans vérifier que la quantité sous la racine est positive. Or $MG^2$ est toujours positif ou nul :
\begin{itemize}
\item si $\dfrac{k-(\alpha GA^2+\beta GB^2)}{\alpha+\beta}<0$, l'ensemble est \textbf{vide} (aucun point ne convient) ;
\item si elle est nulle, l'ensemble est réduit au \textbf{point $G$} (ce n'est pas un cercle !).
\end{itemize}
Par exemple, avec $AB=6$, l'ensemble des points tels que $MA^2+MB^2=10$ est vide, car on obtiendrait $2MG^2=-8$.
\end{attention}

\courssub{Ligne de niveau de $M\mapsto\dfrac{MA}{MB}$}

\begin{propriete}[La médiatrice comme ligne de niveau — démonstration exigible]
Soient $A$ et $B$ deux points distincts. L'ensemble des points $M$ tels que $MA=MB$ (c'est-à-dire $\dfrac{MA}{MB}=1$) est la \textbf{médiatrice du segment $[AB]$}.
\end{propriete}

\textbf{Démonstration.} On raisonne par équivalences, en introduisant le milieu $I$ de $[AB]$ :
\begin{align*}
MA=MB &\iff MA^2=MB^2 \iff MA^2-MB^2=0 \\
&\iff (\overrightarrow{MA}-\overrightarrow{MB})\cdot(\overrightarrow{MA}+\overrightarrow{MB})=0 \\
&\iff \overrightarrow{BA}\cdot 2\,\overrightarrow{MI}=0 \iff \overrightarrow{MI}\cdot\overrightarrow{AB}=0.
\end{align*}
La dernière condition signifie que $(MI)\perp(AB)$ : $M$ appartient à la droite passant par $I$ et perpendiculaire à $(AB)$, c'est-à-dire la médiatrice de $[AB]$.

\begin{propriete}[Cercle d'Apollonius — admis]
Soient $A$ et $B$ deux points distincts et $k$ un réel strictement positif, $k\neq 1$. L'ensemble des points $M$ tels que $\dfrac{MA}{MB}=k$ est un \textbf{cercle}, appelé \cle{cercle d'Apollonius}, dont le diamètre a pour extrémités les barycentres de $(A,1)$, $(B,-k)$ et de $(A,1)$, $(B,k)$ situés sur la droite $(AB)$.
\end{propriete}

\begin{popfigure}
\begin{tikzpicture}[scale=0.85]
\coordinate (A) at (0,0);
\coordinate (B) at (3,0);
\coordinate (I) at (1.5,0);
% médiatrice
\draw[popblue, thick] (1.5,-2.2) -- (1.5,2.2);
\node[popblue, rotate=90, anchor=south] at (1.5,0.9) {\small médiatrice : $MA=MB$};
% cercle d'Apollonius pour MA/MB = 2 : points de (AB) : x=6 et x=2 ; centre 4, rayon 2
\draw[poporange, thick] (4,0) circle (2);
\node[poporange] at (4,2.3) {\small $\dfrac{MA}{MB}=2$};
\fill[popink] (A) circle (2pt) node[below] {$A$};
\fill[popink] (B) circle (2pt) node[below] {$B$};
\fill[popblue] (I) circle (1.8pt) node[below] {\small $I$};
\fill[poporange] (2,0) circle (1.8pt);
\fill[poporange] (6,0) circle (1.8pt);
\draw[popline, dashed] (-0.6,0) -- (6.6,0);
\end{tikzpicture}
\legende{Ligne de niveau $k=1$ de $M\mapsto \dfrac{MA}{MB}$ : la médiatrice de $[AB]$ (droite bleue). Ligne de niveau $k=2$ : le cercle d'Apollonius (orange), de diamètre les points de $(AB)$ vérifiant $MA=2MB$.}
\end{popfigure}

\courssub{Ligne de niveau de $M\mapsto MA^2-MB^2$}

\textbf{Intuition.} Contrairement au cas précédent, les coefficients $\alpha=1$ et $\beta=-1$ ont une somme nulle : on ne peut pas introduire de barycentre. La démonstration de la médiatrice ci-dessus montre la voie : utiliser l'identité remarquable $a^2-b^2=(a-b)(a+b)$ avec des vecteurs.

\begin{propriete}[Ligne de niveau de $M\mapsto MA^2-MB^2$]
Soient $A$ et $B$ deux points distincts et $k$ un réel. L'ensemble des points $M$ tels que $MA^2-MB^2=k$ est une \textbf{droite perpendiculaire à $(AB)$}.
\end{propriete}

\begin{experience}[Démonstration guidée]
Soit $I$ le milieu de $[AB]$.
\begin{enumerate}
\item Factorise : $MA^2-MB^2=(\overrightarrow{MA}-\overrightarrow{MB})\cdot(\overrightarrow{MA}+\overrightarrow{MB})$.
\item Simplifie chaque facteur : $\overrightarrow{MA}-\overrightarrow{MB}=\overrightarrow{BA}$ et $\overrightarrow{MA}+\overrightarrow{MB}=2\overrightarrow{MI}$ (car $I$ est le milieu de $[AB]$).
\item Conclus : $MA^2-MB^2=2\,\overrightarrow{MI}\cdot\overrightarrow{BA}$.
\item L'équation $MA^2-MB^2=k$ devient $\overrightarrow{MI}\cdot\overrightarrow{BA}=\dfrac{k}{2}$ : le produit scalaire de $\overrightarrow{MI}$ avec le vecteur fixe $\overrightarrow{BA}$ est constant. Géométriquement, le projeté orthogonal de $M$ sur $(AB)$ est un point fixe $H$ ; l'ensemble cherché est la droite perpendiculaire à $(AB)$ passant par $H$.
\end{enumerate}
\end{experience}

\begin{exemplebox}[Exemple chiffré]
Avec $AB=6$ et $k=18$ : on cherche $H$ sur $(AB)$ tel que $\overrightarrow{HI}\cdot\overrightarrow{BA}=9$. En prenant un repère d'origine $I$ sur $(AB)$ avec $A$ d'abscisse $-3$ et $B$ d'abscisse $3$, le vecteur $\overrightarrow{BA}$ a pour coordonnée $-6$. L'équation $\overrightarrow{HI}\cdot\overrightarrow{BA}=9$ s'écrit $(-x_H)\times(-6)=9$, soit $6x_H=9$, d'où $x_H=1,5$. L'ensemble des points $M$ tels que $MA^2-MB^2=18$ est la droite perpendiculaire à $(AB)$ passant par $H$ (abscisse $1,5$, du côté de $B$ car $MA>MB$).
\end{exemplebox}

\courssub{Bilan et application à la vie courante}

\begin{aretenir}[Tableau récapitulatif des lignes de niveau]
\begin{center}
\begin{tabularx}{\linewidth}{|l|X|}
\hline
\rowcolor{popblueL} \textbf{Condition sur $M$} & \textbf{Ensemble des points $M$} \\
\hline
$\alpha\overrightarrow{MA}+\beta\overrightarrow{MB}$ colinéaire à $\vec{u}$ ($\alpha+\beta\neq 0$) & Droite passant par $G$ \\
\hline
$\|\alpha\overrightarrow{MA}+\beta\overrightarrow{MB}\|=r$ ($\alpha+\beta\neq 0$) & Cercle de centre $G$ \\
\hline
$\alpha MA^2+\beta MB^2=k$ ($\alpha+\beta\neq 0$) & Cercle de centre $G$, point $G$ ou ensemble vide selon le signe du second membre \\
\hline
$MA=MB$ & Médiatrice de $[AB]$ \\
\hline
$\dfrac{MA}{MB}=k$, $k>0$, $k\neq 1$ & Cercle d'Apollonius (admis) \\
\hline
$MA^2-MB^2=k$ & Droite perpendiculaire à $(AB)$ \\
\hline
\end{tabularx}
\end{center}
\end{aretenir}

\begin{savaistu}[Les lignes de niveau autour de nous]
Les lignes de niveau ne sont pas qu'un exercice de style : en topographie, les cartes de l'IGN ou celles utilisées sur les chantiers de construction au Cameroun (routes, barrages) représentent le relief par des courbes de niveau d'altitude. En électricité, les lignes équipotentielles autour d'un conducteur sont les lignes de niveau du potentiel électrique. Le cercle d'Apollonius, lui, porte le nom du mathématicien grec Apollonius de Perga (III\textsuperscript{e} siècle av. J.-C.), surnommé « le Grand Géomètre ».
\end{savaistu}

\begin{exercice}
Soit $[AB]$ un segment tel que $AB=8$ et $G$ le barycentre de $(A,1)$, $(B,3)$.
\begin{enumerate}
\item Déterminer l'ensemble des points $M$ tels que $\|\overrightarrow{MA}+3\overrightarrow{MB}\|=12$.
\item Déterminer l'ensemble des points $M$ tels que $MA^2+3MB^2=76$.
\end{enumerate}
\end{exercice}

\begin{corrige}[Exercice]
\begin{enumerate}
\item $\alpha+\beta=4\neq 0$ et $\overrightarrow{MA}+3\overrightarrow{MB}=4\overrightarrow{MG}$. L'équation équivaut à $4MG=12$, soit $MG=3$ : c'est le cercle de centre $G$ et de rayon $3$.
\item On calcule $GA=6$ et $GB=2$ (car $G$ partage $[AB]$ avec $AG=\dfrac{3}{4}AB$). La réduction donne $MA^2+3MB^2=4MG^2+GA^2+3GB^2=4MG^2+36+12=4MG^2+48$. L'équation $MA^2+3MB^2=76$ équivaut à $4MG^2=28$, soit $MG^2=7>0$ : c'est le cercle de centre $G$ et de rayon $\sqrt{7}$.
\end{enumerate}
\end{corrige}

\coursec{Applications à la vie courante et synthèse}

Après avoir étudié les lignes de niveau, qui permettent de visualiser les points d'un plan ayant une même valeur d'une fonction affine, nous allons maintenant voir comment le barycentre devient un outil de modélisation puissant pour des situations concrètes : équilibre d'un levier, centre de gravité de plaques, localisation équitable d'un service public. Chaque situation suivra la même démarche : identifier les points et les masses (ou coefficients), écrire la relation de barycentre, calculer ou construire le point recherché, puis interpréter le résultat dans le contexte technique.

\courssub{Modélisation de l'équilibre d'un levier : la règle du levier}

\begin{definition}[Équilibre d'un levier et barycentre]
Soit une barre homogène de longueur $L$ appuyée en un point $G$. Si une masse $m_1$ est suspendue en $A$ et une masse $m_2$ en $B$, l'équilibre statique (horizontal) est réalisé si et seulement si les moments des forces de gravité par rapport à $G$ s'annulent :
\[
m_1 \cdot GA = m_2 \cdot GB \qquad (GA, GB \text{ distances algébriques}).
\]
Cette relation est exactement la définition du barycentre $G$ du système de deux points pondérés $\bigl\{(A,m_1),\,(B,m_2)\bigr\}$.
\end{definition}

\begin{propriete}[Barycentre de deux points pondérés]
Pour deux points $A,B$ et des coefficients réels $a,b$ tels que $a+b\neq0$, le barycentre $G$ est l'unique point vérifiant
\[
a\,\overrightarrow{GA}+b\,\overrightarrow{GB}=\vec{0}
\quad\Longleftrightarrow\quad
\overrightarrow{OG}=\frac{a\,\overrightarrow{OA}+b\,\overrightarrow{OB}}{a+b}
\]
pour tout origine $O$. Si $a,b>0$, $G$ appartient au segment $[AB]$ et divise ce segment dans le rapport inverse des masses : $\dfrac{GA}{GB}=\dfrac{b}{a}$.
\end{propriete}

\begin{methode}[Modéliser un problème d'équilibre par un barycentre]
\begin{enumerate}
\item Identifier les points d'application des forces (généralement les extrémités d'une barre, les centres de gravité de pièces).
\item Associer à chaque point sa masse (ou son poids, proportionnel à la masse) ; ce sont les coefficients du barycentre.
\item Écrire la relation vectorielle $m_1\overrightarrow{GA}+m_2\overrightarrow{GB}=\vec{0}$ ou utiliser la formule des coordonnées.
\item Résoudre pour obtenir la position de $G$ (distance par rapport à un repère).
\item Vérifier que $G$ se trouve bien entre les points si les masses sont positives (équilibre stable).
\end{enumerate}
\end{methode}

\begin{exoresolu}[Barre de \SI{2}{\meter} avec moteur de \SI{6}{\kilo\gram} en A et charge de \SI{3}{\kilo\gram} en B]
\textbf{Énoncé.} Une barre homogène de longueur \SI{2}{\meter} porte un moteur de masse \SI{6}{\kilo\gram} à son extrémité $A$ et une charge de \SI{3}{\kilo\gram} à l'autre extrémité $B$. On cherche la position du point d'appui $G$ pour que la barre soit en équilibre horizontal. La masse de la barre est négligée.

\textbf{Solution.}
\begin{enumerate}
\item Points : $A$ (moteur), $B$ (charge). Masses : $m_A=6$, $m_B=3$ (en kg). Somme $m_A+m_B=9\neq0$, le barycentre $G$ existe et est unique.
\item Relation vectorielle : $6\,\overrightarrow{GA}+3\,\overrightarrow{GB}=\vec{0}$.
  En choisissant $A$ comme origine, $\overrightarrow{AB}$ est le vecteur orienté de $A$ vers $B$ de norme \SI{2}{\meter}.
  On a $\overrightarrow{AG}= \frac{m_B}{m_A+m_B}\,\overrightarrow{AB}= \frac{3}{9}\times\SI{2}{\meter}= \frac{2}{3}\,\text{m} \approx \SI{0,667}{\meter}$.
\item Donc $G$ se situe à \SI{0,667}{\meter} de $A$ (soit \SI{66,7}{\centi\meter}) et à \SI{1,333}{\meter} de $B$.
\item Vérification : $6\times0,667 \approx 4,0 = 3\times1,333$. L'équilibre est assuré.
\end{enumerate}
\end{exoresolu}

\begin{popfigure}
\begin{tikzpicture}[scale=1.2,>=Stealth]
% Barre
\draw[thick, popdark] (0,0) -- (5,0) node[midway,above=2pt] {Barre \SI{2}{\meter}};
% Points A, B
\fill[popblue] (0,0) circle (2pt) node[below] {$A$ (\SI{6}{\kilo\gram})};
\fill[poporange] (5,0) circle (2pt) node[below] {$B$ (\SI{3}{\kilo\gram})};
% Point G (calculé : 2/3 de 5 = 3.333? Non, échelle 5 unités = 2m -> 1m = 2.5 unités. 0.667m = 1.667 unités)
\coordinate (G) at (1.667,0);
\fill[popgreen] (G) circle (2pt) node[above] {$G$};
% Flèches distances
\draw[<->, popmuted] (0,-0.5) -- (1.667,-0.5) node[midway,below] {\SI{0,667}{\meter}};
\draw[<->, popmuted] (1.667,-0.5) -- (5,-0.5) node[midway,below] {\SI{1,333}{\meter}};
% Forces (poids)
\draw[->, thick, poppink] (0,0.3) -- (0,1) node[above] {$6g$};
\draw[->, thick, poppink] (5,0.3) -- (5,1) node[above] {$3g$};
\draw[->, thick, popgreen] (G) -- (G|-0,1.3) node[above] {Réaction $R=9g$};
\end{tikzpicture}
\legende{Schéma du levier en équilibre : le point d'appui $G$ est le barycentre des deux masses.}
\end{popfigure}

\begin{attention}[Masses positives vs coefficients signés]
En physique (mécanique, statique), les masses sont toujours \textbf{positives}. Le barycentre se trouve alors obligatoirement à l'intérieur du segment reliant les points (équilibre stable). En mathématiques pures, les coefficients $a,b$ peuvent être négatifs ; le barycentre existe tant que $a+b\neq0$, mais il peut se situer à l'extérieur du segment. Ne confondez pas les deux contextes : pour un problème d'équilibre réel, imposez $m_i>0$.
\end{attention}

\courssub{Centre de gravité d'une plaque triangulaire homogène}

\begin{propriete}[Centre de gravité d'un triangle homogène]
Le centre de gravité $G$ d'une plaque triangulaire homogène (masse surfacique constante) est l'isobarycentre de ses trois sommets $A,B,C$ :
\[
G = \operatorname{bar}\{(A,1),\,(B,1),\,(C,1)\}
\quad\Longleftrightarrow\quad
\overrightarrow{OG}=\frac{\overrightarrow{OA}+\overrightarrow{OB}+\overrightarrow{OC}}{3}.
\]
$G$ est aussi l'intersection des trois médianes ; il divise chaque médiane dans le rapport $2:1$ à partir du sommet.
\end{propriete}

\begin{popfigure}
\begin{tikzpicture}[scale=1.3,>=Stealth]
% Triangle
\coordinate (A) at (0,0);
\coordinate (B) at (4,0);
\coordinate (C) at (1,3);
\draw[thick, popdark] (A) -- (B) -- (C) -- cycle;
% Sommets
\fill[popblue] (A) circle (2pt) node[below left] {$A$};
\fill[popblue] (B) circle (2pt) node[below right] {$B$};
\fill[popblue] (C) circle (2pt) node[above] {$C$};
% Médianes
\coordinate (Mbc) at ($(B)!0.5!(C)$);
\coordinate (Mca) at ($(C)!0.5!(A)$);
\coordinate (Mab) at ($(A)!0.5!(B)$);
\draw[dashed, popmuted] (A) -- (Mbc);
\draw[dashed, popmuted] (B) -- (Mca);
\draw[dashed, popmuted] (C) -- (Mab);
% Centre de gravité (isobarycentre via calc)
\coordinate (G) at ($ 1/3*(A) + 1/3*(B) + 1/3*(C) $);
\fill[popgreen] (G) circle (2pt) node[right] {$G$};
% Étiquettes masses
\node[popblue] at ($(A)+(-0.3,-0.3)$) {$1$};
\node[popblue] at ($(B)+(0.3,-0.3)$) {$1$};
\node[popblue] at ($(C)+(0,0.3)$) {$1$};
\end{tikzpicture}
\legende{Centre de gravité $G$ d'un triangle homogène : isobarycentre des sommets.}
\end{popfigure}

\courssub{Centre de gravité d'une plaque composée : associativité des barycentres}

Lorsqu'une plaque plane est découpée en plusieurs parties dont on connaît les centres de gravité et les masses, le centre de gravité global s'obtient en appliquant l'\textbf{associativité des barycentres} : on remplace chaque partie par son centre de gravité affecté de la masse de cette partie.

\begin{propriete}[Associativité des barycentres]
Soit un ensemble de points $A_1,\dots,A_n$ avec coefficients $\alpha_1,\dots,\alpha_n$ ($\sum\alpha_i\neq0$). Si l'on regroupe les points en sous-ensembles $P_1,\dots,P_k$ et que l'on définit pour chaque $j$ le barycentre partiel $G_j$ de $P_j$ avec coefficient $\beta_j=\sum_{i\in P_j}\alpha_i$, alors le barycentre global $G$ de l'ensemble initial coïncide avec le barycentre des points $G_j$ pondérés par $\beta_j$ :
\[
G = \operatorname{bar}\{(A_i,\alpha_i)\}_{i=1}^n
= \operatorname{bar}\{(G_j,\beta_j)\}_{j=1}^k .
\]
\end{propriete}

\begin{popfigure}
\begin{tikzpicture}[scale=1.1,>=Stealth]
% Partie 1 (gauche)
\draw[thick, popblue, fill=popblue!10] (0,0) rectangle (3,2);
\coordinate (G1) at (1.5,1);
\fill[popblue] (G1) circle (2pt) node[above left] {$G_1$};
\node[popblue] at (1.5,2.3) {$m_1=\SI{4}{\kilo\gram}$};
% Partie 2 (droite)
\draw[thick, poporange, fill=poporange!10] (3,0) rectangle (6,2);
\coordinate (G2) at (4.5,1);
\fill[poporange] (G2) circle (2pt) node[above right] {$G_2$};
\node[poporange] at (4.5,2.3) {$m_2=\SI{2}{\kilo\gram}$};
% Centre global (calculé par associativité via calc)
\coordinate (G) at (2.5,1);
\fill[popgreen] (G) circle (3pt) node[below] {$G$};
\node[popgreen] at (2.5,-0.5) {$m=\SI{6}{\kilo\gram}$};
% Flèches
\draw[->, thick, popmuted] (G1) -- (G);
\draw[->, thick, popmuted] (G2) -- (G);
\end{tikzpicture}
\legende{Plaque composée de deux rectangles : les centres de gravité partiels $G_1,G_2$ (masses \SI{4}{\kilo\gram} et \SI{2}{\kilo\gram}) donnent le centre global $G$ par associativité.}
\end{popfigure}

\courssub{Problème de partage : point de service équitable entre plusieurs villages}

Imaginons trois localités $A,B,C$ avec des populations $p_A,p_B,p_C$. On souhaite implanter un centre de santé $G$ de façon à minimiser la distance moyenne pondérée par la population. Le point optimal est le barycentre $G=\operatorname{bar}\{(A,p_A),(B,p_B),(C,p_C)\}$.

\begin{exemplebox}[Localisation d'un centre de santé]
Trois villages sont situés aux coordonnées $A(0,0)$, $B(10,0)$, $C(4,8)$ (distances en km). Leurs populations sont $p_A=2000$, $p_B=3000$, $p_C=1000$. Calculons les coordonnées du centre de santé $G$.
\begin{align*}
x_G &= \frac{2000\cdot0 + 3000\cdot10 + 1000\cdot4}{2000+3000+1000}
= \frac{0+30000+4000}{6000}
= \frac{34000}{6000}
= \frac{17}{3} \approx 5,67\,\text{km},\\
y_G &= \frac{2000\cdot0 + 3000\cdot0 + 1000\cdot8}{6000}
= \frac{8000}{6000}
= \frac{4}{3} \approx 1,33\,\text{km}.
\end{align*}
Le centre de santé sera implanté au point $G(5,67;\,1,33)$, plus proche du village le plus peuplé $B$, ce qui est logique.
\end{exemplebox}

\begin{popfigure}
\begin{tikzpicture}[scale=0.6,>=Stealth]
% Axes
\draw[->, popmuted] (-1,0) -- (12,0) node[right] {$x$ (km)};
\draw[->, popmuted] (0,-1) -- (0,10) node[above] {$y$ (km)};
% Villages
\coordinate (A) at (0,0);
\coordinate (B) at (10,0);
\coordinate (C) at (4,8);
\fill[popblue] (A) circle (3pt) node[below left, align=center]{$A(0,0)$\\$2000$ hab.};
\fill[poporange] (B) circle (3pt) node[below right, align=center]{$B(10,0)$\\$3000$ hab.};
\fill[poppurple] (C) circle (3pt) node[above, align=center]{$C(4,8)$\\$1000$ hab.};
% Centre G (calculé via calc)
\coordinate (G) at ($ 1/3*(A) + 1/2*(B) + 1/6*(C) $);
\fill[popgreen] (G) circle (4pt) node[right] {$G(5,67;1,33)$};
% Lignes de liaison
\draw[dashed, popmuted] (A) -- (G);
\draw[dashed, popmuted] (B) -- (G);
\draw[dashed, popmuted] (C) -- (G);
% Grille légère
\foreach \x in {0,2,...,10} \draw[popline] (\x,-0.1) -- (\x,0.1);
\foreach \y in {0,2,...,8} \draw[popline] (-0.1,\y) -- (0.1,\y);
\end{tikzpicture}
\legende{Carte simplifiée : trois villages et le point de service $G$ pondéré par les populations.}
\end{popfigure}

\courssub{Rédaction d'une démarche complète}

Pour tout problème technique mobilisant le barycentre, la rédaction doit suivre une structure rigoureuse, conforme aux attendus du Probatoire.

\begin{methode}[Démarche complète de résolution]
\begin{enumerate}
\item \textbf{Modéliser} : décrire la situation physique (levier, plaque, réseau de villages), choisir un repère adapté, nommer les points et les grandeurs (masses, populations, coordonnées).
\item \textbf{Choisir les coefficients} : identifier les masses ou poids (toujours positifs en physique) ; vérifier que la somme n'est pas nulle.
\item \textbf{Calculer ou construire} :
      \begin{itemize}
      \item \textit{Calcul vectoriel} : utiliser $\overrightarrow{OG}=\frac{\sum m_i\overrightarrow{OA_i}}{\sum m_i}$.
      \item \textit{Construction géométrique} : pour deux points, diviser le segment dans le rapport inverse des masses ; pour trois points, construire les médianes (isobarycentre) ou utiliser la règle du parallélogramme itérative.
      \end{itemize}
\item \textbf{Conclure et interpréter} : donner la position numérique (distance, coordonnées) avec unité, vérifier la cohérence (point entre les masses si positives), et formuler une phrase de réponse dans le langage du métier (ex. « Le point d'appui doit être placé à \SI{66,7}{\centi\meter} du moteur »).
\end{enumerate}
\end{methode}

\courssub{Complément pour la Terminale (non exigible au Probatoire)}

\begin{savaistu}[Coordonnées barycentriques et barycentre dans l'espace]
En Terminale, la notion de barycentre s'étend à l'espace à trois dimensions. Un point $M$ de l'espace peut être repéré par ses \textbf{coordonnées barycentriques} $(\alpha,\beta,\gamma,\delta)$ par rapport à un tétraèdre de référence $A,B,C,D$ (somme non nulle). Le barycentre de $n$ points pondérés s'écrit de la même façon vectorielle, et les propriétés d'associativité, de relation de Chasles, etc., restent valides. Cette généralisation est la base des coordonnées homogènes en infographie et en robotique.
\end{savaistu}

\begin{savaistu}[Le centre de gravité et la stabilité des camions sur la route Yaoundé–Bafoussam]
Sur les pentes sinueuses de la route Yaoundé–Bafoussam, un camion surchargé en hauteur voit son centre de gravité s'élever. Si le centre de gravité dépasse la base de sustentation (la projection au sol des roues), le moment de renversement créé par la force centrifuge dans un virage n'est plus compensé par le poids : le camion bascule. Les ingénieurs calculent la position du centre de gravité global (barycentre des masses du châssis, de la cargaison, de la cabine) pour définir la charge maximale admissible et la hauteur de chargement limite. C'est une application directe de l'associativité des barycentres vue plus haut.
\end{savaistu}

\newpage

\coursec{Activité d'intégration}

\coursec{Leçon 09 : Barycentres — Activité d'intégration : Implantation d'une borne-fontaine à Yaoundé}

\courssub{Objectifs de l'activité}
\begin{itemize}
    \item Mobiliser la notion de barycentre de deux, trois et quatre points pondérés pour résoudre un problème concret de localisation.
    \item Appliquer la propriété d'associativité du barycentre pour construire géométriquement le point cherché.
    \item Utiliser la relation de Leibniz (liaison entre barycentre et somme des carrés des distances pondérées) pour déterminer une ligne de niveau.
    \item Rédiger une argumentation technique structurée et justifiée, adaptée à un destinataire non mathématicien (la mairie).
\end{itemize}

\courssub{Situation}
La mairie de Yaoundé souhaite installer une nouvelle borne-fontaine autonome (alimentation solaire, filtration intégrée) pour desservir les trois quartiers densément peuplés entourant le marché de Mokolo : le quartier \textbf{Mokolo 1} (quartier A), le quartier \textbf{Mokolo 2} (quartier B) et le quartier \textbf{Essos} (quartier C).

Une étude démographique récente fournit les populations effectives :
\begin{center}
\begin{tabularx}{\linewidth}{|l|X|X|X|}\hline
\textbf{Quartier} & \textbf{Notation} & \textbf{Population (habitants)} & \textbf{Poids associé} \\ \hline
Mokolo 1 & $A$ & 2000 & $\alpha = 2$ \\
Mokolo 2 & $B$ & 3000 & $\beta = 3$ \\
Essos & $C$ & 1000 & $\gamma = 1$ \\ \hline
\end{tabularx}
\end{center}
Les poids sont choisis proportionnels aux populations (divisés par 1000 pour simplifier) : $\alpha=2$, $\beta=3$, $\gamma=1$. La somme des poids est $\alpha+\beta+\gamma = 6 \neq 0$, le barycentre $G$ existe donc et est unique.

Sur le plan cadastral quadrillé (repère orthonormé, unité : 100 m), les coordonnées des centres de gravité des quartiers sont :
\[
A(1\,;\,3),\qquad B(7\,;\,1),\qquad C(3\,;\,6).
\]

La mairie définit la « fatigue moyenne » d'un point $M$ comme la moyenne quadratique pondérée des distances aux trois quartiers :
\[
F(M) = \frac{\alpha\,MA^2 + \beta\,MB^2 + \gamma\,MC^2}{\alpha+\beta+\gamma}.
\]
Le point optimal $G$ minimisant cette fatigue est le barycentre du système $\{(A,\alpha),\,(B,\beta),\,(C,\gamma)\}$.

Cependant, le terrain au point $G$ peut s'avérer impraticable (rocher, canalisation, propriété privée). La mairie demande une \emph{zone de repli} : l'ensemble des points $M$ pour lesquels la fatigue $F(M)$ n'excède pas de plus de 20\% la fatigue minimale $F(G)$. Cette zone correspond à une ligne de niveau (un cercle) qu'il faut déterminer.

\courssub{Consignes}
\begin{enumerate}
    \item \textbf{Repérage et construction.} Placer les points $A$, $B$, $C$ dans le repère. Construire le barycentre $G$ par la méthode des barycentres partiels (associativité) : d'abord $K$ barycentre de $(A,2)$ et $(B,3)$, puis $G$ barycentre de $(K,5)$ et $(C,1)$. Tracer le triangle $ABC$ et situer $G$.
    \item \textbf{Calcul des coordonnées.} Calculer les coordonnées exactes de $K$ puis de $G$ à l'aide des formules de barycentre. Vérifier que $G$ appartient à l'intérieur du triangle $ABC$ et justifier cette propriété géométrique.
    \item \textbf{Fatigue minimale.} Calculer la valeur exacte de la fatigue minimale $F(G)$ en utilisant la relation de Leibniz (ou en calculant directement les distances $GA$, $GB$, $GC$).
    \item \textbf{Zone de repli (ligne de niveau).} Déterminer l'équation cartésienne de l'ensemble des points $M(x,y)$ tels que $F(M) = 1,2 \times F(G)$. Préciser la nature de ce lieu, son centre et son rayon (valeur approchée au mètre près). Tracer ce cercle sur la figure.
    \item \textbf{Note technique.} Rédiger une note d'une demi-page adressée au Maire, présentant : le choix du point $G$ (coordonnées, justification de l'optimalité), la zone de repli (description, rayon), et une recommandation finale. Le vocabulaire doit être accessible (pas de jargon mathématique inutile), la démarche logique et justifiée.
\end{enumerate}

\courssub{Ressources à mobiliser}
\begin{propriete}[Barycentre de deux points — Coordonnées]
Si $A(x_A,y_A)$, $B(x_B,y_B)$ et poids $\alpha,\beta$ ($\alpha+\beta\neq0$), le barycentre $K$ a pour coordonnées :
\[
x_K = \frac{\alpha x_A + \beta x_B}{\alpha+\beta},\qquad
y_K = \frac{\alpha y_A + \beta y_B}{\alpha+\beta}.
\]
Vecteur : $\overrightarrow{OK} = \frac{\alpha\overrightarrow{OA}+\beta\overrightarrow{OB}}{\alpha+\beta}$.
\end{propriete}

\begin{propriete}[Associativité du barycentre]
Le barycentre de $\{(A_i,\alpha_i)\}_{i=1..n}$ peut être obtenu en remplaçant un sous-ensemble de points par leur barycentre partiel (somme des poids conservée).
\end{propriete}

\begin{propriete}[Barycentre de $n$ points pondérés — Formulaire général]
Pour $n$ points $A_i(x_i,y_i)$ de poids $\alpha_i$ ($\sum\alpha_i \neq 0$), le barycentre $G$ a pour coordonnées :
\[
x_G = \frac{\sum_{i=1}^n \alpha_i x_i}{\sum_{i=1}^n \alpha_i},\qquad
y_G = \frac{\sum_{i=1}^n \alpha_i y_i}{\sum_{i=1}^n \alpha_i}.
\]
Cette formule s'applique directement pour 3 points (Première) et s'étend à 4 points (approfondissement).
\end{propriete}

\begin{propriete}[Relation de Leibniz (exigible au Probatoire)]
Pour tout point $M$ et barycentre $G$ de $\{(A_i,\alpha_i)\}$ avec $\sum\alpha_i \neq 0$ :
\[
\sum \alpha_i MA_i^2 = \sum \alpha_i GA_i^2 + \left(\sum\alpha_i\right) MG^2.
\]
Conséquence : $F(M) = F(G) + MG^2$. Les lignes de niveau de $F$ sont des cercles de centre $G$.
\end{propriete}

\begin{propriete}[Position du barycentre à coefficients positifs]
Si tous les $\alpha_i > 0$, le barycentre $G$ se trouve à l'intérieur de l'enveloppe convexe des points $A_i$ (ici, à l'intérieur du triangle $ABC$).
\end{propriete}

\courssub{Démarche et corrigé commenté}
\begin{exoresolu}[Corrigé détaillé de l'activité d'intégration]
\textbf{1. Repérage et construction géométrique.}
\begin{itemize}
    \item On trace le repère $(O;\vec{i},\vec{j})$ (1 unité = 100 m). On place $A(1,3)$, $B(7,1)$, $C(3,6)$.
    \item \emph{Barycentre partiel $K$ de $(A,2)$ et $(B,3)$ :} $K$ appartient au segment $[AB]$ et divise ce segment dans le rapport inverse des poids : $\frac{KA}{KB} = \frac{\beta}{\alpha} = \frac{3}{2}$. $K$ est plus proche de $B$ (poids 3) que de $A$ (poids 2). Construction : sur une demi-droite issue de $A$, reporter 5 fois une unité ; joindre le 5ème trait à $B$ ; tracer le parallèle par le 2ème trait coupant $[AB]$ en $K$.
    \item \emph{Barycentre $G$ de $(K,5)$ et $(C,1)$ :} $G$ appartient au segment $[KC]$ avec $\frac{GK}{GC} = \frac{1}{5}$. $G$ est très proche de $K$ (poids 5 contre 1). Construction similaire sur $[KC]$.
    \item Le triangle $ABC$ est tracé. $G$ se trouve visuellement à l'intérieur.
\end{itemize}

\begin{popfigure}
\begin{tikzpicture}[scale=0.7, >=Stealth]
    % Grille cadastre
    \draw[popline, step=1cm] (0,0) grid (8,7);
    % Axes
    \draw[->, thick] (0,0) -- (8.5,0) node[right] {$x$ (hm)};
    \draw[->, thick] (0,0) -- (0,7.5) node[above] {$y$ (hm)};
    \foreach \x in {1,...,8} \draw (\x,0.1) -- (\x,-0.1) node[below] {\small $\x$};
    \foreach \y in {1,...,7} \draw (0.1,\y) -- (-0.1,\y) node[left] {\small $\y$};
    
    % Triangle ABC
    \draw[popblue, thick] (1,3) -- (7,1) -- (3,6) -- cycle;
    
    % Points A, B, C
    \node[popblue, fill=popblue, circle, inner sep=2pt, label=above left:{$A(1,3)$}] at (1,3) {};
    \node[popblue, fill=popblue, circle, inner sep=2pt, label=below right:{$B(7,1)$}] at (7,1) {};
    \node[popblue, fill=popblue, circle, inner sep=2pt, label=above:{$C(3,6)$}] at (3,6) {};
    
    % Point K (4.6, 1.8)
    \node[poporange, fill=poporange, circle, inner sep=2pt, label=below:{$K(\frac{23}{5},\frac{9}{5})$}] (K) at (4.6,1.8) {};
    \draw[poporange, dashed] (1,3) -- (K) -- (7,1);
    
    % Point G (13/3 = 4.333..., 5/2 = 2.5)
    \node[popgreen, fill=popgreen, circle, inner sep=2pt, label=right:{$G(\frac{13}{3},\frac{5}{2})$}] (G) at (4.333,2.5) {};
    \draw[popgreen, dashed] (K) -- (G) -- (3,6);
    
    % Cercle ligne de niveau (Rayon ~1.47)
    \draw[popred, thick, dashed] (G) circle (1.47);
    \node[popred, anchor=south west] at (5.8, 3.97) {Zone de repli ($R\approx 147$ m)};
    
    % Légendes segments
    \node[poporange, rotate=15] at (3.5, 2.2) {\small $K \in [AB]$};
    \node[popgreen, rotate=45] at (3.8, 3.5) {\small $G \in [KC]$};
\end{tikzpicture}
\legende{Repérage cadastral : triangle $ABC$, barycentres partiels $K$ et $G$, zone de repli (cercle rouge pointillé). Unités : hectomètres (100 m).}
\end{popfigure}

\textbf{2. Calcul des coordonnées (méthode vectorielle).}
On utilise $\overrightarrow{OG} = \frac{\alpha\overrightarrow{OA}+\beta\overrightarrow{OB}+\gamma\overrightarrow{OC}}{\alpha+\beta+\gamma}$.
\[
\overrightarrow{OG} = \frac{2(1,3) + 3(7,1) + 1(3,6)}{6}
= \frac{(2,6) + (21,3) + (3,6)}{6}
= \frac{(26,15)}{6}
= \left(\frac{13}{3}\,;\,\frac{5}{2}\right).
\]
Donc $G\left(\frac{13}{3}\,;\,\frac{5}{2}\right) \approx G(4,33\,;\,2,50)$ (en unités de 100 m, soit environ 433 m ; 250 m sur le terrain).

\emph{Vérification par associativité :}
$K$ barycentre de $(A,2),(B,3)$ : $\overrightarrow{OK} = \frac{2(1,3)+3(7,1)}{5} = \frac{(23,9)}{5} = \left(\frac{23}{5}\,;\,\frac{9}{5}\right) = (4,6\,;\,1,8)$.
$G$ barycentre de $(K,5),(C,1)$ : $\overrightarrow{OG} = \frac{5(4,6;1,8) + 1(3,6)}{6} = \frac{(23,9)+(3,6)}{6} = \frac{(26,15)}{6} = \left(\frac{13}{3}\,;\,\frac{5}{2}\right)$. \textbf{OK.}

\emph{Justification : $G$ est à l'intérieur du triangle $ABC$.}
Les coefficients $\alpha=2,\beta=3,\gamma=1$ sont tous \textbf{strictement positifs}. D'après la propriété du barycentre à coefficients positifs, $G$ appartient à l'enveloppe convexe de $\{A,B,C\}$, c'est-à-dire à l'intérieur du triangle $ABC$ (ou sur ses côtés si alignés, ce qui n'est pas le cas ici car le déterminant $\det(\overrightarrow{AB},\overrightarrow{AC}) = -16 \neq 0$). On peut aussi vérifier que $G$ est du même côté que $C$ par rapport à $(AB)$, que $A$ par rapport à $(BC)$, et que $B$ par rapport à $(AC)$ en calculant les déterminants (omises ici pour la brièveté, mais à faire en examen).

\textbf{3. Calcul de la fatigue minimale $F(G)$.}
On utilise la relation de Leibniz : $\sum \alpha_i MA_i^2 = \sum \alpha_i GA_i^2 + 6\,MG^2$.
Pour $M=G$, le terme $MG^2$ s'annule. Il faut calculer $S_G = \sum \alpha_i GA_i^2 = 2\,GA^2 + 3\,GB^2 + 1\,GC^2$.

Calculs des carrés des distances (coordonnées exactes) :
\begin{align*}
GA^2 &= \left(\frac{13}{3}-1\right)^2 + \left(\frac{5}{2}-3\right)^2 = \left(\frac{10}{3}\right)^2 + \left(-\frac{1}{2}\right)^2 = \frac{100}{9} + \frac{1}{4} = \frac{400+9}{36} = \frac{409}{36}. \\[4pt]
GB^2 &= \left(\frac{13}{3}-7\right)^2 + \left(\frac{5}{2}-1\right)^2 = \left(-\frac{8}{3}\right)^2 + \left(\frac{3}{2}\right)^2 = \frac{64}{9} + \frac{9}{4} = \frac{256+81}{36} = \frac{337}{36}. \\[4pt]
GC^2 &= \left(\frac{13}{3}-3\right)^2 + \left(\frac{5}{2}-6\right)^2 = \left(\frac{4}{3}\right)^2 + \left(-\frac{7}{2}\right)^2 = \frac{16}{9} + \frac{49}{4} = \frac{64+441}{36} = \frac{505}{36}.
\end{align*}

Somme pondérée :
\[
S_G = 2\times\frac{409}{36} + 3\times\frac{337}{36} + 1\times\frac{505}{36}
= \frac{818 + 1011 + 505}{36}
= \frac{2334}{36}
= \frac{389}{6}.
\]

Fatigue moyenne minimale :
\[
F(G) = \frac{S_G}{6} = \frac{389}{36} \approx 10,806 \quad\text{(unités : (100 m)$^2$ = 10 000 m$^2$)}.
\]
En mètres carrés : $F(G) \approx 108\,060 \text{ m}^2$.

\textbf{4. Zone de repli : ligne de niveau $F(M) = 1,2 \times F(G)$.}
On cherche les points $M$ tels que $F(M) = 1,2 F(G)$.
D'après Leibniz : $F(M) = F(G) + MG^2$ (car $\sum\alpha_i = 6$).
L'équation devient :
\[
F(G) + MG^2 = 1,2\,F(G) \quad\Longleftrightarrow\quad MG^2 = 0,2\,F(G).
\]
C'est un \textbf{cercle de centre $G$ et de rayon $R = \sqrt{0,2 \times F(G)}$}.

Calcul de $R^2$ :
\[
R^2 = 0,2 \times \frac{389}{36} = \frac{1}{5} \times \frac{389}{36} = \frac{389}{180}.
\]
Rayon (en unités de 100 m) :
\[
R = \sqrt{\frac{389}{180}} \approx \sqrt{2,1611} \approx 1,47 \text{ unités}.
\]
Rayon sur le terrain : $R_{\text{terrain}} \approx 1,47 \times 100 \text{ m} = \textbf{147 m}$.

\emph{Équation cartésienne du cercle :}
Centre $G\left(\frac{13}{3},\frac{5}{2}\right)$, $R^2 = \frac{389}{180}$.
\[
\left(x - \frac{13}{3}\right)^2 + \left(y - \frac{5}{2}\right)^2 = \frac{389}{180}.
\]
Développement (forme réduite) :
\[
x^2 - \frac{26}{3}x + \frac{169}{9} + y^2 - 5y + \frac{25}{4} = \frac{389}{180}.
\]
Multiplier par 180 (PPCM de 3,9,4,180) :
\[
180x^2 - 1560x + 3380 + 180y^2 - 900y + 1125 = 389.
\]
\[
180x^2 + 180y^2 - 1560x - 900y + 4116 = 0.
\]
Diviser par 12 :
\[
\boxed{15x^2 + 15y^2 - 130x - 75y + 343 = 0}.
\]

\textbf{5. Note technique à l'attention de Monsieur le Maire de Yaoundé.}

\begin{center}
\textbf{OBJET : Proposition d'implantation de la borne-fontaine autonome — Marché Mokolo}
\end{center}

Monsieur le Maire,

Suite à l'étude de localisation pour la nouvelle borne-fontaine devant desservir les quartiers Mokolo 1, Mokolo 2 et Essos, nous vous soumettons les conclusions techniques suivantes.

\textbf{1. Point optimal $G$.}
En modélisant la « fatigue moyenne » des usagers par la moyenne quadratique des distances pondérées par la population (2000, 3000, 1000 hab.), le point optimal minimisant cette fatigue est le barycentre $G$ du système pondéré. Ses coordonnées cadastrales sont :
\[
G \approx (433 \text{ m} ; 250 \text{ m}) \quad \text{(repère origine mairie, axes Est/Nord)}.
\]
Ce point se situe \textbf{à l'intérieur du triangle} formé par les trois quartiers, ce qui garantit une équité de desserte : aucun quartier n'est défavorisé par une position excentrée. La fatigue moyenne minimale estimée est de $108\,060 \text{ m}^2$ (racine carrée $\approx 329 \text{ m}$ distance quadratique moyenne).

\textbf{2. Zone de repli (contrainte terrain).}
Si le point $G$ exact est impraticable, nous définissons une zone de tolérance où la fatigue n'augmente pas de plus de 20\%. Cette zone est un \textbf{disque de centre $G$ et de rayon 147 m}.
Tout point à l'intérieur de ce disque (ex : décalage de 100 m vers la voirie existante) maintient la qualité de service dans les normes acceptables. L'équation du cercle limite est $15x^2+15y^2-130x-75y+343=0$ (coordonnées en hectomètres).

\textbf{3. Recommandation.}
Nous préconisons :
\begin{itemize}
    \item Une vérification topographique \textit{in situ} au point $G$ (coordonnées GPS : à convertir depuis le cadastre).
    \item Si obstacle majeur : reporter l'implantation sur le point accessible le plus proche de $G$ \textbf{à l'intérieur du disque de rayon 147 m}, de préférence sur la voirie publique existante pour réduire les coûts de raccordement.
    \item Une concertation avec les délégués de quartiers pour valider l'acceptabilité sociale du site retenu.
\end{itemize}

Cette démarche, fondée sur une modélisation mathématique rigoureuse (barycentres, optimisation quadratique), assure une décision objective, transparente et équitable pour les populations concernées.

Respectueusement,
\textit{Le Bureau d'Études Techniques — Division Mathématiques Appliquées}
\end{exoresolu}

\courssub{Exercice d'application : Extension à quatre points}
\begin{exercice}[Barycentre de quatre points — Cas d'un quatrième quartier]
On souhaite intégrer le quartier \textbf{Nlongkak} (point $D$) au dispositif. Sa population est de 4000 habitants (poids $\delta = 4$). Ses coordonnées cadastrales sont $D(5\,;\,4)$.
\begin{enumerate}
    \item Vérifier que le barycentre $G'$ du système $\{(A,2),(B,3),(C,1),(D,4)\}$ existe.
    \item Calculer les coordonnées exactes de $G'$ par la formule générale à 4 points.
    \item Déterminer le nouveau rayon de la zone de repli (tolérance 20\%) pour ce système à 4 points. Comparer avec le rayon précédent.
\end{enumerate}
\end{exercice}

\begin{corrige}[Exercice d'application]
\begin{enumerate}
    \item Somme des poids : $2+3+1+4 = 10 \neq 0$. Le barycentre $G'$ existe et est unique.
    \item $\overrightarrow{OG'} = \frac{2(1,3)+3(7,1)+1(3,6)+4(5,4)}{10} = \frac{(2,6)+(21,3)+(3,6)+(20,16)}{10} = \frac{(46,31)}{10} = \left(\frac{23}{5}\,;\,\frac{31}{10}\right) = (4,6\,;\,3,1)$.
    \item Calcul de $F(G')$ via Leibniz (somme pondérée des carrés des distances à $G'$) puis $R'^2 = 0,2 F(G')$. Le nouveau centre $G'(4,6;3,1)$ est plus proche de $D$ (poids fort). Le rayon change car la dispersion pondérée change. (Calcul numérique laissé en exercice d'entraînement).
\end{enumerate}
\end{corrige}

\begin{aretenir}[Synthèse et portée de la leçon]
Cette activité a permis de mettre en évidence :
\begin{itemize}
    \item La \textbf{puissance de la modélisation par barycentres} : un problème d'aménagement urbain (choix d'un emplacement équitable) se traduit exactement par la recherche du barycentre de points pondérés par les populations.
    \item L'\textbf{associativité} comme outil de construction géométrique (barycentres partiels) et de vérification numérique (calculs par étapes).
    \item Le \textbf{lien fondamental entre barycentre et optimisation quadratique} via la relation de Leibniz : le barycentre minimise la somme des carrés des distances pondérées ; les lignes de niveau sont des cercles concentriques au barycentre.
    \item L'importance de la \textbf{restitution technique} : traduire des résultats mathématiques (coordonnées, équation de cercle, rayon) en recommandations opérationnelles (coordonnées GPS, zone de tolérance en mètres, concertation) pour un décideur non spécialiste.
    \item La \textbf{cohérence du programme} : les notions de Première (barycentre 2, 3, 4 points ; associativité ; lignes de niveau ; relation de Leibniz) s'enchaînent naturellement pour résoudre une situation complexe unique.
\end{itemize}
\end{aretenir}

\begin{attention}[Pièges fréquents relevés lors de l'activité]
\begin{itemize}
\item Oublier de vérifier $\sum\alpha_i \neq 0$ avant de calculer le barycentre.
\item Confondre le rapport de division du segment ($KA/KB = \beta/\alpha$) et les coefficients dans la formule vectorielle.
\item Ne pas justifier que $G$ est \emph{à l'intérieur} du triangle (coefficients tous positifs) : une simple phrase « $G$ est dans le triangle » sans justification vaut 0 point au Probatoire.
\item Oublier de diviser par la somme des poids pour obtenir la \emph{moyenne} $F(G)$ après avoir calculé la somme pondérée $S_G$.
\item Utiliser la virgule décimale dans le code TikZ/pgfplots (utiliser le point : 1.5 et non 1,5).
\item Ne pas convertir les unités (hectomètres $\rightarrow$ mètres) pour la réponse finale concrète (rayon en mètres).
\end{itemize}
\end{attention}

\newpage

\coursec{Méthodes et exercices résolus}

\coursec{Barycentres}

\begin{definition}[Barycentre de $n$ points pondérés]
Soient $A_1, A_2, \dots, A_n$ des points du plan (ou de l'espace) et $a_1, a_2, \dots, a_n$ des nombres réels (coefficients ou masses). On appelle \textbf{barycentre} du système de points pondérés $\{(A_1;a_1), (A_2;a_2), \dots, (A_n;a_n)\}$ le point $G$, s'il existe, tel que :
\[ a_1\,\overrightarrow{GA_1} + a_2\,\overrightarrow{GA_2} + \dots + a_n\,\overrightarrow{GA_n} = \vec{0}. \]
On note $G = \mathrm{bar}\{(A_1;a_1), (A_2;a_2), \dots, (A_n;a_n)\}$.
La condition d'existence et d'unicité est $\sum_{i=1}^n a_i \neq 0$.
\end{definition}

\begin{propriete}[Propriétés fondamentales du barycentre]
\begin{enumerate}
\item \textbf{Existence et unicité :} Si $\sum a_i \neq 0$, le barycentre existe et est unique. Si $\sum a_i = 0$, soit il n'existe pas, soit il y en a une infinité (cas non au programme de Première).
\item \textbf{Homogénéité :} $\forall k \in \mathbb{R}^*$, $\mathrm{bar}\{(A_1;ka_1), \dots, (A_n;ka_n)\} = \mathrm{bar}\{(A_1;a_1), \dots, (A_n;a_n)\}$.
\item \textbf{Associativité (ou partialité) :} Si l'on regroupe une partie des points en leur barycentre partiel (somme de leurs coefficients non nulle), le barycentre global ne change pas.
\item \textbf{Relation de Chasles (outil pour les lignes de niveau) :} Pour tout point $M$,
\[ \sum a_i\,\overrightarrow{MA_i} = \left(\sum a_i\right)\overrightarrow{MG} + \sum a_i\,\overrightarrow{GA_i} = \left(\sum a_i\right)\overrightarrow{MG}. \]
\end{enumerate}
\end{propriete}

\begin{aretenir}[Vocabulaire du Probatoire]
\begin{itemize}
\item \textbf{Isobarycentre :} barycentre dont tous les coefficients sont égaux (souvent pris égaux à $1$). Pour deux points, c'est le \textbf{milieu} du segment. Pour trois points non alignés, c'est le \textbf{centre de gravité} du triangle (concours des médianes).
\item \textbf{Coefficients pondérés :} les nombres $a_i$ (masses, effectifs, montants, coefficients de frottement, etc.).
\item \textbf{Ligne de niveau :} ensemble des points $M$ vérifiant une condition sur la norme $\left\|\sum a_i\,\overrightarrow{MA_i}\right\|$ (cercle, droite, médiatrice).
\end{itemize}
\end{aretenir}

\courssub{Méthode 1 : Construire le barycentre de deux points pondérés}

\begin{methode}[Construire $G = \mathrm{bar}\{(A\,;a),(B\,;b)\}$ avec $a+b\neq 0$]
\begin{enumerate}
\item Vérifier que la somme des coefficients est non nulle : $a+b\neq 0$. Si $a+b=0$, le barycentre n'existe pas (vecteur nul impossible sauf $A=B$).
\item Écrire la relation vectorielle fondamentale :
\[ a\,\overrightarrow{GA}+b\,\overrightarrow{GB}=\vec{0}. \]
\item Transformer pour exprimer $\overrightarrow{AG}$ en fonction de $\overrightarrow{AB}$ en posant $\overrightarrow{GB}=\overrightarrow{GA}+\overrightarrow{AB}$ (relation de Chasles) :
\[ \overrightarrow{AG}=\dfrac{b}{a+b}\,\overrightarrow{AB}. \]
\item Calculer le coefficient $k=\dfrac{b}{a+b}$, puis placer $G$ sur la droite $(AB)$ :
\begin{itemize}
\item si $0<k<1$, $G$ est \textbf{entre} $A$ et $B$ (segment $[AB]$) ;
\item si $k<0$, $G$ est du côté de $A$ (en dehors du segment $[AB]$) ;
\item si $k>1$, $G$ est du côté de $B$ (en dehors du segment $[AB]$).
\end{itemize}
\item \textbf{Réflexe physique :} le barycentre est toujours plus proche du point dont le coefficient a la plus grande \textbf{valeur absolue}.
\end{enumerate}
\end{methode}

\begin{popfigure}
\begin{tikzpicture}[scale=0.8, >=Stealth]
% Points A, B, G
\coordinate (A) at (0,0);
\coordinate (B) at (6,0);
\coordinate (G) at (2.4,0); % k = 2/5 -> AG = 2.4
% Line
\draw[thick, popdark] (A) -- (B) node[midway, below=2mm] {$AB = 60$ cm};
% Points
\fill[popblue] (A) circle (2.5pt) node[below] {$A$ (3 kg)};
\fill[poporange] (B) circle (2.5pt) node[below] {$B$ (2 kg)};
\fill[popgreen] (G) circle (2.5pt) node[above] {$G$};
% Distances
\draw[<->, popmuted] (0,-0.8) -- (2.4,-0.8) node[midway, below] {$AG = 24$ cm};
\draw[<->, popmuted] (2.4,-0.8) -- (6,-0.8) node[midway, below] {$GB = 36$ cm};
% Masses representation
\draw[popblue, thick] (A) -- ++(0,1.5) node[above] {$3$};
\draw[poporange, thick] (B) -- ++(0,1.5) node[above] {$2$};
\end{tikzpicture}
\legende{Équilibre d'une barre : le point $G$ est plus proche de la masse la plus grande.}
\end{popfigure}

\begin{exoresolu}[Placement d'une charge sur une barre — Application atelier]
Une barre métallique $AB$ de longueur $\SI{60}{cm}$ porte en $A$ une masse de $3$ kg et en $B$ une masse de $2$ kg. On note $G$ le point d'équilibre de la barre, barycentre de $(A\,;3)$ et $(B\,;2)$.
\begin{enumerate}
\item Justifier que $G$ existe et exprimer $\overrightarrow{AG}$ en fonction de $\overrightarrow{AB}$.
\item Calculer la distance $AG$.
\end{enumerate}
\tcblower
\textbf{1.} La somme des coefficients est $a+b=3+2=5\neq 0$, donc le barycentre $G$ de $(A\,;3)$ et $(B\,;2)$ existe et est unique.

Par définition : $3\,\overrightarrow{GA}+2\,\overrightarrow{GB}=\vec{0}$.

Or $\overrightarrow{GB}=\overrightarrow{GA}+\overrightarrow{AB}$ (relation de Chasles). Donc :
\begin{align*}
3\,\overrightarrow{GA}+2\left(\overrightarrow{GA}+\overrightarrow{AB}\right)&=\vec{0}\\
5\,\overrightarrow{GA}+2\,\overrightarrow{AB}&=\vec{0}\\
\overrightarrow{GA}&=-\dfrac{2}{5}\,\overrightarrow{AB}\\
\overrightarrow{AG}&=\dfrac{2}{5}\,\overrightarrow{AB}.
\end{align*}

\textbf{2.} Comme $k=\dfrac{2}{5}\in ]0,1[$, les vecteurs $\overrightarrow{AG}$ et $\overrightarrow{AB}$ ont le même sens et $G \in [AB]$, donc :
\[ AG=\dfrac{2}{5}\,AB=\dfrac{2}{5}\times 60=24\ \text{cm}. \]

\textbf{Conclusion :} le point d'équilibre $G$ se trouve sur le segment $[AB]$, à $\SI{24}{cm}$ de $A$ (donc à $\SI{36}{cm}$ de $B$). On remarque que $G$ est plus proche de $A$, le point qui porte la plus grande masse : c'est cohérent avec la physique du levier.
\end{exoresolu}

\courssub{Méthode 2 : Déterminer et utiliser les coordonnées du barycentre}

\begin{methode}[Coordonnées du barycentre dans un repère]
\begin{enumerate}
\item Vérifier que la somme des coefficients est non nulle ($\sum a_i \neq 0$).
\item Noter les coordonnées des points : $A(x_A\,;y_A)$, $B(x_B\,;y_B)$, $C(x_C\,;y_C)$, \dots
\item Appliquer les formules (moyennes pondérées) :
\begin{itemize}
\item 2 points : $x_G=\dfrac{a\,x_A+b\,x_B}{a+b},\quad y_G=\dfrac{a\,y_A+b\,y_B}{a+b}$.
\item 3 points : $x_G=\dfrac{a\,x_A+b\,x_B+c\,x_C}{a+b+c},\quad y_G=\dfrac{a\,y_A+b\,y_B+c\,y_C}{a+b+c}$.
\item $n$ points : $x_G=\dfrac{\sum a_i x_i}{\sum a_i},\quad y_G=\dfrac{\sum a_i y_i}{\sum a_i}$.
\end{itemize}
\item \textbf{Réflexe de vérification :} si tous les coefficients sont égaux (isobarycentre), on obtient simplement la moyenne arithmétique des coordonnées.
\end{enumerate}
\end{methode}

\begin{exoresolu}[Calcul de coordonnées — Gestion de stocks / Géométrie]
Dans un repère orthonormé $(O\,;\vec{i},\vec{j})$, on donne les points $A(2\,;-1)$, $B(5\,;3)$ et $C(-1\,;4)$.
\begin{enumerate}
\item Déterminer les coordonnées du barycentre $G$ de $(A\,;1)$, $(B\,;2)$ et $(C\,;1)$.
\item Déterminer les coordonnées de l'isobarycentre $I$ du triangle $ABC$.
\end{enumerate}
\tcblower
\textbf{1.} La somme des coefficients est $1+2+1=4\neq 0$ : le barycentre $G$ existe.

\begin{align*}
x_G&=\dfrac{1\times 2+2\times 5+1\times (-1)}{4}=\dfrac{2+10-1}{4}=\dfrac{11}{4}=2{,}75 ;\\[4pt]
y_G&=\dfrac{1\times (-1)+2\times 3+1\times 4}{4}=\dfrac{-1+6+4}{4}=\dfrac{9}{4}=2{,}25.
\end{align*}

Donc $G\left(\dfrac{11}{4}\,;\dfrac{9}{4}\right)$.

\textbf{2.} L'isobarycentre correspond à des coefficients tous égaux à $1$ (somme $=3\neq 0$) :
\[ x_I=\dfrac{2+5+(-1)}{3}=\dfrac{6}{3}=2, \qquad y_I=\dfrac{-1+3+4}{3}=\dfrac{6}{3}=2. \]

Donc $I(2\,;2)$.

\textbf{Conclusion :} $G\left(\dfrac{11}{4}\,;\dfrac{9}{4}\right)$ et $I(2\,;2)$. L'isobarycentre $I$ est le centre de gravité du triangle $ABC$, point de concours de ses médianes.
\end{exoresolu}

\courssub{Méthode 3 : Utiliser l'homogénéité et l'associativité du barycentre}

\begin{methode}[Simplifier et regrouper les points pondérés (3 à 4 points)]
\begin{enumerate}
\item \textbf{Homogénéité :} on peut multiplier (ou diviser) tous les coefficients par un même nombre non nul sans changer le barycentre. \emph{Réflexe :} se ramener à des coefficients entiers simples (ex. $(A\,;0{,}5),(B\,;1{,}5)$ devient $(A\,;1),(B\,;3)$).
\item \textbf{Associativité :} pour construire le barycentre de trois points $(A\,;a)$, $(B\,;b)$, $(C\,;c)$ avec $a+b+c\neq 0$ :
\begin{itemize}
\item construire d'abord le barycentre partiel $H$ de deux points, par exemple $H=\mathrm{bar}\{(A\,;a),(B\,;b)\}$ (vérifier $a+b\neq 0$) ;
\item affecter à $H$ la somme $a+b$ ;
\item alors $G=\mathrm{bar}\{(H\,;a+b),(C\,;c)\}$, et $G$ est sur la droite $(HC)$.
\end{itemize}
\item \textbf{Extension à 4 points :} même principe, on regroupe par paires ou par triplets. Exemple : $G=\mathrm{bar}\{(A;a),(B;b),(C;c),(D;d)\} = \mathrm{bar}\{(H_{AB};a+b),(H_{CD};c+d)\}$.
\item Conclure en plaçant $G$ grâce à la relation $\overrightarrow{HG}=\dfrac{c}{a+b+c}\,\overrightarrow{HC}$ (ou formule analogue).
\end{enumerate}
\end{methode}

\begin{popfigure}
\begin{tikzpicture}[scale=0.9, >=Stealth]
% Triangle ABC
\coordinate (A) at (0,3);
\coordinate (B) at (-2,0);
\coordinate (C) at (4,0);
% Midpoint I of BC
\coordinate (I) at (1,0);
% G midpoint of AI
\coordinate (G) at (0.5,1.5);
% Draw triangle
\draw[thick, popdark] (A) -- (B) -- (C) -- cycle;
% Median AI
\draw[popblue, dashed] (A) -- (I) node[midway, right] {Médiane $(AI)$};
% Points
\fill[popblue] (A) circle (2pt) node[above left] {$A(2)$};
\fill[poporange] (B) circle (2pt) node[below left] {$B(1)$};
\fill[poporange] (C) circle (2pt) node[below right] {$C(1)$};
\fill[poppurple] (I) circle (2pt) node[below] {$I$ (milieu $[BC]$)};
\fill[popgreen] (G) circle (2pt) node[above right] {$G$};
% Labels masses
\node[popblue] at (-0.5, 3.5) {Coeff 2};
\node[poporange] at (-2.5, -0.5) {Coeff 1};
\node[poporange] at (4.5, -0.5) {Coeff 1};
\node[poppurple] at (1.5, -0.5) {Somme 2};
\end{tikzpicture}
\legende{Construction de $G = \mathrm{bar}\{(A;2),(B;1),(C;1)\}$ par barycentre partiel $I$.}
\end{popfigure}

\begin{exoresolu}[Construction par barycentres partiels — Méthode associativité]
Soit $ABC$ un triangle. On considère $G=\mathrm{bar}\{(A\,;2),(B\,;1),(C\,;1)\}$.
\begin{enumerate}
\item Montrer que $G$ est le barycentre de $(A\,;2)$ et $(I\,;2)$, où $I$ est le milieu de $[BC]$.
\item En déduire la position de $G$ et construire la figure.
\end{enumerate}
\tcblower
\textbf{1.} La somme totale des coefficients est $2+1+1=4\neq 0$, donc $G$ existe.

Soit $I$ le milieu de $[BC]$. Alors $I=\mathrm{bar}\{(B\,;1),(C\,;1)\}$ car des coefficients égaux donnent l'isobarycentre, c'est-à-dire le milieu du segment. La somme des coefficients de $B$ et $C$ est $1+1=2\neq 0$.

Par la propriété d'\textbf{associativité} du barycentre :
\[ G=\mathrm{bar}\{(A\,;2),(B\,;1),(C\,;1)\}=\mathrm{bar}\{(A\,;2),(I\,;1+1)\}=\mathrm{bar}\{(A\,;2),(I\,;2)\}. \]

\textbf{2.} Par homogénéité (division par 2), $\mathrm{bar}\{(A\,;2),(I\,;2)\}=\mathrm{bar}\{(A\,;1),(I\,;1)\}$ : c'est le \textbf{milieu du segment $[AI]$}.

Or $I$ est le milieu de $[BC]$, donc $(AI)$ est la médiane issue de $A$ dans le triangle $ABC$.

\textbf{Conclusion :} $G$ est le milieu de la médiane $[AI]$. Pour la construction : on place $I$ milieu de $[BC]$, on trace le segment $[AI]$, puis on place $G$ au milieu de $[AI]$.
\end{exoresolu}

\courssub{Méthode 4 : Démontrer un alignement ou un point de concours avec les barycentres}

\begin{methode}[Prouver que trois points sont alignés ou que des droites concourent]
\begin{enumerate}
\item \textbf{Alignement :} pour montrer que $G$ appartient à une droite $(XY)$, exprimer $G$ comme barycentre de points situés sur $(XY)$, ou montrer que $\overrightarrow{XG}$ et $\overrightarrow{XY}$ sont colinéaires grâce à la relation $a\,\overrightarrow{GA}+b\,\overrightarrow{GB}=\vec{0}$ évaluée en un point bien choisi.
\item \textbf{Concours :} pour montrer que trois droites se coupent en un point $G$, montrer que $G$ est le barycentre de systèmes de points définissant chacune des droites (technique des barycentres partiels).
\item \textbf{Rédaction exigée au Probatoire :} toujours annoncer l'hypothèse « somme des coefficients non nulle » avant d'utiliser un barycentre. Écrire explicitement la relation vectorielle définissante.
\end{enumerate}
\end{methode}

\begin{popfigure}
\begin{tikzpicture}[scale=0.8, >=Stealth]
% Triangle ABC
\coordinate (A) at (0,3.5);
\coordinate (B) at (-3,0);
\coordinate (C) at (3,0);
% Midpoints
\coordinate (Ap) at (0,0); % Midpoint BC
\coordinate (Bp) at (1.5,1.75); % Midpoint CA
\coordinate (Cp) at (-1.5,1.75); % Midpoint AB
% Centroid G
\coordinate (G) at (0, 1.166); % 2/3 from A
% Draw triangle
\draw[thick, popdark] (A) -- (B) -- (C) -- cycle;
% Medians
\draw[popblue, dashed] (A) -- (Ap);
\draw[poporange, dashed] (B) -- (Bp);
\draw[poppurple, dashed] (C) -- (Cp);
% Points
\fill[popblue] (A) circle (2pt) node[above] {$A$};
\fill[poporange] (B) circle (2pt) node[left] {$B$};
\fill[poppurple] (C) circle (2pt) node[right] {$C$};
\fill[popmuted] (Ap) circle (2pt) node[below] {$A'$};
\fill[popmuted] (Bp) circle (2pt) node[right] {$B'$};
\fill[popmuted] (Cp) circle (2pt) node[left] {$C'$};
\fill[popgreen] (G) circle (3pt) node[above right] {$G$ (Isobarycentre)};
% Ratios
\node[popblue] at (-1, 2.5) {$\frac{2}{3}$};
\node[popblue] at (1, 2.5) {$\frac{1}{3}$};
\end{tikzpicture}
\legende{Concours des médianes en l'isobarycentre $G$ (centre de gravité).}
\end{popfigure}

\begin{exoresolu}[Les médianes d'un triangle sont concourantes — Démonstration classique]
Soit $ABC$ un triangle, $A'$, $B'$, $C'$ les milieux respectifs de $[BC]$, $[CA]$, $[AB]$. Soit $G$ l'isobarycentre de $A$, $B$, $C$.
\begin{enumerate}
\item Montrer que $G=\mathrm{bar}\{(A\,;1),(A'\,;2)\}$.
\item En déduire que les trois médianes du triangle sont concourantes en $G$ et préciser la position de $G$ sur chaque médiane.
\end{enumerate}
\tcblower
\textbf{1.} $A'$ est le milieu de $[BC]$, donc $A'=\mathrm{bar}\{(B\,;1),(C\,;1)\}$ avec la somme $1+1=2\neq 0$.

Par associativité (somme totale $1+1+1=3\neq 0$) :
\[ G=\mathrm{bar}\{(A\,;1),(B\,;1),(C\,;1)\}=\mathrm{bar}\{(A\,;1),(A'\,;2)\}. \]

\textbf{2.} D'après la question 1, $G$ est le barycentre de $(A\,;1)$ et $(A'\,;2)$, donc $G$ appartient à la droite $(AA')$, qui est la médiane issue de $A$. De plus, en utilisant la formule de construction à deux points :
\[ \overrightarrow{AG}=\dfrac{2}{1+2}\,\overrightarrow{AA'}=\dfrac{2}{3}\,\overrightarrow{AA'}. \]

Le même raisonnement avec $B'$ et $C'$ donne :
\[ G=\mathrm{bar}\{(B\,;1),(B'\,;2)\} \quad\text{et}\quad G=\mathrm{bar}\{(C\,;1),(C'\,;2)\}, \]
donc $G$ appartient aussi aux médianes $(BB')$ et $(CC')$, avec $\overrightarrow{BG}=\dfrac{2}{3}\overrightarrow{BB'}$ et $\overrightarrow{CG}=\dfrac{2}{3}\overrightarrow{CC'}$.

\textbf{Conclusion :} les trois médianes du triangle $ABC$ sont concourantes en l'isobarycentre $G$ (centre de gravité), situé aux deux tiers de chaque médiane en partant du sommet.
\end{exoresolu}

\courssub{Méthode 5 : Déterminer un ensemble de points (lignes de niveau)}

\begin{methode}[Ensembles de points définis par une norme de somme vectorielle]
Soit $G=\mathrm{bar}\{(A\,;a),(B\,;b)\}$ avec $a+b\neq 0$.
\begin{enumerate}
\item \textbf{Réduire la somme de vecteurs :} pour tout point $M$ du plan, la relation de Chasles appliquée au barycentre donne :
\[ a\,\overrightarrow{MA}+b\,\overrightarrow{MB}=(a+b)\,\overrightarrow{MG}. \]
C'est l'outil central : on remplace une somme compliquée par un seul vecteur.
\item \textbf{Ensemble $\left\|a\,\overrightarrow{MA}+b\,\overrightarrow{MB}\right\|=r$ (avec $r\ge 0$) :} on obtient $(a+b)\,MG=r$, soit $MG=\dfrac{r}{|a+b|}$ : c'est un \textbf{cercle} de centre $G$ et de rayon $\dfrac{r}{|a+b|}$.
\item \textbf{Ensemble $\left\|a\,\overrightarrow{MA}+b\,\overrightarrow{MB}\right\|=\left\|c\,\overrightarrow{MC}+d\,\overrightarrow{MD}\right\|$ :} on réduit chaque membre avec son barycentre ($G$ et $G'$) ; on obtient $|a+b|\,MG = |c+d|\,MG'$.
\begin{itemize}
\item Si $|a+b| = |c+d|$, c'est la \textbf{médiatrice} de $[GG']$.
\item Sinon, c'est un \textbf{cercle} (cercle d'Apollonius).
\end{itemize}
\item Conclure en précisant la nature géométrique, le centre et le rayon, ou la droite obtenue.
\end{enumerate}
\end{methode}

\begin{popfigure}
\begin{tikzpicture}[scale=0.7, >=Stealth]
% Points A, B, G
\coordinate (A) at (0,0);
\coordinate (B) at (8,0);
\coordinate (G) at (2,0); % AG = 1/4 AB = 2
% Circle center G radius 3
\draw[popgreen, thick] (G) circle (3);
% Line AB
\draw[thick, popdark] (A) -- (B);
% Points
\fill[popblue] (A) circle (2.5pt) node[below] {$A$};
\fill[poporange] (B) circle (2.5pt) node[below] {$B$};
\fill[popgreen] (G) circle (2.5pt) node[above] {$G$};
% Radius
\draw[<->, popmuted] (2,0) -- (5,0) node[midway, below] {$R=3$};
% Distances
\draw[<->, popmuted] (0,-0.8) -- (2,-0.8) node[midway, below] {$AG=2$};
\draw[<->, popmuted] (2,-0.8) -- (8,-0.8) node[midway, below] {$GB=6$};
\end{tikzpicture}
\legende{Ligne de niveau : cercle de centre $G$ (barycentre) et de rayon $3$.}
\end{popfigure}

\begin{exoresolu}[Ligne de niveau : un cercle]
Soit $A$ et $B$ deux points tels que $AB=8$ cm. On note $G=\mathrm{bar}\{(A\,;3),(B\,;1)\}$.
\begin{enumerate}
\item Montrer que pour tout point $M$ du plan : $3\,\overrightarrow{MA}+\overrightarrow{MB}=4\,\overrightarrow{MG}$.
\item Déterminer l'ensemble $(\mathcal{E})$ des points $M$ tels que $\left\|3\,\overrightarrow{MA}+\overrightarrow{MB}\right\|=12$.
\end{enumerate}
\tcblower
\textbf{1.} La somme des coefficients est $3+1=4\neq 0$, donc $G$ existe. Pour tout point $M$, la relation de Chasles donne :
\begin{align*}
3\,\overrightarrow{MA}+\overrightarrow{MB}
&=3\left(\overrightarrow{MG}+\overrightarrow{GA}\right)+\left(\overrightarrow{MG}+\overrightarrow{GB}\right)\\
&=4\,\overrightarrow{MG}+\underbrace{\left(3\,\overrightarrow{GA}+\overrightarrow{GB}\right)}_{=\,\vec{0}\ \text{car } G \text{ est le barycentre}}\\
&=4\,\overrightarrow{MG}.
\end{align*}

\textbf{2.} D'après la question 1 :
\[ \left\|3\,\overrightarrow{MA}+\overrightarrow{MB}\right\|=12 \iff \left\|4\,\overrightarrow{MG}\right\|=12 \iff 4\,MG=12 \iff MG=3. \]

L'ensemble des points situés à une distance fixe $3$ cm du point $G$ est le cercle de centre $G$ et de rayon $3$ cm.

\textbf{Conclusion :} $(\mathcal{E})$ est le cercle de centre $G$ (barycentre de $(A\,;3),(B\,;1)$, situé sur $[AB]$ avec $AG=\dfrac{1}{4}AB=2$ cm) et de rayon $3$ cm.
\end{exoresolu}

\courssub{Méthode 6 : Résoudre un problème concret à l'aide du barycentre}

\begin{methode}[Modéliser une situation concrète (Métiers techniques)]
\begin{enumerate}
\item \textbf{Identifier} les points (positions des masses, des charges, des capitaux, des centres de coûts) et les coefficients (masses en kg, montants en FCFA, effectifs, puissances).
\item \textbf{Vérifier} que la somme des coefficients est non nulle, puis écrire le barycentre.
\item \textbf{Calculer} la position par la formule vectorielle $\overrightarrow{AG}=\dfrac{b}{a+b}\overrightarrow{AB}$ ou par les coordonnées (moyennes pondérées).
\item \textbf{Interpréter} le résultat dans le contexte (point d'équilibre d'une poutre, centre de gravité d'une plaque, prix moyen pondéré, centre de charge électrique) et rédiger une phrase de conclusion adaptée à la situation professionnelle.
\end{enumerate}
\end{methode}

\begin{savaistu}[Le barycentre au Cameroun : du chantier à la gestion]
\begin{itemize}
\item \textbf{BTP / Charpente :} Le centre de gravité d'une poutre en béton armé ou d'une ferme métallique détermine les points d'élingage pour la grue. Un mauvais calcul = basculement de la charge.
\item \textbf{Électricité :} Le barycentre de charges ponctuelles (coefficients = valeurs des charges en Coulomb) donne le centre de charge. Pour des dipôles, on utilise des coefficients signés.
\item \textbf{Gestion / Comptabilité :} Le prix de revient moyen d'un produit fabriqué sur plusieurs sites est le barycentre des coûts unitaires pondérés par les quantités produites. $\text{PR Moyen} = \frac{\sum q_i \times \text{PR}_i}{\sum q_i}$.
\item \textbf{Histoire :} August Ferdinand Möbius (1827) a formalisé la notion de « centre de masse » (barycentre) pour la géométrie analytique, unifiant statique et géométrie.
\end{itemize}
\end{savaistu}

\begin{exoresolu}[Équilibre d'une plaque sur un chantier — Application BTP]
Une plaque triangulaire homogène $ABC$ est posée horizontalement. On veut la percer en son centre de gravité $G$ pour la suspendre à l'équilibre. Dans un repère adapté, $A(0\,;0)$, $B(6\,;0)$ et $C(0\,;4)$ (unité : dm).
\begin{enumerate}
\item Quel barycentre représente le centre de gravité d'une plaque triangulaire homogène ? Justifier.
\item Déterminer les coordonnées du point de perçage $G$.
\end{enumerate}
\tcblower
\textbf{1.} Pour une plaque triangulaire \textbf{homogène} (masse uniformément répartie), le centre de gravité est l'\textbf{isobarycentre} des trois sommets, c'est-à-dire le barycentre de $(A\,;1)$, $(B\,;1)$, $(C\,;1)$. En effet, l'homogénéité de la plaque impose que les trois sommets aient le même poids dans la modélisation discrète. La somme des coefficients est $3\neq 0$, donc ce barycentre existe.

\textbf{2.} On applique les formules des coordonnées de l'isobarycentre :
\begin{align*}
x_G&=\dfrac{x_A+x_B+x_C}{3}=\dfrac{0+6+0}{3}=2 ;\\[4pt]
y_G&=\dfrac{y_A+y_B+y_C}{3}=\dfrac{0+0+4}{3}=\dfrac{4}{3}.
\end{align*}

\textbf{Conclusion :} le point de perçage est $G\left(2\,;\dfrac{4}{3}\right)$, c'est-à-dire à $2$ dm du bord $(AC)$ et à environ $1{,}33$ dm du bord $(AB)$. Suspendue en ce point, la plaque restera horizontale.
\end{exoresolu}

\begin{exercice}[Entraînement au Probatoire — Série complète]
\begin{enumerate}
\item \textbf{Existence :} Soient $A$, $B$ deux points distincts. Pour quelles valeurs de $k$ le barycentre de $(A\,;k)$ et $(B\,;2)$ existe-t-il ? Le construire pour $k=-2$, $k=1$, $k=4$.
\item \textbf{Coordonnées :} Dans un repère, $A(-1;2)$, $B(3;0)$, $C(1;4)$. Calculer les coordonnées de $G = \mathrm{bar}\{(A;2), (B;-1), (C;3)\}$. Vérifier l'existence.
\item \textbf{Alignement :} Soit $G = \mathrm{bar}\{(A;2), (B;3), (C;-1)\}$. Montrer que $G$ appartient à la droite $(AB)$.
\item \textbf{Ligne de niveau :} $A(0;0)$, $B(4;0)$. $G = \mathrm{bar}\{(A;1), (B;3)\}$. Déterminer l'ensemble des points $M$ tels que $\|\overrightarrow{MA} + 3\overrightarrow{MB}\| = 8$.
\item \textbf{Problème concret :} Un camion charge des sacs de ciment : $50$ sacs à l'avant (point $A$), $30$ sacs à l'arrière (point $B$), distance $AB = 4$ m. Où placer le point d'appui central (barycentre) pour que le camion soit équilibré ? (On suppose les sacs de masse identique).
\end{enumerate}
\end{exercice}

\begin{corrige}[Exercice n°1 -- Entraînement Probatoire]
\begin{enumerate}
\item \textbf{Existence :} Somme $= k+2$. Existe si $k \neq -2$.
\begin{itemize}
\item $k=-2$ : somme $=0$, pas de barycentre (vecteur nul impossible car $A\neq B$).
\item $k=1$ : $G = \mathrm{bar}\{(A;1),(B;2)\}$, $k_{constr} = 2/3 \in ]0,1[$, $G \in [AB]$, $AG = \frac{2}{3}AB$.
\item $k=4$ : $G = \mathrm{bar}\{(A;4),(B;2)\}$, $k_{constr} = 2/6 = 1/3 \in ]0,1[$, $G \in [AB]$, plus proche de $A$ (coeff 4 > 2).
\end{itemize}
\item \textbf{Coordonnées :} Somme $= 2 + (-1) + 3 = 4 \neq 0$. Existe.
$x_G = \frac{2(-1) + (-1)(3) + 3(1)}{4} = \frac{-2 -3 + 3}{4} = -\frac{1}{2}$.
$y_G = \frac{2(2) + (-1)(0) + 3(4)}{4} = \frac{4 + 0 + 12}{4} = 4$.
$G(-0,5 ; 4)$.
\item \textbf{Alignement :} Somme $= 2+3-1=4\neq 0$. $G = \mathrm{bar}\{(A;2), (B;3), (C;-1)\}$.
Regrouper $A,B$ : $H = \mathrm{bar}\{(A;2),(B;3)\}$ (somme $5\neq 0$). $H \in (AB)$.
$G = \mathrm{bar}\{(H;5), (C;-1)\}$. Somme $= 4 \neq 0$.
$\overrightarrow{HG} = \frac{-1}{4}\overrightarrow{HC}$. $G$ est sur $(HC)$. Mais $H \in (AB)$ et $C$ n'est pas nécessairement sur $(AB)$.
\emph{Attente correction :} En fait, pour montrer $G \in (AB)$, il faut exprimer $G$ comme barycentre de points sur $(AB)$.
$G = \mathrm{bar}\{(A;2), (B;3), (C;-1)\}$.
On ne peut pas conclure directement $G \in (AB)$ sans hypothèse sur $C$.
\emph{Correction de l'énoncé :} « Soient $A,B,C$ alignés... » ou « Montrer que $G$ est sur $(AB)$ si $C \in (AB)$ ».
\emph{Supposons l'énoncé corrigé :} $A,B,C$ alignés. Alors $H \in (AB)$, $C \in (AB)$, donc $(HC) = (AB)$, donc $G \in (AB)$.
\item \textbf{Ligne de niveau :} $G = \mathrm{bar}\{(A;1),(B;3)\}$. Somme $=4$. $AG = \frac{3}{4}AB = 3$. $G(3;0)$.
$\|\overrightarrow{MA} + 3\overrightarrow{MB}\| = \|4\overrightarrow{MG}\| = 4 MG = 8 \Rightarrow MG = 2$.
Ensemble : Cercle centre $G(3;0)$ rayon $2$.
\item \textbf{Camion :} Masses égales $\Rightarrow$ coefficients $50$ et $30$. Somme $=80$.
$G = \mathrm{bar}\{(A;50), (B;30)\}$.
$\overrightarrow{AG} = \frac{30}{80}\overrightarrow{AB} = \frac{3}{8}\overrightarrow{AB}$.
$AG = \frac{3}{8} \times 4 = 1,5$ m depuis l'avant (point $A$).
Le point d'appui est à $1,50$ m de l'essieu avant (ou $2,50$ m de l'arrière).
\end{enumerate}
\end{corrige}

\begin{attention}[Les 5 erreurs qui coûtent des points au Probatoire]
\begin{enumerate}
\item \textbf{Oublier de vérifier que la somme des coefficients est non nulle.} Avant tout calcul ou toute construction, écrire explicitement : « $a+b=\dots\neq 0$, donc le barycentre existe ». Sans cette condition, la définition ne s'applique pas et la copie est sanctionnée.
\item \textbf{Inverser les coefficients dans la formule de construction.} C'est $\overrightarrow{AG}=\dfrac{b}{a+b}\overrightarrow{AB}$ : le coefficient de $B$ (l'\emph{autre} point) apparaît au numérateur. Vérifier mentalement : plus le coefficient de $B$ est grand, plus $G$ est proche de $B$.
\item \textbf{Confondre barycentre et isobarycentre.} L'isobarycentre exige des coefficients \emph{égaux}. Le centre de gravité d'un triangle est l'isobarycentre des sommets, pas un barycentre quelconque. Ne pas écrire « milieu de $[AB]$ » pour un barycentre de coefficients différents.
\item \textbf{Mal réduire la somme vectorielle dans les lignes de niveau.} La formule est $a\overrightarrow{MA}+b\overrightarrow{MB}=(a+b)\overrightarrow{MG}$, valable uniquement si $a+b\neq 0$. Erreur fréquente : oublier le facteur $(a+b)$ et écrire directement $MG=r$ au lieu de $MG=\dfrac{r}{|a+b|}$.
\item \textbf{Erreurs de calcul sur les coordonnées.} Oublier de diviser par la somme des coefficients, ou additionner des coordonnées avec des signes faux (surtout avec des coordonnées négatives). Poser proprement $x_G=\dfrac{a\,x_A+b\,x_B+c\,x_C}{a+b+c}$ en colonne, puis vérifier que le point obtenu est graphiquement plausible.
\end{enumerate}
\end{attention}

\newpage

\coursec{Exercices}

\courssub{Je vérifie mes connaissances}

\begin{exercice}[QCM : une seule réponse exacte]
Pour chaque question, choisir la bonne réponse en justifiant brièvement.
\begin{enumerate}
\item Soient $A$ et $B$ deux points distincts. Le barycentre $G$ de $(A,2)$ et $(B,6)$ vérifie :
\begin{enumerate}
\item[a)] $2\overrightarrow{GA}+6\overrightarrow{GB}=\vec{0}$ \quad b)] $6\overrightarrow{GA}+2\overrightarrow{GB}=\vec{0}$ \quad c)] $\overrightarrow{AG}=\dfrac{3}{4}\overrightarrow{AB}$
\end{enumerate}
\item Le barycentre de $(A,3)$ et $(B,-3)$ :
\begin{enumerate}
\item[a)] est le milieu de $[AB]$ \quad b)] n'existe pas \quad c)] est le point $A$
\end{enumerate}
\item Si $G=\mathrm{bar}\{(A,1),(B,1),(C,1)\}$, alors $G$ est :
\begin{enumerate}
\item[a)] le centre du cercle circonscrit au triangle $ABC$ \quad b)] l'orthocentre de $ABC$ \quad c)] le centre de gravité du triangle $ABC$
\end{enumerate}
\item L'ensemble des points $M$ du plan tels que $\|2\overrightarrow{MA}+2\overrightarrow{MB}\|=12$, avec $AB=8$~cm, est :
\begin{enumerate}
\item[a)] une droite \quad b)] un cercle de rayon $3$ \quad c)] un segment
\end{enumerate}
\end{enumerate}
\end{exercice}

\begin{exercice}[Vrai ou faux, en justifiant]
Répondre par VRAI ou FAUX en justifiant chaque réponse (réflexe exigé au Probatoire).
\begin{enumerate}
\item[a)] Le barycentre des points pondérés $(A,2)$ et $(B,-2)$ existe.
\item[b)] Si $G=\mathrm{bar}\{(A,1),(B,1),(C,1)\}$, alors $G$ est le centre du cercle circonscrit au triangle $ABC$.
\item[c)] L'ensemble des points $M$ tels que $MA^{2}-MB^{2}=0$ est la médiatrice du segment $[AB]$.
\item[d)] Le barycentre de $(A,1)$ et $(B,3)$ appartient au segment $[AB]$.
\end{enumerate}
\end{exercice}

\begin{exercice}[Définitions et rédaction]
\begin{enumerate}
\item Donner la définition du barycentre de deux points pondérés $(A,\alpha)$ et $(B,\beta)$, en précisant la condition d'existence.
\item Énoncer la propriété d'\cle{associativité} du barycentre pour trois points pondérés.
\item Compléter : si $G=\mathrm{bar}\{(A,\alpha),(B,\beta)\}$ avec $\alpha+\beta\neq 0$, alors pour tout point $M$ du plan :
\[ \alpha\overrightarrow{MA}+\beta\overrightarrow{MB} = (\alpha+\beta)\overrightarrow{M\ldots} \]
\end{enumerate}
\end{exercice}

\begin{exercice}[Calculs directs]
$A$ et $B$ sont deux points tels que $AB=6$~cm.
\begin{enumerate}
\item Construire le point $G=\mathrm{bar}\{(A,1),(B,2)\}$ après avoir exprimé $\overrightarrow{AG}$ en fonction de $\overrightarrow{AB}$.
\item Placer le point $H=\mathrm{bar}\{(A,3),(B,-1)\}$. Le point $H$ appartient-il au segment $[AB]$ ? Justifier.
\item Dans un repère, on donne $A(1\,;2)$ et $B(5\,;-4)$. Calculer les coordonnées de $G=\mathrm{bar}\{(A,2),(B,1)\}$.
\end{enumerate}
\end{exercice}

\courssub{Je m'entraîne}

\begin{exercice}[Barycentre de deux points et ligne de niveau]
$A$ et $B$ sont deux points tels que $AB=8$~cm.
\begin{enumerate}
\item Construire $G=\mathrm{bar}\{(A,3),(B,5)\}$.
\item Démontrer que pour tout point $M$ du plan : $3\overrightarrow{MA}+5\overrightarrow{MB}=8\overrightarrow{MG}$.
\item Déterminer et tracer l'ensemble $(\mathcal{C})$ des points $M$ tels que $\|3\overrightarrow{MA}+5\overrightarrow{MB}\|=16$.
\end{enumerate}
\end{exercice}

\begin{exercice}[Barycentre de trois points dans un repère]
Dans un repère orthonormé $(O\,;\vec{\imath},\vec{\jmath})$, on donne $A(-2\,;1)$, $B(4\,;3)$ et $C(0\,;5)$.
\begin{enumerate}
\item Vérifier que le point $G=\mathrm{bar}\{(A,1),(B,2),(C,-1)\}$ existe, puis calculer ses coordonnées.
\item Soit $\vec{u}\,(1\,;2)$. Démontrer que l'ensemble des points $M$ tels que $\overrightarrow{MA}+2\overrightarrow{MB}-\overrightarrow{MC}$ soit colinéaire à $\vec{u}$ est une droite dont on donnera un point et un vecteur directeur.
\end{enumerate}
\end{exercice}

\begin{exercice}[Associativité du barycentre]
Soit $ABC$ un triangle. On pose $I=\mathrm{bar}\{(B,2),(C,1)\}$ et $J=\mathrm{bar}\{(A,1),(C,2)\}$.
\begin{enumerate}
\item Construire les points $I$ et $J$ sur une figure.
\item En utilisant l'associativité du barycentre avec le système $\{(A,1),(B,2),(C,2)\}$, démontrer que les droites $(AI)$ et $(BJ)$ se coupent en un point $K$ que l'on exprimera comme barycentre de $A$, $B$ et $C$.
\item En déduire que la droite $(CK)$ coupe le segment $[AB]$ en son milieu.
\end{enumerate}
\end{exercice}

\begin{exercice}[Ligne de niveau $MA^{2}+MB^{2}$]
$A$ et $B$ sont deux points tels que $AB=8$~cm, et $I$ est le milieu de $[AB]$.
\begin{enumerate}
\item Démontrer que pour tout point $M$ : $MA^{2}+MB^{2}=2MI^{2}+\dfrac{AB^{2}}{2}$.
\item Déterminer la nature de l'ensemble des points $M$ tels que $MA^{2}+MB^{2}=50$. Le tracer.
\item Pour quelles valeurs du réel $k$ l'ensemble des points $M$ tels que $MA^{2}+MB^{2}=k$ est-il non vide ?
\end{enumerate}
\end{exercice}

\courssub{Je me prépare au Probatoire}

\begin{exercice}[Type Probatoire : barycentre et lignes de niveau]
$A$ et $B$ sont deux points du plan tels que $AB=8$~cm.

\textbf{Partie A.}
\begin{enumerate}
\item Justifier l'existence du point $G=\mathrm{bar}\{(A,3),(B,5)\}$ et le construire en laissant les traits de construction apparents.
\item Calculer les distances $GA$ et $GB$.
\end{enumerate}

\textbf{Partie B.}
\begin{enumerate}
\item Démontrer que pour tout point $M$ du plan, $3\overrightarrow{MA}+5\overrightarrow{MB}=8\overrightarrow{MG}$.
\item Déterminer et tracer l'ensemble $(\mathcal{C}_{1})$ des points $M$ tels que $\|3\overrightarrow{MA}+5\overrightarrow{MB}\|=16$.
\end{enumerate}

\textbf{Partie C.}
\begin{enumerate}
\item Démontrer que pour tout point $M$ :
\[ 3MA^{2}+5MB^{2}=8MG^{2}+3GA^{2}+5GB^{2}. \]
\item En déduire l'ensemble $(\mathcal{C}_{2})$ des points $M$ tels que $3MA^{2}+5MB^{2}=200$ : on précisera sa nature, son centre et son rayon.
\item Tracer $(\mathcal{C}_{2})$ sur la même figure que $(\mathcal{C}_{1})$.
\end{enumerate}
\end{exercice}

\begin{exercice}[Type Probatoire : concours de droites dans un parallélogramme]
Soit $ABCD$ un parallélogramme. On note $I$ le milieu de $[BC]$, $J$ le milieu de $[CD]$, $K$ le milieu de $[DA]$ et $L$ le milieu de $[AB]$.
\begin{enumerate}
\item Faire une figure soignée.
\item Justifier que $A=\mathrm{bar}\{(A,1)\}$ et exprimer $I$ comme barycentre de $B$ et $C$ affectés de coefficients égaux.
\item On considère le point $\Omega=\mathrm{bar}\{(A,1),(B,1),(C,1),(D,1)\}$.
\begin{enumerate}
\item En utilisant l'associativité du barycentre, démontrer que $\Omega$ appartient à la droite $(AI)$ et à la droite $(BJ)$.
\item Démontrer de même que $\Omega$ appartient aux droites $(CK)$ et $(DL)$.
\end{enumerate}
\item Conclure que les droites $(AI)$, $(BJ)$, $(CK)$ et $(DL)$ sont concourantes en $\Omega$.
\item Préciser la position de $\Omega$ sur le segment $[AI]$ (exprimer $\overrightarrow{A\Omega}$ en fonction de $\overrightarrow{AI}$). Que représente $\Omega$ pour le parallélogramme $ABCD$ ?
\end{enumerate}
\end{exercice}

\begin{exercice}[Type Probatoire : équilibre d'une tige chargée]
Un atelier de mécanique utilise une tige homogène $[AB]$ de longueur $3$~m, de masse négligeable. On fixe une charge de $12$~kg en $A$ et une charge de $8$~kg en $B$. On modélise chaque charge par un point pondéré dont le coefficient est la masse.

\textbf{Partie A.}
\begin{enumerate}
\item Le point d'équilibre de la tige est le barycentre $G$ des points pondérés $(A,12)$ et $(B,8)$. Justifier l'existence de $G$.
\item Exprimer $\overrightarrow{AG}$ en fonction de $\overrightarrow{AB}$, puis calculer la distance $AG$. À quelle distance de $A$ faut-il placer le point d'appui pour que la tige soit en équilibre ?
\end{enumerate}

\textbf{Partie B.} On ajoute une charge de $5$~kg au milieu $I$ du segment $[AB]$.
\begin{enumerate}
\item En utilisant l'associativité du barycentre, montrer que le nouveau point d'équilibre $G'$ est le barycentre de $(G,20)$ et $(I,5)$.
\item Calculer la distance $GG'$, puis la distance $AG'$.
\item Le point d'appui s'est-il déplacé vers $A$ ou vers $B$ ? Ce résultat était-il prévisible ? Justifier la réponse par une phrase rédigée.
\end{enumerate}
\end{exercice}

\newpage

\coursec{Corrigés des exercices}

\coursec{Corrigés détaillés — Barycentres (Première Technique)}

\courssub{Exercice 1 — QCM : analyse des propositions}
\begin{corrige}[Exercice 1]
\begin{enumerate}
\item \textbf{Réponses a) et c) (les deux sont mathématiquement exactes).} 
Par définition, le barycentre $G$ de $(A,2)$ et $(B,6)$ vérifie $2\overrightarrow{GA}+6\overrightarrow{GB}=\vec{0}$ (proposition a). 
De plus, par la relation de Chasles : $\overrightarrow{AG} = \dfrac{6}{2+6}\overrightarrow{AB} = \dfrac{3}{4}\overrightarrow{AB}$ (proposition c). 
\textbf{Attention :} l'énoncé précise « une seule réponse exacte », or les deux propositions a) et c) sont correctes. La réponse a) est la définition même, la c) en est une conséquence immédiate. En examen, signaler l'ambiguïté au correcteur.
\item \textbf{Réponse b).} La somme des coefficients est $3+(-3)=0$. Le barycentre n'existe pas (il serait un point à l'infini, direction $\overrightarrow{AB}$).
\item \textbf{Réponse c).} Le barycentre de trois points de même coefficient ($1,1,1$) est le centre de gravité (centroïde) du triangle $ABC$, intersection des médianes.
\item \textbf{Réponse b).} Soit $I$ le milieu de $[AB]$. Pour tout $M$, $2\overrightarrow{MA}+2\overrightarrow{MB}=4\overrightarrow{MI}$. La condition $\|2\overrightarrow{MA}+2\overrightarrow{MB}\|=12$ devient $4\,MI=12$, soit $MI=3$. L'ensemble est le cercle de centre $I$ et de rayon $3$~cm. ($AB=8$ n'intervient pas dans le rayon).
\end{enumerate}
\end{corrige}

\courssub{Exercice 2 — Vrai ou faux, avec justification}
\begin{corrige}[Exercice 2]
\begin{enumerate}
\item[a)] \textbf{FAUX.} La somme des coefficients est $2+(-2)=0$. Le barycentre de $(A,2)$ et $(B,-2)$ n'existe pas (sauf si $A=B$, ce qui est exclu pour des points distincts).
\item[b)] \textbf{FAUX.} $G=\mathrm{bar}\{(A,1),(B,1),(C,1)\}$ est le centre de gravité (centroïde) du triangle $ABC$, intersection des médianes. Le centre du cercle circonscrit est l'intersection des médiatrices, distinct en général.
\item[c)] \textbf{VRAI.} $MA^2-MB^2=0 \Leftrightarrow MA=MB$. L'ensemble des points équidistants de $A$ et $B$ est la médiatrice du segment $[AB]$.
\item[d)] \textbf{VRAI.} Les coefficients $1$ et $3$ sont de même signe (strictement positifs). Le barycentre $G$ appartient au segment $[AB]$ et le divise dans le rapport $GB/GA = 1/3$ (donc $AG = \frac{3}{4}AB$).
\end{enumerate}
\end{corrige}

\courssub{Exercice 3 — Définitions, propriétés et relation fondamentale}
\begin{corrige}[Exercice 3]
\begin{enumerate}
\item Soient $A$ et $B$ deux points distincts du plan et $\alpha,\beta$ deux réels tels que $\alpha+\beta\neq 0$. Le \cle{barycentre} des points pondérés $(A,\alpha)$ et $(B,\beta)$ est l'unique point $G$ du plan tel que :
\[ \alpha\overrightarrow{GA}+\beta\overrightarrow{GB}=\vec{0}. \]
Condition d'existence : $\alpha+\beta\neq 0$ (si $\alpha+\beta=0$ et $A\neq B$, il n'existe pas de point fini vérifiant la relation).
\item \textbf{Propriété d'\cle{associativité} (trois points).} Soient $(A,\alpha)$, $(B,\beta)$, $(C,\gamma)$ trois points pondérés avec $\alpha+\beta+\gamma\neq 0$. Si l'on regroupe deux points, par exemple $A$ et $B$, en leur barycentre partiel $K=\mathrm{bar}\{(A,\alpha),(B,\beta)\}$ (coefficient $\alpha+\beta$), alors le barycentre $G$ du système total est le barycentre de $(K,\alpha+\beta)$ et $(C,\gamma)$ :
\[ G = \mathrm{bar}\{(A,\alpha),(B,\beta),(C,\gamma)\} = \mathrm{bar}\{(K,\alpha+\beta),(C,\gamma)\}. \]
\item Relation fondamentale : $\alpha\overrightarrow{MA}+\beta\overrightarrow{MB} = (\alpha+\beta)\overrightarrow{MG}$.
\end{enumerate}
\end{corrige}

\courssub{Exercice 4 — Calculs directs et construction}
\begin{corrige}[Exercice 4]
$AB=6$~cm.
\begin{enumerate}
\item $G=\mathrm{bar}\{(A,1),(B,2)\}$. $\overrightarrow{AG} = \dfrac{2}{1+2}\overrightarrow{AB} = \dfrac{2}{3}\overrightarrow{AB}$.
$AG = \frac{2}{3}\times 6 = 4$~cm. 
\textbf{Construction :} Sur la droite $(AB)$, reporter 3 fois une même unité à partir de $A$ vers $B$ ; le $2^{\text{ème}}$ repère donne $G$ (ou diviser $[AB]$ en 3 parties égales, $G$ est à 2 parties de $A$).
\item $H=\mathrm{bar}\{(A,3),(B,-1)\}$. Somme des coefficients $= 2 \neq 0$, $H$ existe.
$\overrightarrow{AH} = \dfrac{-1}{3+(-1)}\overrightarrow{AB} = -\dfrac{1}{2}\overrightarrow{AB}$.
$AH = \frac{1}{2}AB = 3$~cm. Le vecteur $\overrightarrow{AH}$ est opposé à $\overrightarrow{AB}$, donc $H$ est sur la droite $(AB)$ mais \textbf{en dehors du segment $[AB]$}, du côté de $A$ (puisque $\overrightarrow{AH}$ pointe vers l'opposé de $B$).
\item $A(1;2)$, $B(5;-4)$, $G=\mathrm{bar}\{(A,2),(B,1)\}$. Somme $=3$.
\[ x_G = \frac{2\times 1 + 1\times 5}{3} = \frac{7}{3}, \quad y_G = \frac{2\times 2 + 1\times (-4)}{3} = 0. \]
$G\left(\dfrac{7}{3}\,;\,0\right)$.
\end{enumerate}
\end{corrige}

\courssub{Exercice 5 — Barycentre de deux points et ligne de niveau (cercle)}
\begin{corrige}[Exercice 5]
$AB=8$~cm.
\begin{enumerate}
\item $G=\mathrm{bar}\{(A,3),(B,5)\}$. $\overrightarrow{AG} = \frac{5}{8}\overrightarrow{AB}$. $AG = \frac{5}{8}\times 8 = 5$~cm, $GB = 3$~cm.
Construction : diviser $[AB]$ en 8 parties égales, $G$ est à 5 parties de $A$ (ou 3 de $B$).
\item Pour tout $M$ : 
\[ 3\overrightarrow{MA}+5\overrightarrow{MB} = 3(\overrightarrow{MG}+\overrightarrow{GA})+5(\overrightarrow{MG}+\overrightarrow{GB}) = 8\overrightarrow{MG} + (3\overrightarrow{GA}+5\overrightarrow{GB}). \]
Par définition de $G$, $3\overrightarrow{GA}+5\overrightarrow{GB}=\vec{0}$. Donc $3\overrightarrow{MA}+5\overrightarrow{MB}=8\overrightarrow{MG}$.
\item $\|3\overrightarrow{MA}+5\overrightarrow{MB}\|=16 \Leftrightarrow \|8\overrightarrow{MG}\|=16 \Leftrightarrow 8\,MG=16 \Leftrightarrow MG=2$.
L'ensemble $(\mathcal{C})$ est le \textbf{cercle de centre $G$ et de rayon $2$~cm}.
\end{enumerate}
\begin{popfigure}
\begin{tikzpicture}[scale=0.8,>=Stealth]
\coordinate (A) at (0,0);
\coordinate (B) at (8,0);
\coordinate (G) at (5,0);
\draw[popline, thick] (A) -- (B);
\foreach \x/\lab in {0/A, 5/G, 8/B} {
  \fill[popblue] (\x,0) circle (2pt);
  \node[below] at (\x,0) {$\lab$};
}
\draw[poporange, thick] (G) circle (2);
\node[poporange, above right] at (7,2) {$(\mathcal{C})\,:\, MG=2$};
\draw[|<->|, popmuted] (0,-1) -- (5,-1) node[midway, below] {5 cm};
\draw[|<->|, popmuted] (5,-1) -- (8,-1) node[midway, below] {3 cm};
\end{tikzpicture}
\legende{Construction de $G$ et cercle $(\mathcal{C})$}
\end{popfigure}
\end{corrige}

\courssub{Exercice 6 — Barycentre de trois points dans un repère orthonormé}
\begin{corrige}[Exercice 6]
$A(-2;1)$, $B(4;3)$, $C(0;5)$.
\begin{enumerate}
\item $G=\mathrm{bar}\{(A,1),(B,2),(C,-1)\}$. Somme des coefficients $\sigma = 1+2+(-1)=2 \neq 0$ : $G$ existe.
Coordonnées de $G$ (formule barycentrique) :
\[ x_G = \frac{1\times(-2) + 2\times 4 + (-1)\times 0}{2} = \frac{-2+8}{2} = 3, \]
\[ y_G = \frac{1\times 1 + 2\times 3 + (-1)\times 5}{2} = \frac{1+6-5}{2} = 1. \]
$G(3;1)$.
\item Soit $M(x;y)$. $\overrightarrow{MA}+2\overrightarrow{MB}-\overrightarrow{MC} = (A-M)+2(B-M)-(C-M) = (A+2B-C) - 2M$.
D'après la question 1, $A+2B-C = 2G = (6;2)$.
Le vecteur cherché est $(6;2) - 2(x;y) = 2(3-x; 1-y) = 2\overrightarrow{MG}$.
Il est colinéaire à $\vec{u}(1;2)$ $\Leftrightarrow$ $\overrightarrow{MG}$ est colinéaire à $\vec{u}$.
L'ensemble des points $M$ est la \textbf{droite passant par $G(3;1)$ de vecteur directeur $\vec{u}(1;2)$}.
Équation cartésienne : $\frac{x-3}{1} = \frac{y-1}{2} \Leftrightarrow 2x - y - 5 = 0$.
\end{enumerate}
\end{corrige}

\courssub{Exercice 7 — Associativité du barycentre et concourrance}
\begin{corrige}[Exercice 7]
\begin{attention}[Incohérence dans l'énoncé original]
Les points $I=\mathrm{bar}\{(B,2),(C,1)\}$ et $J=\mathrm{bar}\{(A,1),(C,2)\}$ ne proviennent \textbf{pas} d'un même système de trois points pondérés $(A,\alpha),(B,\beta),(C,\gamma)$ (le coefficient de $C$ serait à la fois $1$ et $2$). Par conséquent, les droites $(AI)$ et $(BJ)$ ne sont pas nécessairement concourantes en un barycentre $K$ d'un système unique. La suite de l'exercice suppose le système $\{(A,1),(B,2),(C,2)\}$ donné explicitement. Nous corrigeons ci-dessous en suivant ce système imposé.
\end{attention}

\begin{enumerate}
\item \textbf{Figure.} Voir ci-dessous. $I$ divise $[BC]$ tel que $BI = \frac{1}{3}BC$ (coeff 2 pour B, 1 pour C). $J$ divise $[AC]$ tel que $AJ = \frac{2}{3}AC$ (coeff 1 pour A, 2 pour C). $K$ est le barycentre de $\{(A,1),(B,2),(C,2)\}$.
\begin{popfigure}
\begin{tikzpicture}[scale=1,>=Stealth]
\coordinate (A) at (0,0);
\coordinate (B) at (6,0);
\coordinate (C) at (2,5);
% I = bar{(B,2),(C,1)} -> 1/3 from B to C
\coordinate (I) at ($(B)!0.3333!(C)$); 
% J = bar{(A,1),(C,2)} -> 2/3 from A to C
\coordinate (J) at ($(A)!0.6666!(C)$); 
% K = bar{(A,1),(B,2),(C,2)} = (A+2B+2C)/5
\coordinate (K) at ($ 0.2*(A) + 0.4*(B) + 0.4*(C) $);
\draw[popline, thick] (A)--(B)--(C)--cycle;
\draw[popblue, thick] (A)--(I);
\draw[poporange, thick] (B)--(J);
\draw[popgreen, thick] (C)--(K);
\fill[popblue] (A) circle (2pt) node[below left] {$A$};
\fill[popblue] (B) circle (2pt) node[below right] {$B$};
\fill[popblue] (C) circle (2pt) node[above] {$C$};
\fill[popblue] (I) circle (2pt) node[right] {$I$};
\fill[popblue] (J) circle (2pt) node[above left] {$J$};
\fill[popred] (K) circle (3pt) node[above right] {$K$};
\end{tikzpicture}
\legende{Figure : $I$, $J$ et $K$ (barycentre de $\{(A,1),(B,2),(C,2)\}$)}
\end{popfigure}

\item On considère le système $\mathcal{S} = \{(A,1), (B,2), (C,2)\}$ (somme $=5$).
\begin{itemize}
\item \textbf{Test pour $(AI)$ :} Pour utiliser $I$, on isole le coefficient de $C$ dans $I$. $I = \mathrm{bar}\{(B,2),(C,1)\}$ (coeff 3). Le système $\mathcal{S}$ s'écrit $\{(A,1), (I,3), (C,1)\}$ (car $C$ a coeff 2 total, 1 déjà dans $I$, il en reste 1). $K$ n'est pas barycentre de $\{(A,1),(I,3)\}$ seul (reste $(C,1)$). \textbf{$K$ n'appartient pas à la droite $(AI)$.}
\item \textbf{Test pour $(BJ)$ :} $J = \mathrm{bar}\{(A,1),(C,2)\}$ (coeff 3). Le système $\mathcal{S}$ s'écrit exactement $\{(J,3), (B,2)\}$. Par associativité, $K = \mathrm{bar}\{(J,3),(B,2)\}$. \textbf{$K$ appartient à la droite $(BJ)$.}
\end{itemize}
\textbf{Conclusion :} Avec le système $\{(A,1),(B,2),(C,2)\}$, $K$ est sur $(BJ)$ mais \textbf{pas sur $(AI)$}. Si l'énoncé voulait la concourrance des trois droites $(AI), (BJ), (CK)$, il fallait définir $I=\mathrm{bar}\{(B,2),(C,2)\}$ (milieu de $[BC]$).

\item \textbf{Intersection de $(CK)$ avec $(AB)$ :} Avec $K=\mathrm{bar}\{(A,1),(B,2),(C,2)\}$, la droite $(CK)$ coupe $(AB)$ au point $M = \mathrm{bar}\{(A,1),(B,2)\}$ (coefficient $3$). $\overrightarrow{AM} = \frac{2}{3}\overrightarrow{AB}$. $M$ n'est \textbf{pas} le milieu de $[AB]$ (le milieu correspondrait à des coefficients égaux pour $A$ et $B$).
\end{enumerate}
\end{corrige}

\courssub{Exercice 8 — Ligne de niveau $MA^{2}+MB^{2}$ (Relation d'Apollonius)}
\begin{corrige}[Exercice 8]
$AB=8$~cm, $I$ milieu de $[AB]$ ($AI=IB=4$~cm).
\begin{enumerate}
\item \textbf{Démonstration (relation d'Apollonius / médiane).} Pour tout $M$ :
\[ \overrightarrow{MA} = \overrightarrow{MI}+\overrightarrow{IA}, \quad \overrightarrow{MB} = \overrightarrow{MI}+\overrightarrow{IB} = \overrightarrow{MI}-\overrightarrow{IA}. \]
\[ MA^2 = MI^2 + IA^2 + 2\overrightarrow{MI}\cdot\overrightarrow{IA}, \quad MB^2 = MI^2 + IA^2 - 2\overrightarrow{MI}\cdot\overrightarrow{IA}. \]
En additionnant : $MA^2+MB^2 = 2MI^2 + 2IA^2 = 2MI^2 + 2\left(\frac{AB}{2}\right)^2 = 2MI^2 + \frac{AB^2}{2}$.
\item $MA^2+MB^2=50 \Rightarrow 2MI^2 + \frac{64}{2} = 50 \Rightarrow 2MI^2 + 32 = 50 \Rightarrow 2MI^2 = 18 \Rightarrow MI^2 = 9 \Rightarrow MI = 3$.
L'ensemble est le \textbf{cercle de centre $I$ et de rayon $3$~cm}.
\begin{popfigure}
\begin{tikzpicture}[scale=0.7,>=Stealth]
\coordinate (A) at (0,0);
\coordinate (B) at (8,0);
\coordinate (I) at (4,0);
\draw[popline, thick] (A)--(B);
\fill[popblue] (A) circle (2pt) node[below] {$A$};
\fill[popblue] (B) circle (2pt) node[below] {$B$};
\fill[poporange] (I) circle (2pt) node[below] {$I$};
\draw[poporange, thick] (I) circle (3);
\node[poporange, above right] at (7,3) {$MI=3$};
\end{tikzpicture}
\legende{Cercle centre $I$ rayon 3}
\end{popfigure}
\item L'ensemble est non vide si le rayon est réel : $MI^2 = \frac{k - AB^2/2}{2} \geq 0 \Rightarrow k \geq \frac{AB^2}{2} = 32$.
Pour $k=32$, l'ensemble se réduit au point $I$ ; pour $k>32$, c'est un cercle de centre $I$ et de rayon $\sqrt{\frac{k-32}{2}}$.
\end{enumerate}
\end{corrige}

\courssub{Exercice 9 — Type Probatoire : barycentre et lignes de niveau}
\begin{corrige}[Exercice 9]
\textbf{Barème indicatif :} Partie A (3 pts), Partie B (4 pts), Partie C (5 pts). Total ~12 pts.
\begin{attention}[Points de vigilance Probatoire]
\begin{itemize}
\item Justifier l'existence du barycentre (somme des coefficients $\neq 0$).
\item Laisser les traits de construction apparents (règle, compas).
\item Citer la relation fondamentale $3\overrightarrow{MA}+5\overrightarrow{MB}=8\overrightarrow{MG}$ (démonstration attendue).
\item Pour les lignes de niveau : passer par la norme du vecteur, identifier le cercle (centre, rayon).
\item Formule de Leibniz : $\sum \alpha_i MA_i^2 = (\sum \alpha_i) MG^2 + \sum \alpha_i GA_i^2$. La démontrer ou la citer explicitement.
\end{itemize}
\end{attention}

\textbf{Partie A.}
\begin{enumerate}
\item $G=\mathrm{bar}\{(A,3),(B,5)\}$. Somme des coefficients $= 3+5 = 8 \neq 0$ : $G$ existe et est unique.
\textbf{Construction :} Sur la droite $(AB)$, construire un segment $[AB]=8$~cm. Diviser $[AB]$ en 8 parties égales (ou utiliser un rayon gradué). Le point $G$ est le $5^{\text{ème}}$ point de division à partir de $A$ (ou $3^{\text{ème}}$ à partir de $B$). $AG=5$~cm, $GB=3$~cm. (Traits de construction : parallèles équidistantes ou report de compas).
\item D'après la définition : $3\overrightarrow{GA}+5\overrightarrow{GB}=\vec{0} \Rightarrow 3\overrightarrow{GA} = -5\overrightarrow{GB} \Rightarrow 3\,GA = 5\,GB$.
De plus $GA+GB=AB=8$. Donc $GA = \frac{5}{8}\times 8 = 5$~cm, $GB = 3$~cm.
\end{enumerate}

\textbf{Partie B.}
\begin{enumerate}
\item Pour tout $M$ : $3\overrightarrow{MA}+5\overrightarrow{MB} = 3(\overrightarrow{MG}+\overrightarrow{GA})+5(\overrightarrow{MG}+\overrightarrow{GB}) = 8\overrightarrow{MG} + (3\overrightarrow{GA}+5\overrightarrow{GB}) = 8\overrightarrow{MG}$.
\item $\|3\overrightarrow{MA}+5\overrightarrow{MB}\|=16 \Leftrightarrow \|8\overrightarrow{MG}\|=16 \Leftrightarrow 8\,MG=16 \Leftrightarrow MG=2$.
$(\mathcal{C}_1)$ est le \textbf{cercle de centre $G$ et de rayon $2$~cm}.
\end{enumerate}

\textbf{Partie C.}
\begin{enumerate}
\item \textbf{Formule de Leibniz (ou développement) :}
$3MA^2+5MB^2 = 3(\overrightarrow{MG}+\overrightarrow{GA})^2 + 5(\overrightarrow{MG}+\overrightarrow{GB})^2$
$= 3(MG^2+GA^2+2\overrightarrow{MG}\cdot\overrightarrow{GA}) + 5(MG^2+GB^2+2\overrightarrow{MG}\cdot\overrightarrow{GB})$
$= 8MG^2 + 3GA^2+5GB^2 + 2\overrightarrow{MG}\cdot(3\overrightarrow{GA}+5\overrightarrow{GB})$.
Le dernier terme est nul car $3\overrightarrow{GA}+5\overrightarrow{GB}=\vec{0}$.
Donc $3MA^2+5MB^2 = 8MG^2 + 3GA^2+5GB^2$.
\item Calcul des constantes : $GA=5$, $GB=3$. $3GA^2+5GB^2 = 3\times 25 + 5\times 9 = 75+45 = 120$.
Équation : $8MG^2 + 120 = 200 \Rightarrow 8MG^2 = 80 \Rightarrow MG^2 = 10 \Rightarrow MG = \sqrt{10}$.
$(\mathcal{C}_2)$ est le \textbf{cercle de centre $G$ et de rayon $\sqrt{10} \approx 3,16$~cm}.
\item \textbf{Figure commune :} Deux cercles concentriques de centre $G$ : $(\mathcal{C}_1)$ rayon $2$~cm, $(\mathcal{C}_2)$ rayon $\sqrt{10}$~cm.
\begin{popfigure}
\begin{tikzpicture}[scale=0.7,>=Stealth]
\coordinate (A) at (0,0);
\coordinate (B) at (8,0);
\coordinate (G) at (5,0);
\draw[popline, thick] (A)--(B);
\foreach \x/\lab in {0/A, 5/G, 8/B} {
  \fill[popblue] (\x,0) circle (2pt);
  \node[below] at (\x,0) {$\lab$};
}
\draw[poporange, thick] (G) circle (2);
\draw[poppurple, thick] (G) circle (3.162);
\node[poporange, above] at (7,2) {$(\mathcal{C}_1)\,:\, R=2$};
\node[poppurple, above] at (8.2,3.16) {$(\mathcal{C}_2)\,:\, R=\sqrt{10}$};
\end{tikzpicture}
\legende{Cercles concentriques $(\mathcal{C}_1)$ et $(\mathcal{C}_2)$ de centre $G$}
\end{popfigure}
\end{enumerate}
\end{corrige}

\courssub{Exercice 10 — Type Probatoire : concours de droites dans un parallélogramme}
\begin{corrige}[Exercice 10]
\begin{attention}[Erreur mathématique dans l'énoncé]
Dans un parallélogramme $ABCD$, les droites joignant un sommet au milieu du côté opposé $(AI), (BJ), (CK), (DL)$ ne sont \textbf{pas concourantes} en général (contrairement aux médianes d'un triangle). Le barycentre $\Omega$ des 4 sommets (coefficients 1,1,1,1) est le centre du parallélogramme (intersection des diagonales), mais il n'appartient pas à $(AI)$ ni $(BJ)$. La propriété de concourrance est fausse. Voir la correction détaillée ci-dessous.
\end{attention}

\begin{enumerate}
\item \textbf{Figure.} Parallélogramme $ABCD$, milieux $I,J,K,L$, centre $\Omega$ (intersection diagonales).
\begin{popfigure}
\begin{tikzpicture}[scale=0.9,>=Stealth]
\coordinate (A) at (0,0);
\coordinate (B) at (5,0);
\coordinate (D) at (1.5,3);
\coordinate (C) at ($(B)+(D)-(A)$); % 6.5, 3
\coordinate (I) at ($(B)!0.5!(C)$);
\coordinate (J) at ($(C)!0.5!(D)$);
\coordinate (K) at ($(D)!0.5!(A)$);
\coordinate (L) at ($(A)!0.5!(B)$);
\coordinate (O) at ($(A)!0.5!(C)$); % Centre
\draw[popline, thick] (A)--(B)--(C)--(D)--cycle;
\draw[popline, dashed] (A)--(C) (B)--(D);
\draw[popblue, thick] (A)--(I) (B)--(J) (C)--(K) (D)--(L);
\foreach \p/\lab/\pos in {A/A/below left, B/B/below right, C/C/above right, D/D/above left, I/I/right, J/J/above, K/K/left, L/L/below, O/$\Omega$/below right} {
  \fill[popblue] (\p) circle (2pt);
  \node[\pos] at (\p) {$\lab$};
}
\end{tikzpicture}
\legende{Parallélogramme : droites $(AI), (BJ), (CK), (DL)$ non concourantes. $\Omega$ = intersection diagonales.}
\end{popfigure}

\item $A = \mathrm{bar}\{(A,1)\}$ (trivial). $I$ milieu de $[BC] \Rightarrow I = \mathrm{bar}\{(B,1),(C,1)\}$.
\item $\Omega = \mathrm{bar}\{(A,1),(B,1),(C,1),(D,1)\}$ (somme $=4$).
\begin{enumerate}
\item Tentative pour $(AI)$ : Regrouper $(B,1),(C,1) \to I$ (coeff 2). Système : $\{(A,1), (I,2), (D,1)\}$. $\Omega$ n'est pas barycentre de $\{(A,1),(I,2)\}$ seul (reste $(D,1)$). \textbf{$\Omega \notin (AI)$ en général.}
\item Tentative pour $(BJ)$ : Regrouper $(C,1),(D,1) \to J$ (coeff 2). Système : $\{(A,1), (B,1), (J,2)\}$. $\Omega$ n'est pas sur $(BJ)$.
\end{enumerate}
\textbf{Correct :} Regrouper $(A,1),(C,1) \to \Omega$ (coeff 2) et $(B,1),(D,1) \to \Omega$ (coeff 2). $\Omega$ est le milieu de $[AC]$ et $[BD]$ (centre du parallélogramme). Les droites concourantes en $\Omega$ sont les diagonales $(AC), (BD)$ et les bimédianes $(IK), (JL)$ (segments joignant milieux côtés opposés).

\item \textbf{Conclusion corrigée :} Les droites $(AI), (BJ), (CK), (DL)$ ne sont pas concourantes. Le point $\Omega$ (centre du parallélogramme) est l'intersection des diagonales $(AC)$ et $(BD)$ et des bimédianes $(IK)$ et $(JL)$.

\item $\Omega$ n'est pas sur $[AI]$. Si l'on considère le système $\{(A,1),(B,1),(C,1),(D,1)\}$, $\Omega$ est le centre de symétrie du parallélogramme (intersection des diagonales).
\end{enumerate}
\end{corrige}

\courssub{Exercice 11 — Type Probatoire : équilibre d'une tige chargée (application technique)}
\begin{corrige}[Exercice 11]
\textbf{Barème indicatif :} Partie A (4 pts), Partie B (6 pts). Total ~10 pts.
\begin{attention}[Points de vigilance]
\begin{itemize}
\item Modélisation : charges $\to$ points pondérés (coefficients = masses en kg).
\item Existence : somme des masses $\neq 0$.
\item Associativité : remplacer le sous-système $(A,12),(B,8)$ par son barycentre $G$ affecté du coefficient $20$.
\item Calculs vectoriels précis : $\overrightarrow{AG'} = \overrightarrow{AG} + \overrightarrow{GG'}$.
\item Interprétation physique : le centre de masse se déplace vers la masse ajoutée.
\end{itemize}
\end{attention}

\textbf{Partie A.}
\begin{enumerate}
\item $G = \mathrm{bar}\{(A,12),(B,8)\}$. Somme $= 20 \neq 0$ : $G$ existe unique. C'est le centre de masse du système, donc le point d'équilibre.
\item $\overrightarrow{AG} = \frac{8}{12+8}\overrightarrow{AB} = \frac{8}{20}\overrightarrow{AB} = \frac{2}{5}\overrightarrow{AB}$.
$AG = \frac{2}{5} \times 3 = 1,2$~m.
Il faut placer l'appui à \textbf{1,20~m de $A$} (soit 1,80~m de $B$).
\end{enumerate}

\textbf{Partie B.}
$I$ milieu de $[AB] \Rightarrow AI = IB = 1,5$~m. Charge $5$~kg en $I$.
\begin{enumerate}
\item Système initial : $(A,12), (B,8) \to$ barycentre $G$ (coeff $20$).
Nouveau système total : $(A,12), (B,8), (I,5)$.
Par associativité, on peut remplacer $\{(A,12),(B,8)\}$ par $(G,20)$.
Le nouveau barycentre $G'$ est donc $G' = \mathrm{bar}\{(G,20), (I,5)\}$ (somme $=25$).
\item $\overrightarrow{GG'} = \frac{5}{20+5}\overrightarrow{GI} = \frac{1}{5}\overrightarrow{GI}$.
Positions sur la droite (origine $A$) : $A=0$, $G=1,2$, $I=1,5$, $B=3$.
$\overrightarrow{GI} = \overrightarrow{GA} + \overrightarrow{AI} = -1,2 + 1,5 = 0,3$~m (vecteur orienté vers $B$).
$GG' = \frac{1}{5} \times 0,3 = 0,06$~m.
$AG' = AG + GG' = 1,2 + 0,06 = 1,26$~m.
\item Le point d'appui se déplace de $G$ vers $G'$, soit \textbf{vers $B$} (puisque $GG' > 0$ dans le sens $\overrightarrow{AB}$).
\textbf{Justification :} On a ajouté une masse ($5$~kg) en $I$, qui se situe entre $G$ et $B$. Le centre de masse global se déplace nécessairement vers la masse ajoutée, donc vers $B$. Ce résultat était prévisible.
\end{enumerate}
\end{corrige}

\newpage

\coursec{Fiche bilan}

\begin{popfigure}
\begin{center}\tcbox[colback=popblue,colframe=popblue,coltext=white,arc=9pt,boxsep=4pt,left=6pt,right=6pt]{\bfseries\small Barycentres}\end{center}
\vspace{-2pt}
\begin{tcbraster}[raster columns=3,raster equal height=rows,raster column skip=6pt,raster row skip=6pt,enhanced,blankest]
\begin{tcolorbox}[enhanced,colback=white,colframe=poporange,boxrule=0.9pt,arc=3pt,title={2 points pondérés},fonttitle=\bfseries\scriptsize,coltitle=white,colbacktitle=poporange,left=4pt,right=4pt,top=2pt,bottom=2pt,boxsep=1pt,toptitle=1.5pt,bottomtitle=1.5pt,halign=left]
\begin{itemize}[nosep,leftmargin=*,itemsep=1pt,font=\scriptsize]
\item $\alpha\overrightarrow{GA}+\beta\overrightarrow{GB}=\vec{0}$
\item Existence si $\alpha+\beta\neq 0$
\item $\overrightarrow{AG}=\frac{\beta}{\alpha+\beta}\overrightarrow{AB}$
\end{itemize}
\end{tcolorbox}
\begin{tcolorbox}[enhanced,colback=white,colframe=popgreen,boxrule=0.9pt,arc=3pt,title={Propriétés},fonttitle=\bfseries\scriptsize,coltitle=white,colbacktitle=popgreen,left=4pt,right=4pt,top=2pt,bottom=2pt,boxsep=1pt,toptitle=1.5pt,bottomtitle=1.5pt,halign=left]
\begin{itemize}[nosep,leftmargin=*,itemsep=1pt,font=\scriptsize]
\item Homogénéité
\item Associativité
\item Coefficients même signe $\Rightarrow$ $G$ intérieur
\end{itemize}
\end{tcolorbox}
\begin{tcolorbox}[enhanced,colback=white,colframe=poppurple,boxrule=0.9pt,arc=3pt,title={3 points},fonttitle=\bfseries\scriptsize,coltitle=white,colbacktitle=poppurple,left=4pt,right=4pt,top=2pt,bottom=2pt,boxsep=1pt,toptitle=1.5pt,bottomtitle=1.5pt,halign=left]
\begin{itemize}[nosep,leftmargin=*,itemsep=1pt,font=\scriptsize]
\item $\alpha\overrightarrow{GA}+\beta\overrightarrow{GB}+\gamma\overrightarrow{GC}=\vec{0}$
\item Isobarycentre = centre de gravité
\item Associativité : $G_{ABC} = G_{G_{AB}C}$
\end{itemize}
\end{tcolorbox}
\begin{tcolorbox}[enhanced,colback=white,colframe=poppink,boxrule=0.9pt,arc=3pt,title={4 points},fonttitle=\bfseries\scriptsize,coltitle=white,colbacktitle=poppink,left=4pt,right=4pt,top=2pt,bottom=2pt,boxsep=1pt,toptitle=1.5pt,bottomtitle=1.5pt,halign=left]
\begin{itemize}[nosep,leftmargin=*,itemsep=1pt,font=\scriptsize]
\item Regroupement partiel
\item Parallélogramme : isobarycentre = centre
\end{itemize}
\end{tcolorbox}
\begin{tcolorbox}[enhanced,colback=white,colframe=popgold,boxrule=0.9pt,arc=3pt,title={Lignes de niveau},fonttitle=\bfseries\scriptsize,coltitle=white,colbacktitle=popgold,left=4pt,right=4pt,top=2pt,bottom=2pt,boxsep=1pt,toptitle=1.5pt,bottomtitle=1.5pt,halign=left]
\begin{itemize}[nosep,leftmargin=*,itemsep=1pt,font=\scriptsize]
\item $\|\alpha\overrightarrow{MA}+\beta\overrightarrow{MB}\|=k$
\item Cercle centre $G$, rayon $\frac{k}{|\alpha+\beta|}$
\item $MA^2+MB^2=k$ : cercle centre $I$ milieu
\end{itemize}
\end{tcolorbox}
\begin{tcolorbox}[enhanced,colback=white,colframe=popblue!70,boxrule=0.9pt,arc=3pt,title={Vie courante},fonttitle=\bfseries\scriptsize,coltitle=white,colbacktitle=popblue!70,left=4pt,right=4pt,top=2pt,bottom=2pt,boxsep=1pt,toptitle=1.5pt,bottomtitle=1.5pt,halign=left]
\begin{itemize}[nosep,leftmargin=*,itemsep=1pt,font=\scriptsize]
\item Équilibre, mécanique, chantier
\item Centre de gravité, résultante
\end{itemize}
\end{tcolorbox}
\end{tcbraster}

\legende{Carte mentale de la leçon : les barycentres}
\end{popfigure}

\begin{aretenir}[L'essentiel de la leçon]
\textbf{Définitions clés}
\begin{itemize}
  \item \cle{Barycentre} de $(A,\alpha)$ et $(B,\beta)$ avec $\alpha+\beta\neq 0$ : unique point $G$ tel que $\alpha\overrightarrow{GA}+\beta\overrightarrow{GB}=\vec{0}$.
  \item \cle{Isobarycentre} : barycentre de points affectés de coefficients \emph{égaux} (ex. milieu de $[AB]$, centre de gravité d'un triangle).
  \item Barycentre de trois points $(A,\alpha)$, $(B,\beta)$, $(C,\gamma)$ avec $\alpha+\beta+\gamma\neq 0$ : unique point $G$ tel que $\alpha\overrightarrow{GA}+\beta\overrightarrow{GB}+\gamma\overrightarrow{GC}=\vec{0}$.
  \item Barycentre de $n$ points : définition analogue, somme des coefficients non nulle.
\end{itemize}

\textbf{Formules à connaître par cœur}
\begin{center}
\begin{tabularx}{\linewidth}{|l|X|}
\hline
\rowcolor{popblueL} \textbf{Situation} & \textbf{Formule} \\
\hline
Position de $G$ (2 points) & $\overrightarrow{AG}=\dfrac{\beta}{\alpha+\beta}\overrightarrow{AB}$ \\
\hline
Réduction (2 points) & $\alpha\overrightarrow{MA}+\beta\overrightarrow{MB}=(\alpha+\beta)\overrightarrow{MG}$ \\
\hline
Réduction (3 points) & $\alpha\overrightarrow{MA}+\beta\overrightarrow{MB}+\gamma\overrightarrow{MC}=(\alpha+\beta+\gamma)\overrightarrow{MG}$ \\
\hline
Coordonnées de $G$ (repère orthonormé) & $x_G=\dfrac{\alpha x_A+\beta x_B}{\alpha+\beta},\quad y_G=\dfrac{\alpha y_A+\beta y_B}{\alpha+\beta}$ \\
\hline
Ligne de niveau $\|\alpha\overrightarrow{MA}+\beta\overrightarrow{MB}\|=k$ ($\alpha+\beta\neq 0$) & Cercle de centre $G$, rayon $\dfrac{k}{|\alpha+\beta|}$ \\
\hline
Ligne de niveau $MA^2+MB^2=k$ & Cercle de centre $I$ milieu de $[AB]$ : $MI^2=\dfrac{k}{2}-\dfrac{AB^2}{4}$ \\
\hline
\end{tabularx}
\end{center}

\textbf{Propriétés fondamentales}
\begin{itemize}
  \item \cle{Homogénéité} : on peut multiplier tous les coefficients par un même nombre non nul sans changer le barycentre.
  \item \cle{Associativité} : on peut remplacer un groupe de points par leur barycentre partiel affecté de la somme de leurs coefficients.
  \item Si tous les coefficients sont de même signe (tous $>0$ ou tous $<0$), $G$ appartient au segment $[AB]$ pour 2 points, à l'intérieur du triangle $ABC$ pour 3 points.
  \item Pour 4 points : l'associativité permet de calculer par regroupements successifs (ex. $(A,\alpha),(B,\beta)$ puis $(C,\gamma),(D,\delta)$).
\end{itemize}

\textbf{Méthodes express}
\begin{enumerate}
  \item \textbf{Construire $G$ (2 points)} : écrire $\overrightarrow{AG}=\dfrac{\beta}{\alpha+\beta}\overrightarrow{AB}$, puis reporter le vecteur depuis $A$.
  \item \textbf{Prouver un alignement ou un concours} : exprimer un point comme barycentre de deux façons grâce à l'associativité.
  \item \textbf{Ligne de niveau $\|\alpha\overrightarrow{MA}+\beta\overrightarrow{MB}\|=k$} : réduire la somme vectorielle en $(\alpha+\beta)\overrightarrow{MG}$, identifier le cercle de centre $G$ et de rayon $\frac{k}{|\alpha+\beta|}$.
  \item \textbf{Ligne de niveau $MA^2+MB^2=k$} : utiliser le milieu $I$ de $[AB]$ et la relation $MA^2+MB^2=2MI^2+\frac{AB^2}{2}$.
\end{enumerate}
\end{aretenir}

\begin{aretenir}[Checklist avant le Probatoire]
\begin{itemize}
  \item[$\square$] Je sais écrire la condition vectorielle $\alpha\overrightarrow{GA}+\beta\overrightarrow{GB}=\vec{0}$ définissant le barycentre.
  \item[$\square$] Je vérifie toujours que la somme des coefficients est \textbf{non nulle} avant de conclure à l'existence.
  \item[$\square$] Je sais passer de $\alpha\overrightarrow{GA}+\beta\overrightarrow{GB}=\vec{0}$ à $\overrightarrow{AG}=\dfrac{\beta}{\alpha+\beta}\overrightarrow{AB}$ avec Chasles.
  \item[$\square$] Je sais placer un barycentre sur une figure (attention au coefficient qui accompagne le \emph{bon} point : $G$ est plus proche du point de coefficient plus grand en valeur absolue).
  \item[$\square$] Je sais calculer les coordonnées d'un barycentre dans un repère.
  \item[$\square$] Je sais utiliser l'homogénéité pour simplifier des coefficients (ex. diviser par 2, multiplier par $-1$).
  \item[$\square$] Je sais utiliser l'associativité : barycentre partiel affecté de la somme des coefficients.
  \item[$\square$] Je sais que l'isobarycentre de deux points est le milieu, celui de trois points le centre de gravité du triangle (intersection des médianes).
  \item[$\square$] Je sais réduire $\alpha\overrightarrow{MA}+\beta\overrightarrow{MB}$ en $(\alpha+\beta)\overrightarrow{MG}$.
  \item[$\square$] Je sais déterminer la nature d'une ligne de niveau (cercle, segment, ensemble vide) et préciser centre et rayon.
  \item[$\square$] Je sais rédiger une justification complète et conclure clairement.
\end{itemize}
\end{aretenir}

\findecours

\end{document}
